% 农业上市公司画像：畜禽饲料、水产饲料（共 15 家，按流通市值降序）
% 由 scripts/make_company_profiles.py 自动生成，请勿手工修改

\subsubsection{海大集团（002311）}
\label{co:002311}

\pfl{公司基本情况} 海大集团成立于 2004 年，以水产饲料为核心、覆盖饲料—种苗—动保—养殖的全链饲料龙头，饲料销量与海外扩张领先；2026 年上半年生猪养殖亏损拖累盈利，增收不增利。 公司于 2009-11-27 在深交所主板首发上市，注册资本 16.57 亿元，总股本 16.64 亿股。 截至 2026 年 9 月 24 日收盘价 41.10 元，流通市值 680.4 亿元，近一年-34.1\%，流通市值在三级行业“水产饲料”的 4 家公司中按由大到小排第\zhnumber{1}位，近一年涨跌幅低于全样本中位数（-10.1\%）24.0 个百分点。 按 2026 年半年报披露的行业构成，收入占比前两位为饲料行业 87.6\%；养殖行业 12.4\%。

\pfl{历史股价} 自 2021 年 1 月至 2026 年 9 月，股价由 65.81 元变为 41.10 元，累计 -37.5\%，区间最高 83.40 元（2021 年 4 月）、最低 37.80 元（2024 年 1 月）；当前价相当于区间最高价的 49.3\%，为区间最低价的 1.09 倍。 月度收益的年化波动率 30.2\%，低于全样本中位数 37.4\%，在三级行业“水产饲料”的 4 家公司中波动率由高到低排第\zhnumber{3}位。

\begin{center}
\begin{minipage}{\textwidth}\centering
\begin{tikzpicture}
\begin{axis}[
  width=12.6cm, height=3.9cm, scale only axis,
  xmin=2020.922, xmax=2026.888, ymin=35.15, ymax=89.24,
  xtick={2021,2022,2023,2024,2025,2026},
  yearticks,
  yticklabel style={font=\scriptsize},
  ylabel={元/股}, ylabel style={font=\scriptsize},
  grid=major, grid style={LightGray, line width=0.3pt},
  axis line style={InkGray, line width=0.5pt},
  every axis plot/.append style={line width=0.9pt},
]
\addplot[AgriGreen] table[x=x,y=y]{data/clean/profiles/px_002311.dat};
\addplot[SoftRed, dashed, line width=0.7pt] table[x=x,y=y]{data/clean/profiles/pxzz_002311.dat};
\addplot[only marks, mark=*, mark size=1.1pt, SoftRed] table[x=x,y=y]{data/clean/profiles/pxzz_002311.dat};
\end{axis}
\end{tikzpicture}

\captionof{figure}[海大集团（002311）股价与周期骨架]{海大集团（002311）月度收盘价（元/股）与 ZigZag 周期骨架（阈值 25\%，圆点为周期转折点）}
\end{minipage}
\end{center}

\pfl{过去周期分析} 以 25\% 的 ZigZag 阈值划分，样本区间内股价形成 2 个主升段与 2 个主跌段。 最大主跌段为 2021 年 4 月 至 2024 年 1 月，33 个月-54.7\%；最大主升段出现在 2024 年 1 月 至 2025 年 9 月，20 个月+68.7\%。 最近一次转折出现在 2026 年 9 月（下跌段自 2025 年 9 月起，12 个月-35.5\%），当前处于该段的延续之中。 公司股价较申万农林牧渔指数提前约 28 个月见底，是板块中较早定价周期反转的公司之一。 水产饲料的需求取决于投苗量与存塘鱼价，第 \ref{sec:other-livestock} 章显示 2026 年淡水鱼价格与投苗量已率先回升，其需求驱动与生猪存栏的联系弱于畜禽饲料，公司 2026 年上半年实现盈利、半年 ROE 6.5\%（未年化），好于全样本中位数（0.75\%），下行周期对其报表的冲击小于同业。 综合区间形态看，该股单段涨跌幅度偏大、以趋势段为主（最大主升 +68.7\%、最大主跌 -54.7\%），年化波动率 30.2\% 与全样本中位数 37.4\% 相比波动小于样本中位数。

\pfl{盈利情况} 2026 年上半年营业收入 642.40 亿元（同比 +9.2\%），归母净利润 16.11 亿元，同比 -39.0\%。 以年报为趋势对照，2023---2025 年营业收入由 1,161.17 亿元变为 1,284.68 亿元（两年年均复合增速 +5.2\%），归母净利润由 27.41 亿元变为 42.80 亿元。 按半年报口径，2026 年上半年的 ROE 为 6.5\%（未年化）、销售净利率 2.7\%、毛利率 9.3\%，ROE 在“水产饲料”4 家公司中按数值由高到低排第\zhnumber{1}位。 同期全样本半年 ROE 中位数 0.75\%，公司与之相差 5.72 个百分点（略好于中位数公司）。 从公司披露的原因看，业绩变动主要来自：2026 年上半年，饲料行业营收 562.67 亿元、同比增长 14.61\%，为收入核心。 综合看，公司在报告期维持盈利（高于全样本半年 ROE 中位数 0.75\%），利润的可持续性取决于水产饲料景气的延续与成本管控的执行。

\begin{center}\small\setlength{\tabcolsep}{3.4pt}
\captionof{table}[海大集团（002311）关键财务指标]{海大集团（002311）关键财务指标（合并报表口径；两个半年报期均未年化；现金短债比 $=$ 货币资金 $\div$ 短期债务）}
\begin{adjustbox}{max width=\textwidth}
\begin{tabular}{@{}lrrrrr@{}}
\toprule
指标 & 2023 年 & 2024 年 & 2025 年 & 2025 年半年报 & 2026 年半年报 \\
\midrule
营业收入（亿元） & 1,161.17 & 1,146.01 & 1,284.68 & 588.31 & 642.40 \\
归母净利润（亿元） & 27.41 & 45.04 & 42.80 & 26.39 & 16.11 \\
毛利率（\%） & 8.5 & 11.3 & 10.3 & 11.7 & 9.3 \\
ROE（\%） & 14.5 & 20.9 & 17.4 & 10.7 & 6.5 \\
资产负债率（\%） & 53.3 & 47.9 & 46.4 & 49.2 & 54.6 \\
经营活动现金流净额（亿元） & 126.98 & 79.96 & 62.60 & 30.48 & 39.24 \\
货币资金（亿元） & 54.76 & 34.78 & 32.03 & 34.81 & 34.25 \\
有息负债合计（亿元） & 70.99 & 52.67 & 40.82 & 44.63 & 64.39 \\
现金短债比（倍） & 2.23 & 2.35 & 1.91 & 2.19 & 0.84 \\
\bottomrule
\end{tabular}
\end{adjustbox}
\end{center}

\begin{center}
\begin{minipage}{\textwidth}\centering
\begin{tikzpicture}
\begin{groupplot}[
  group style={group size=2 by 1, horizontal sep=0.60cm},
  width=6.85cm, height=4.60cm,
  tick label style={font=\tiny},
  title style={font=\scriptsize\heiti},
  yticklabel style={font=\tiny},
  ylabel style={font=\tiny},
  grid=major, grid style={LightGray, line width=0.3pt},
  axis line style={InkGray, line width=0.5pt},
  enlarge x limits={abs=0.8},
  legend style={font=\scriptsize, at={(0.5,-0.20)}, anchor=north,
                legend columns=2, draw=none, fill=none, column sep=6pt},
]
\nextgroupplot[
  ybar, bar width=4pt,
  ymin=-128.47, ymax=1438.84,
  xtick={0,1,2,3,4,5.7},
  xticklabels={2021,2022,2023,2024,2025,2026H1},
  ylabel={亿元},
]
\addplot[fill=AgriGreen!75, draw=AgriGreen, bar shift=-2.6pt] table[x=i, y=rev]{data/clean/profiles/fin_002311.dat};
\addplot[fill=SoftRed!60, draw=SoftRed, bar shift=2.6pt] table[x=i, y=ni]{data/clean/profiles/fin_002311.dat};
\addplot[fill=AgriGreen!30, draw=AgriGreen, pattern=north east lines, bar shift=-2.6pt] table[x=x, y=rev]{data/clean/profiles/finh_002311.dat};
\addplot[fill=SoftRed!20, draw=SoftRed, pattern=north east lines, bar shift=2.6pt] table[x=x, y=ni]{data/clean/profiles/finh_002311.dat};
\legend{营业收入（年/半年）, 归母净利润（年/半年）}
\nextgroupplot[
  ymin=-4.5, ymax=61.9,
  xtick={0,1,2,3,4,5.7},
  xticklabels={2021,2022,2023,2024,2025,2026H1},
  ylabel={\%},
]
\addplot[AgriGold, mark=*, mark size=1.3pt] table[x=i, y=roe]{data/clean/profiles/fin_002311.dat};
\addplot[OilBlue, mark=square*, mark size=1.3pt] table[x=i, y=debt]{data/clean/profiles/fin_002311.dat};
\addplot[only marks, AgriGold, mark=*, mark size=2.0pt] table[x=x, y=roe]{data/clean/profiles/finh_002311.dat};
\addplot[only marks, OilBlue, mark=square*, mark size=2.0pt] table[x=x, y=debt]{data/clean/profiles/finh_002311.dat};
\legend{ROE, 资产负债率}
\end{groupplot}
\end{tikzpicture}
\begin{tikzpicture}
\begin{axis}[
  name=finbot,
  width=13.75cm, height=3.10cm, scale only axis,
  ymin=-10.85, ymax=121.48,
  xtick={0,1,2,3,4,5.7},
  xticklabels={2021,2022,2023,2024,2025,2026H1},
  tick label style={font=\tiny},
  yticklabel style={font=\tiny},
  ylabel={亿元}, ylabel style={font=\tiny},
  ybar, bar width=4pt,
  grid=major, grid style={LightGray, line width=0.3pt},
  axis line style={InkGray, line width=0.5pt},
  enlarge x limits={abs=0.8},
  legend style={font=\scriptsize, at={(0.5,-0.26)}, anchor=north,
                legend columns=2, draw=none, fill=none, column sep=10pt},
]
\addplot[fill=AgriGreen!55, draw=AgriGreen, bar shift=-3.0pt] table[x=i, y=cash]{data/clean/profiles/fin_002311.dat};
\addplot[fill=SoftRed!45, draw=SoftRed, bar shift=3.0pt] table[x=i, y=ibd]{data/clean/profiles/fin_002311.dat};
\addplot[fill=AgriGreen!25, draw=AgriGreen, pattern=north east lines, bar shift=-3.0pt] table[x=x, y=cash]{data/clean/profiles/finh_002311.dat};
\addplot[fill=SoftRed!20, draw=SoftRed, pattern=north east lines, bar shift=3.0pt] table[x=x, y=ibd]{data/clean/profiles/finh_002311.dat};
\legend{货币资金（左轴）, 有息负债（左轴）}
\end{axis}
\begin{axis}[
  at={(finbot.south east)}, anchor=south east,
  width=13.75cm, height=3.10cm, scale only axis,
  axis y line*=right, axis x line*=none, xtick=\empty,
  ymin=-33.02, ymax=165.08,
  tick label style={font=\tiny},
  yticklabel style={font=\tiny}, scaled y ticks=false,
  legend style={font=\scriptsize, at={(0.5,-0.44)}, anchor=north,
                legend columns=1, draw=none, fill=none},
]
\addplot[OilBlue, mark=*, mark size=2.2pt, line width=1.3pt] table[x=i, y=ocf]{data/clean/profiles/fin_002311.dat};
\addplot[only marks, OilBlue, mark=*, mark size=3.2pt] table[x=x, y=ocf]{data/clean/profiles/finh_002311.dat};
\legend{经营活动现金流净额（右轴，亿元，蓝点）}
\end{axis}
\end{tikzpicture}

\captionof{figure}[海大集团（002311）近年主要财务指标]{海大集团（002311）近年主要财务指标（左上：营业收入与归母净利润；右上：ROE 与资产负债率；下：货币资金、有息负债为左轴柱状，经营活动现金流净额为其右轴的蓝点折线。
2021---2025 年为年报口径，2026H1 为 2026 年半年报/半年末，与其前一年报之间留出间隔、以斜纹或空心点区分；半年报数据未年化）}
\end{minipage}
\end{center}

\pfl{财务分析} 2026 年 6 月末资产负债率 54.6\%，较 2025 年末上升 8.2 个百分点（样本中位数 52.2\%，所处三级行业内按由低到高排第\zhnumber{2}位）。 2026 年 6 月末货币资金 34.25 亿元、短期债务（短期借款与一年内到期非流动负债）40.97 亿元、长期债务 23.42 亿元，有息负债合计 64.39 亿元，现金短债比 0.84 倍——覆盖偏紧。 净有息负债 30.14 亿元，相当于其 2026 年上半年营业收入的 5\%。 流动比率 1.06、速动比率 0.55，短期指标尚可。 2026 年上半年经营活动现金流净额 39.24 亿元，经营现金流与利润同向为正，内生资金可以支撑运营。 2026 年 6 月末在建工程 8.71 亿元，较上年末增加 2.67 亿元，仍有在建投入。 2026 年 6 月末商誉 3.88 亿元，占归母净资产的 1.5\%，是净资产中最脆弱的一块。 综合看，现金短债比 0.84 倍，短期资金安排偏紧；资产负债率高于样本中位数 2.4 个百分点；有息负债中短期债务占比更高，期限结构仍以滚动续作为主。 公司披露的股东与质押风险为：2026 年 7 月，控股股东广州市海灏投资有限公司质押公司 320 万股；2026 年 9 月，控股股东海灏投资再质押 300 万股（占其所持 0.33\%）。

\pfl{重大战略转型} 近三年公司的战略动作集中在：其他动作（2 项）：2024 年，股票期权激励计划推进，公司调整该计划股票期权行权价格；2026 年 8 月，半年报提出“饲料、种苗和动保三体联动”核心业务模式，以全球化布局与技术服务优势对冲原料价格波动；海外与国际化（1 项）：2024---2026 年，2026 年上半年海外饲料销量增速超 25\%。 公告口径上，近三年（2023 年 9 月---2026 年 9 月）公司披露重大重组与并购类公告 0 条、对外投资与转型类公告 0 条、回购与股权激励类公告 54 条，合计公告 448 条。 第一大行业仍是“饲料行业”，收入占比 87.6\%，与 2021 年年报相比没有方向性变化。 综合判断，报告期内公司未披露实质性的重组、并购或转型安排，资本开支以维持现有产能与业务为主。

\begin{center}\small\setlength{\tabcolsep}{4pt}
\captionof{table}[海大集团（002311）历史融资与资本运作]{海大集团（002311）历史融资与资本运作（据公司公告与公开报道整理；金额为披露值，含首次公开发行、再融资、债券、银行借款与重整投资等）}
\begin{adjustbox}{max width=\textwidth}
\begin{tabular}{@{}p{1.45cm}p{2.55cm}p{4.05cm}p{4.95cm}@{}}
\toprule
时间 & 方式 & 规模 & 用途与说明 \\
\midrule
2009-11 & IPO（深交所中小板） & 发行 5，\allowbreak{}600 万股、\allowbreak{}28.00 元/\allowbreak{}股 & 上市日 2009-11-27，\allowbreak{}发行市盈率 48.39 倍 \\
2013-11 & 定向增发（非公开发行） & 6,\allowbreak{}850 万股、\allowbreak{}11.34 元/\allowbreak{}股，\allowbreak{}募集资金净额 7.58 亿元 & 新增股份 2013-12-06 上市，\allowbreak{}限售 12 个月 \\
2020-03 & 公开发行可转换公司债券 & 30.00 亿元（2,\allowbreak{}830 万张，\allowbreak{}面值 100.00 元） & 2020-03-19 发行、\allowbreak{}2020-04-16 上市 \\
2020-03 & 可转债（募集资金专户销户） & 28.30 亿元 & 2023-12-28 公告可转债募集资金专户销户完成，\allowbreak{}募集资金已使用完毕 \\
2024 & 现金分红（2024 年度） & 每 10 股派 11.00 元（含税），\allowbreak{}合计派现 18.30 亿元 & 股东回报；上市以来累计实施 19 次现金分红，\allowbreak{}形成“年度高分红+中期稳回馈”模式；2025 年半年报再推 10 派 2.00 元共 3.33 亿元 \\
2026-09 & 2026 年中期权益分派 & 每 10 股派现金红利 1.50 元（含税），\allowbreak{}回购专户 7.56 万股不参与分派 & 股东回报；2026 年 9 月 10 日发布实施公告 \\
\bottomrule
\end{tabular}
\end{adjustbox}
\end{center}

\pfl{历史融投资情况} 从现金流量表看，2025 年公司取得借款收到的现金 50.54 亿元、偿还债务支付的现金 65.82 亿元、吸收权益性投资 0.73 亿元，筹资活动现金流净额 -51.76 亿元（2023 年为取得借款 98.05 亿元、偿还债务 135.73 亿元）。 2026 年上半年筹资活动现金流净额为 -7.25 亿元（取得借款 65.37 亿元、偿还债务 42.72 亿元，未年化），已转为净流出，处于净还债状态。 2025 年分配股利、利润或偿付利息支付的现金合计 26.24 亿元。 公开披露的融资记录共 6 笔，工具类型以 IPO、定向增发、可转债为主；金额与用途见下表。 其中规模最大的两笔为：2020 年 3 月，可转债，30.00 亿元（2,830 万张，面值 100.00 元）；2020 年 3 月，可转债，28.30 亿元。 股本方面，近三年披露股本变动 6 次（期权行权 2 次、其他 1 次）；总股本由 2021 年末的 16.61 亿股变为 16.64 亿股（+0.16\%）。 分红方面，近三年现金分红 6 次、合计每股派现 3.50 元，最近一次实施在 2026 年 6 月。 近三年“再融资”类公告 3 条，涉及定向增发、可转债、募集资金使用。 综合看，公司融资结构上以债务滚动为主、股权融资为补充；2025 年筹资活动现金流净额 -51.76 亿元，融资节奏仍以维持存量债务周转为主，与前述经营与投资现金流的方向基本一致。

\pfl{未来潜在融资需求分析} 2026 年 6 月末货币资金 34.25 亿元，而短期债务 40.97 亿元（现金短债比 0.84 倍），2026 年上半年经营现金流净额 39.24 亿元。按“短期债务减去账面货币资金”的滚动偿付口径，静态缺口约 6.72 亿元；上半年盈利或接近盈亏平衡，该缺口不随经营亏损进一步扩大。 按第 \ref{sec:financing-need} 节的三档划分，公司属于第三档“扩张型”：2026 年上半年 ROE 6.5\%（未年化，折合约 12.9\%）、6 月末资产负债率 54.6\%，资产负债表仍具备加杠杆空间；融资用途以产能扩张、品类并购与海外布局为主，单笔规模通常在 5---20 亿元量级，会被市场定价为成长而非补血。 公开披露的后续资金安排线索包括：股权融资线索：2022 年 4 月，公司募集资金总额不超过 15.00 亿元；2023 年 3 月，广发证券出具 2022 年度向特定对象发行 A 股股票发行保荐书并取得证监会同意注册决定。 资金需求还要叠加在建工程：期末在建工程 8.71 亿元、2026 年上半年购建固定资产等支出 16.93 亿元（未年化），这部分支出在周期底部通常会被主动放缓。 综合看，公司的外部资金需求以营运资金周转与在建项目投入为主，滚动偿付口径下的静态缺口约 6.72 亿元，融资的时点与规模更多取决于信用条件与资本市场窗口。




\subsubsection{新希望（000876）}
\label{co:000876}

\pfl{公司基本情况} 1998 年设立至今，新希望是农牧全产业链龙头（饲料+生猪养殖与屠宰），2023 年底白羽肉禽与食品深加工引战出表以聚焦主业；2026 年上半年饲料创新高，猪价探底致亏损 17.02 亿元。 公司于 1998-03-11 在深交所主板通过重大资产重组（借壳）上市，注册资本 49.62 亿元，总股本 45.46 亿股。 截至 2026 年 9 月 24 日收盘价 6.92 元，流通市值 311.4 亿元，近一年-30.1\%，流通市值在三级行业“畜禽饲料”的 11 家公司中按由大到小排第\zhnumber{1}位，近一年涨跌幅低于全样本中位数（-10.1\%）20.0 个百分点。 按 2026 年半年报披露的行业构成，收入占比前三位为饲料 74.5\%；猪产业 23.3\%；其他 2.2\%。

\pfl{历史股价} 本报告样本区间（2021 年 1 月 至 2026 年 9 月）内，股价由 21.15 元变为 6.92 元，累计 -67.3\%，区间最高 24.07 元（2021 年 2 月）、最低 6.27 元（2026 年 6 月）；当前价相当于区间最高价的 28.7\%，为区间最低价的 1.10 倍。 月度收益的年化波动率 31.5\%，低于全样本中位数 37.4\%，在三级行业“畜禽饲料”的 11 家公司中波动率由高到低排第\zhnumber{10}位。 把股价阶段与公司事件对齐看，2022 年 3 月 至 2024 年 1 月股价下跌-51.9\%，同期（2023 年 12 月）董事会通过白羽肉禽业务引进战略投资者议案：中国牧工商集团有限公司以现金约 27.00 亿元受让禽产业链运营主体山东中新食品集团 51\% 股权。 从时间对应关系看，公司层面的资本运作与股权、风险事件解释了区间内多数急涨急跌，与行业指数的同步性相对有限。

\begin{center}
\begin{minipage}{\textwidth}\centering
\begin{tikzpicture}
\begin{axis}[
  width=12.6cm, height=3.9cm, scale only axis,
  xmin=2020.922, xmax=2026.888, ymin=5.83, ymax=25.75,
  xtick={2021,2022,2023,2024,2025,2026},
  yearticks,
  yticklabel style={font=\scriptsize},
  ylabel={元/股}, ylabel style={font=\scriptsize},
  grid=major, grid style={LightGray, line width=0.3pt},
  axis line style={InkGray, line width=0.5pt},
  every axis plot/.append style={line width=0.9pt},
]
\addplot[AgriGreen] table[x=x,y=y]{data/clean/profiles/px_000876.dat};
\addplot[SoftRed, dashed, line width=0.7pt] table[x=x,y=y]{data/clean/profiles/pxzz_000876.dat};
\addplot[only marks, mark=*, mark size=1.1pt, SoftRed] table[x=x,y=y]{data/clean/profiles/pxzz_000876.dat};
\end{axis}
\end{tikzpicture}

\captionof{figure}[新希望（000876）股价与周期骨架]{新希望（000876）月度收盘价（元/股）与 ZigZag 周期骨架（阈值 25\%，圆点为周期转折点）}
\end{minipage}
\end{center}

\pfl{过去周期分析} 以 25\% 的 ZigZag 阈值划分，样本区间内股价形成 2 个主升段与 3 个主跌段。 最大主跌段为 2022 年 3 月 至 2024 年 1 月，22 个月-51.9\%；最大主升段出现在 2021 年 8 月 至 2022 年 3 月，7 个月+55.2\%。 最近一次转折出现在 2026 年 6 月（下跌段自 2024 年 9 月起，21 个月-39.8\%），当前处于该段的延续之中。 公司股价与申万农林牧渔指数几乎同步见底，个股并未走出独立行情。 饲料夹在原料成本与养殖需求之间（第 \ref{sec:feed} 节），量的部分取决于存栏、利的部分取决于原料与成品的价差，公司 2026 年 6 月末资产负债率 70.7\%，高于全样本中位数 52.2\%，其股价因此更多反映资金面与信用风险，而不是价格弹性本身。 综合区间形态看，该股单段涨跌幅度偏大、以趋势段为主（最大主升 +55.2\%、最大主跌 -51.9\%），年化波动率 31.5\% 与全样本中位数 37.4\% 相比波动小于样本中位数。

\pfl{盈利情况} 2026 年上半年营业收入 555.89 亿元（同比 +7.7\%），归母净利润 -17.02 亿元，由上年同期的盈利转为亏损。 以年报为趋势对照，2023---2025 年营业收入由 1,417.03 亿元变为 1,068.56 亿元（两年年均复合增速 -13.2\%），归母净利润由 2.49 亿元变为 -17.84 亿元。 按半年报口径，2026 年上半年的 ROE 为 -9.1\%（未年化）、销售净利率 -3.0\%、毛利率 4.7\%，ROE 在“畜禽饲料”11 家公司中按数值由高到低排第\zhnumber{8}位。 同期全样本半年 ROE 中位数 0.75\%，公司与之相差 9.87 个百分点（弱于中位数公司）。 公司给出的应对方向是：2026 年 8 月，董事会通过使用向特定对象发行募集资金置换预先投入募投项目的自筹资金 7.95 亿元及已支付发行费用 92.11 万元。 综合看，公司仍处亏损区间、由盈转亏，半年 ROE 在三级行业内按由高到低排第\zhnumber{8}位，单位盈利能力的修复依赖行业景气与自身成本端的改善，报表弹性主要体现为价格与销量的方向性变化。

\begin{center}\small\setlength{\tabcolsep}{3.4pt}
\captionof{table}[新希望（000876）关键财务指标]{新希望（000876）关键财务指标（合并报表口径；两个半年报期均未年化；现金短债比 $=$ 货币资金 $\div$ 短期债务）}
\begin{adjustbox}{max width=\textwidth}
\begin{tabular}{@{}lrrrrr@{}}
\toprule
指标 & 2023 年 & 2024 年 & 2025 年 & 2025 年半年报 & 2026 年半年报 \\
\midrule
营业收入（亿元） & 1,417.03 & 1,030.63 & 1,068.56 & 516.25 & 555.89 \\
归母净利润（亿元） & 2.49 & 4.74 & -17.84 & 7.55 & -17.02 \\
毛利率（\%） & 2.8 & 6.7 & 6.3 & 7.9 & 4.7 \\
ROE（\%） & 0.7 & 1.7 & -8.8 & 2.9 & -9.1 \\
资产负债率（\%） & 72.3 & 69.0 & 69.9 & 68.8 & 70.7 \\
经营活动现金流净额（亿元） & 139.04 & 91.27 & 93.78 & 51.53 & 16.46 \\
货币资金（亿元） & 108.50 & 86.99 & 86.27 & 70.72 & 79.63 \\
有息负债合计（亿元） & 639.65 & 614.09 & 582.06 & 589.73 & 579.09 \\
现金短债比（倍） & 0.41 & 0.30 & 0.37 & 0.25 & 0.39 \\
\bottomrule
\end{tabular}
\end{adjustbox}
\end{center}

\begin{center}
\begin{minipage}{\textwidth}\centering
\begin{tikzpicture}
\begin{groupplot}[
  group style={group size=2 by 1, horizontal sep=0.60cm},
  width=6.85cm, height=4.60cm,
  tick label style={font=\tiny},
  title style={font=\scriptsize\heiti},
  yticklabel style={font=\tiny},
  ylabel style={font=\tiny},
  grid=major, grid style={LightGray, line width=0.3pt},
  axis line style={InkGray, line width=0.5pt},
  enlarge x limits={abs=0.8},
  legend style={font=\scriptsize, at={(0.5,-0.20)}, anchor=north,
                legend columns=2, draw=none, fill=none, column sep=6pt},
]
\nextgroupplot[
  ybar, bar width=4pt,
  ymin=-247.20, ymax=1598.59,
  xtick={0,1,2,3,4,5.7},
  xticklabels={2021,2022,2023,2024,2025,2026H1},
  ylabel={亿元},
]
\addplot[fill=AgriGreen!75, draw=AgriGreen, bar shift=-2.6pt] table[x=i, y=rev]{data/clean/profiles/fin_000876.dat};
\addplot[fill=SoftRed!60, draw=SoftRed, bar shift=2.6pt] table[x=i, y=ni]{data/clean/profiles/fin_000876.dat};
\addplot[fill=AgriGreen!30, draw=AgriGreen, pattern=north east lines, bar shift=-2.6pt] table[x=x, y=rev]{data/clean/profiles/finh_000876.dat};
\addplot[fill=SoftRed!20, draw=SoftRed, pattern=north east lines, bar shift=2.6pt] table[x=x, y=ni]{data/clean/profiles/finh_000876.dat};
\legend{营业收入（年/半年）, 归母净利润（年/半年）}
\nextgroupplot[
  ymin=-36.6, ymax=82.4,
  xtick={0,1,2,3,4,5.7},
  xticklabels={2021,2022,2023,2024,2025,2026H1},
  ylabel={\%},
]
\addplot[AgriGold, mark=*, mark size=1.3pt] table[x=i, y=roe]{data/clean/profiles/fin_000876.dat};
\addplot[OilBlue, mark=square*, mark size=1.3pt] table[x=i, y=debt]{data/clean/profiles/fin_000876.dat};
\addplot[only marks, AgriGold, mark=*, mark size=2.0pt] table[x=x, y=roe]{data/clean/profiles/finh_000876.dat};
\addplot[only marks, OilBlue, mark=square*, mark size=2.0pt] table[x=x, y=debt]{data/clean/profiles/finh_000876.dat};
\legend{ROE, 资产负债率}
\end{groupplot}
\end{tikzpicture}
\begin{tikzpicture}
\begin{axis}[
  name=finbot,
  width=13.75cm, height=3.10cm, scale only axis,
  ymin=-69.36, ymax=776.85,
  xtick={0,1,2,3,4,5.7},
  xticklabels={2021,2022,2023,2024,2025,2026H1},
  tick label style={font=\tiny},
  yticklabel style={font=\tiny},
  ylabel={亿元}, ylabel style={font=\tiny},
  ybar, bar width=4pt,
  grid=major, grid style={LightGray, line width=0.3pt},
  axis line style={InkGray, line width=0.5pt},
  enlarge x limits={abs=0.8},
  legend style={font=\scriptsize, at={(0.5,-0.26)}, anchor=north,
                legend columns=2, draw=none, fill=none, column sep=10pt},
]
\addplot[fill=AgriGreen!55, draw=AgriGreen, bar shift=-3.0pt] table[x=i, y=cash]{data/clean/profiles/fin_000876.dat};
\addplot[fill=SoftRed!45, draw=SoftRed, bar shift=3.0pt] table[x=i, y=ibd]{data/clean/profiles/fin_000876.dat};
\addplot[fill=AgriGreen!25, draw=AgriGreen, pattern=north east lines, bar shift=-3.0pt] table[x=x, y=cash]{data/clean/profiles/finh_000876.dat};
\addplot[fill=SoftRed!20, draw=SoftRed, pattern=north east lines, bar shift=3.0pt] table[x=x, y=ibd]{data/clean/profiles/finh_000876.dat};
\legend{货币资金（左轴）, 有息负债（左轴）}
\end{axis}
\begin{axis}[
  at={(finbot.south east)}, anchor=south east,
  width=13.75cm, height=3.10cm, scale only axis,
  axis y line*=right, axis x line*=none, xtick=\empty,
  ymin=-36.15, ymax=180.75,
  tick label style={font=\tiny},
  yticklabel style={font=\tiny}, scaled y ticks=false,
  legend style={font=\scriptsize, at={(0.5,-0.44)}, anchor=north,
                legend columns=1, draw=none, fill=none},
]
\addplot[OilBlue, mark=*, mark size=2.2pt, line width=1.3pt] table[x=i, y=ocf]{data/clean/profiles/fin_000876.dat};
\addplot[only marks, OilBlue, mark=*, mark size=3.2pt] table[x=x, y=ocf]{data/clean/profiles/finh_000876.dat};
\legend{经营活动现金流净额（右轴，亿元，蓝点）}
\end{axis}
\end{tikzpicture}

\captionof{figure}[新希望（000876）近年主要财务指标]{新希望（000876）近年主要财务指标（左上：营业收入与归母净利润；右上：ROE 与资产负债率；下：货币资金、有息负债为左轴柱状，经营活动现金流净额为其右轴的蓝点折线。
2021---2025 年为年报口径，2026H1 为 2026 年半年报/半年末，与其前一年报之间留出间隔、以斜纹或空心点区分；半年报数据未年化）}
\end{minipage}
\end{center}

\pfl{财务分析} 2026 年 6 月末资产负债率 70.7\%，较 2025 年末上升 0.8 个百分点（样本中位数 52.2\%，所处三级行业内按由低到高排第\zhnumber{9}位）。 2026 年 6 月末货币资金 79.63 亿元、短期债务（短期借款与一年内到期非流动负债）205.48 亿元、长期债务 373.61 亿元，有息负债合计 579.09 亿元，现金短债比 0.39 倍——覆盖偏紧。 净有息负债 499.46 亿元，相当于其 2026 年上半年营业收入的 90\%。 流动比率 0.55、速动比率 0.28，短期偿付压力较大。 2026 年上半年经营活动现金流净额 16.46 亿元，净利润为负而经营现金流为正，差额主要来自折旧摊销与营运资本变动，说明亏损尚未完全转化为现金流出。 2026 年 6 月末在建工程 35.71 亿元，较上年末增加 3.71 亿元，仍有在建投入。 2026 年 6 月末商誉 11.27 亿元，占归母净资产的 3.3\%，是净资产中最脆弱的一块。 综合看，账面货币资金对有息负债与短期债务的覆盖不足，滚动偿付对新增融资的依赖度较高；资产负债率高于样本中位数 18.5 个百分点；有息负债中长期债务占比更高，期限结构对短债滚动的压力较小。

\pfl{重大战略转型} 近三年公司的战略动作集中在：并购与重组（2 项）：2023 年 12 月，董事会通过白羽肉禽业务引进战略投资者议案：中国牧工商集团有限公司以现金约 27.00 亿元受让禽产业链运营主体山东中新食品集团 51\% 股权。 公告口径上，近三年（2023 年 9 月---2026 年 9 月）公司披露重大重组与并购类公告 4 条、对外投资与转型类公告 5 条、回购与股权激励类公告 28 条，合计公告 667 条。 第一大行业“饲料”的收入占比由 2021 年年报的 56.1\% 变为 74.5\%（18.5 个百分点），收入结构出现再平衡。 综合判断，报告期内公司有 9 条与战略相关的公告落地（重大重组与并购 4 条、对外投资与转型 5 条），资本运作与主业调整之间存在直接关联。

\begin{center}\small\setlength{\tabcolsep}{4pt}
\captionof{table}[新希望（000876）历史融资与资本运作]{新希望（000876）历史融资与资本运作（据公司公告与公开报道整理；金额为披露值，含首次公开发行、再融资、债券、银行借款与重整投资等）}
\begin{adjustbox}{max width=\textwidth}
\begin{tabular}{@{}p{1.45cm}p{2.55cm}p{4.05cm}p{4.95cm}@{}}
\toprule
时间 & 方式 & 规模 & 用途与说明 \\
\midrule
1998-03 & IPO（深交所） & 公众股 3，\allowbreak{}600 万股、\allowbreak{}内部职工股 400 万股 & 公司原名四川新希望农业股份有限公司 \\
2011 & 发行股份购买资产 & 新增股份 9.05 亿股、\allowbreak{}8.00 元/\allowbreak{}股 & 收购饲料及农牧相关资产（资产出售、\allowbreak{}资产置换及发行股份购买资产一揽子方案）；2011-10-14 完成股份登记不构成借壳上市 \\
2020-01 & 公开发行可转换公司债券（“希望转债”127015） & 40.00 亿元（4,\allowbreak{}000 万张，\allowbreak{}面值 100.00 元），\allowbreak{}期限 6 年 & 2020-01-03 发行、\allowbreak{}2020-02-04 起深交所挂牌 \\
2021-11 & 公开发行可转换公司债券（“希望转2”127049） & 50.00 亿元（8,\allowbreak{}150 万张，\allowbreak{}面值 100.00 元） & 投资生猪养殖项目 57.05 亿元 \\
2023-12 & 资产出售（白羽肉禽与食品深加工引战） & 两笔交易金额合计约 42.00 亿元（27.00 亿元 + 15.01 亿元） & 聚焦核心主业、\allowbreak{}降低资产负债率；2023 年该两项引战影响归母净利润增加约 51.00 亿元至 52.00 亿元 \\
2026-08 & 向特定对象发行股票（定向增发） & 4.59 亿股、\allowbreak{}5.66 元/\allowbreak{}股，\allowbreak{}募集资金总额 26.00 亿元（26.00 亿元） & 猪场生物安全防控及数智化升级项目、\allowbreak{}偿还银行债务；；26 家机构投资者认购，\allowbreak{}总股本增至 49.62 亿股 \\
\bottomrule
\end{tabular}
\end{adjustbox}
\end{center}

\pfl{历史融投资情况} 从现金流量表看，2025 年公司取得借款收到的现金 464.49 亿元、偿还债务支付的现金 518.94 亿元、吸收权益性投资 37.32 亿元，筹资活动现金流净额 -82.59 亿元（2023 年为取得借款 432.13 亿元、偿还债务 485.38 亿元）。 2026 年上半年筹资活动现金流净额为 -13.37 亿元（取得借款 208.32 亿元、偿还债务 229.84 亿元，未年化），已转为净流出，处于净还债状态。 2025 年分配股利、利润或偿付利息支付的现金合计 22.38 亿元。 公开披露的融资记录共 6 笔，工具类型以 IPO、发行股份购买资产、可转债、定向增发为主；金额与用途见下表。 其中规模最大的两笔为：2026 年 8 月，定向增发，4.59 亿股、5.66 元/股，募集资金总额 26.00 亿元（26.00 亿元）；2021 年 11 月，可转债，50.00 亿元（8,150 万张，面值 100.00 元）。 股本方面，近三年披露股本变动 16 次（可转债转股 11 次、股份回购 2 次、增发新股上市 1 次、股权激励 1 次）；总股本由 2021 年末的 45.05 亿股变为 45.46 亿股（+0.90\%）。 分红方面，近三年现金分红 1 次、合计每股派现 0.02 元，最近一次实施在 2024 年 12 月。 近三年“再融资”类公告 95 条，涉及定向增发、可转债、募集资金使用。 综合看，公司融资结构上以债务滚动为主、股权融资为补充；2025 年筹资活动现金流净额 -82.59 亿元，融资节奏仍以维持存量债务周转为主，与前述经营与投资现金流的方向基本一致。

\pfl{未来潜在融资需求分析} 2026 年 6 月末货币资金 79.63 亿元，而短期债务 205.48 亿元（现金短债比 0.39 倍），2026 年上半年经营现金流净额 16.46 亿元。按“短期债务减去账面货币资金”的滚动偿付口径，静态缺口约 125.85 亿元；若再计入上半年亏损的年化额（34.04 亿元），维持一年正常经营所需的外部资金约 125.85---159.89 亿元。 按第 \ref{sec:financing-need} 节的三档划分，公司属于第一档“补血型”：2026 年上半年归母净利润 -17.02 亿元、6 月末资产负债率 70.7\%。 可行的融资路径按现实性排序为：定向增发（需股价配合，且控股股东需放弃优先认购或同步增资）、出售非核心资产与参股股权、引入国资或产业投资人（部分公司以控制权让渡为对价）。 公开披露的后续资金安排线索包括：股权融资线索：2023 年 12 月，定增方案三次调整：2023 年 12 月预案拟募不超过 73.50 亿元 → 2024 年 8 月修订为不超过 38.00 亿元；2026 年 8 月，规范 26.00 亿元定增募集资金存放与使用；债务与授信安排：2025 年，实际控制人下辖财务公司给予公司综合授信额度 130.00 亿元，为流动性补充；2026 年 6 月，截至评级时公司获得授信总额 798.17 亿元、尚未使用额度 337.96 亿元。 资金需求还要叠加在建工程：期末在建工程 35.71 亿元、2026 年上半年购建固定资产等支出 11.35 亿元（未年化），这部分支出在周期底部通常会被主动放缓。 综合看，公司的外部资金需求以债务接续为主，滚动偿付口径下的静态缺口约 125.85 亿元，融资的时点与规模更多取决于信用条件与资本市场窗口。




\subsubsection{大北农（002385）}
\label{co:002385}

\pfl{公司基本情况} 2010 年设立至今，大北农是以饲料、种业（含转基因生物育种）、生猪养殖与动保为主业的农业科技公司；2026 年上半年猪价低迷致亏损 6.77 亿元，实控人邵根伙去世后控制权由配偶莫云继承。 公司于 2010-04-09 在深交所主板首发上市，注册资本 42.80 亿元，总股本 43.24 亿股。 截至 2026 年 9 月 24 日收盘价 3.09 元，流通市值 106.4 亿元，近一年-25.4\%，流通市值在三级行业“畜禽饲料”的 11 家公司中按由大到小排第\zhnumber{2}位，近一年涨跌幅低于全样本中位数（-10.1\%）15.3 个百分点。 按 2026 年半年报披露的产品构成，收入占比前三位为饲料产品 67.7\%；养猪产品 21.7\%；其他产品 7.3\%。

\pfl{历史股价} 自 2021 年 1 月至 2026 年 9 月，股价由 9.27 元变为 3.09 元，累计 -66.7\%，区间最高 10.49 元（2021 年 12 月）、最低 2.76 元（2026 年 6 月）；当前价相当于区间最高价的 29.5\%，为区间最低价的 1.12 倍。 月度收益的年化波动率 33.3\%，低于全样本中位数 37.4\%，在三级行业“畜禽饲料”的 11 家公司中波动率由高到低排第\zhnumber{9}位。

\begin{center}
\begin{minipage}{\textwidth}\centering
\begin{tikzpicture}
\begin{axis}[
  width=12.6cm, height=3.9cm, scale only axis,
  xmin=2020.922, xmax=2026.888, ymin=2.57, ymax=11.22,
  xtick={2021,2022,2023,2024,2025,2026},
  yearticks,
  yticklabel style={font=\scriptsize},
  ylabel={元/股}, ylabel style={font=\scriptsize},
  grid=major, grid style={LightGray, line width=0.3pt},
  axis line style={InkGray, line width=0.5pt},
  every axis plot/.append style={line width=0.9pt},
]
\addplot[AgriGreen] table[x=x,y=y]{data/clean/profiles/px_002385.dat};
\addplot[SoftRed, dashed, line width=0.7pt] table[x=x,y=y]{data/clean/profiles/pxzz_002385.dat};
\addplot[only marks, mark=*, mark size=1.1pt, SoftRed] table[x=x,y=y]{data/clean/profiles/pxzz_002385.dat};
\end{axis}
\end{tikzpicture}

\captionof{figure}[大北农（002385）股价与周期骨架]{大北农（002385）月度收盘价（元/股）与 ZigZag 周期骨架（阈值 25\%，圆点为周期转折点）}
\end{minipage}
\end{center}

\pfl{过去周期分析} 以 25\% 的 ZigZag 阈值划分，样本区间内股价形成 2 个主升段与 3 个主跌段。 最大主升段出现在 2021 年 6 月 至 2021 年 12 月，6 个月+54.3\%；最大主跌段为 2022 年 11 月 至 2026 年 6 月，43 个月-70.4\%。 最近一次转折出现在 2026 年 6 月（下跌段自 2022 年 11 月起，43 个月-70.4\%），当前处于该段的延续之中。 公司股价与申万农林牧渔指数几乎同步见底，个股并未走出独立行情。 饲料夹在原料成本与养殖需求之间（第 \ref{sec:feed} 节），量的部分取决于存栏、利的部分取决于原料与成品的价差，公司 2026 年 6 月末资产负债率 69.9\%，高于全样本中位数 52.2\%，其股价因此更多反映资金面与信用风险，而不是价格弹性本身。 综合区间形态看，该股单段涨跌幅度偏大、以趋势段为主（最大主升 +54.3\%、最大主跌 -70.4\%），年化波动率 33.3\% 与全样本中位数 37.4\% 相比波动小于样本中位数。

\pfl{盈利情况} 2026 年上半年营业收入 138.11 亿元（同比 +1.9\%），归母净利润 -6.77 亿元，由上年同期的盈利转为亏损。 以年报为趋势对照，2023---2025 年营业收入由 333.90 亿元变为 291.19 亿元（两年年均复合增速 -6.6\%），归母净利润由 -21.74 亿元变为 -6.42 亿元。 按半年报口径，2026 年上半年的 ROE 为 -9.3\%（未年化）、销售净利率 -6.0\%、毛利率 8.7\%，ROE 在“畜禽饲料”11 家公司中按数值由高到低排第\zhnumber{9}位。 同期全样本半年 ROE 中位数 0.75\%，公司与之相差 10.04 个百分点（弱于中位数公司）。 从公司披露的原因看，业绩变动主要来自：2026 年 8 月，公司投资者关系活动披露：以饲料、种业、动保、食品多板块协同应对周期波动。 面对这一局面，公司披露的举措包括：2026 年 7 月，实施股份回购，并对部分限制性股票回购注销、减少注册资本。 综合看，公司仍处亏损区间、由盈转亏，半年 ROE 在三级行业内按由高到低排第\zhnumber{9}位，单位盈利能力的修复依赖行业景气与自身成本端的改善，报表弹性主要体现为价格与销量的方向性变化。

\begin{center}\small\setlength{\tabcolsep}{3.4pt}
\captionof{table}[大北农（002385）关键财务指标]{大北农（002385）关键财务指标（合并报表口径；两个半年报期均未年化；现金短债比 $=$ 货币资金 $\div$ 短期债务）}
\begin{adjustbox}{max width=\textwidth}
\begin{tabular}{@{}lrrrrr@{}}
\toprule
指标 & 2023 年 & 2024 年 & 2025 年 & 2025 年半年报 & 2026 年半年报 \\
\midrule
营业收入（亿元） & 333.90 & 287.67 & 291.19 & 135.59 & 138.11 \\
归母净利润（亿元） & -21.74 & 3.46 & -6.42 & 2.35 & -6.77 \\
毛利率（\%） & 9.6 & 15.1 & 12.7 & 14.1 & 8.7 \\
ROE（\%） & -23.0 & 4.2 & -7.8 & 2.7 & -9.3 \\
资产负债率（\%） & 66.3 & 63.5 & 67.0 & 64.0 & 69.9 \\
经营活动现金流净额（亿元） & 2.45 & 24.83 & 8.29 & 0.24 & -12.25 \\
货币资金（亿元） & 46.41 & 38.24 & 42.77 & 45.62 & 32.81 \\
有息负债合计（亿元） & 131.20 & 115.30 & 133.70 & 130.82 & 135.45 \\
现金短债比（倍） & 0.46 & 0.42 & 0.41 & 0.45 & 0.31 \\
\bottomrule
\end{tabular}
\end{adjustbox}
\end{center}

\begin{center}
\begin{minipage}{\textwidth}\centering
\begin{tikzpicture}
\begin{groupplot}[
  group style={group size=2 by 1, horizontal sep=0.60cm},
  width=6.85cm, height=4.60cm,
  tick label style={font=\tiny},
  title style={font=\scriptsize\heiti},
  yticklabel style={font=\tiny},
  ylabel style={font=\tiny},
  grid=major, grid style={LightGray, line width=0.3pt},
  axis line style={InkGray, line width=0.5pt},
  enlarge x limits={abs=0.8},
  legend style={font=\scriptsize, at={(0.5,-0.20)}, anchor=north,
                legend columns=2, draw=none, fill=none, column sep=6pt},
]
\nextgroupplot[
  ybar, bar width=4pt,
  ymin=-57.30, ymax=376.58,
  xtick={0,1,2,3,4,5.7},
  xticklabels={2021,2022,2023,2024,2025,2026H1},
  ylabel={亿元},
]
\addplot[fill=AgriGreen!75, draw=AgriGreen, bar shift=-2.6pt] table[x=i, y=rev]{data/clean/profiles/fin_002385.dat};
\addplot[fill=SoftRed!60, draw=SoftRed, bar shift=2.6pt] table[x=i, y=ni]{data/clean/profiles/fin_002385.dat};
\addplot[fill=AgriGreen!30, draw=AgriGreen, pattern=north east lines, bar shift=-2.6pt] table[x=x, y=rev]{data/clean/profiles/finh_002385.dat};
\addplot[fill=SoftRed!20, draw=SoftRed, pattern=north east lines, bar shift=2.6pt] table[x=x, y=ni]{data/clean/profiles/finh_002385.dat};
\legend{营业收入（年/半年）, 归母净利润（年/半年）}
\nextgroupplot[
  ymin=-30.4, ymax=79.2,
  xtick={0,1,2,3,4,5.7},
  xticklabels={2021,2022,2023,2024,2025,2026H1},
  ylabel={\%},
]
\addplot[AgriGold, mark=*, mark size=1.3pt] table[x=i, y=roe]{data/clean/profiles/fin_002385.dat};
\addplot[OilBlue, mark=square*, mark size=1.3pt] table[x=i, y=debt]{data/clean/profiles/fin_002385.dat};
\addplot[only marks, AgriGold, mark=*, mark size=2.0pt] table[x=x, y=roe]{data/clean/profiles/finh_002385.dat};
\addplot[only marks, OilBlue, mark=square*, mark size=2.0pt] table[x=x, y=debt]{data/clean/profiles/finh_002385.dat};
\legend{ROE, 资产负债率}
\end{groupplot}
\end{tikzpicture}
\begin{tikzpicture}
\begin{axis}[
  name=finbot,
  width=13.75cm, height=3.10cm, scale only axis,
  ymin=-13.55, ymax=151.71,
  xtick={0,1,2,3,4,5.7},
  xticklabels={2021,2022,2023,2024,2025,2026H1},
  tick label style={font=\tiny},
  yticklabel style={font=\tiny},
  ylabel={亿元}, ylabel style={font=\tiny},
  ybar, bar width=4pt,
  grid=major, grid style={LightGray, line width=0.3pt},
  axis line style={InkGray, line width=0.5pt},
  enlarge x limits={abs=0.8},
  legend style={font=\scriptsize, at={(0.5,-0.26)}, anchor=north,
                legend columns=2, draw=none, fill=none, column sep=10pt},
]
\addplot[fill=AgriGreen!55, draw=AgriGreen, bar shift=-3.0pt] table[x=i, y=cash]{data/clean/profiles/fin_002385.dat};
\addplot[fill=SoftRed!45, draw=SoftRed, bar shift=3.0pt] table[x=i, y=ibd]{data/clean/profiles/fin_002385.dat};
\addplot[fill=AgriGreen!25, draw=AgriGreen, pattern=north east lines, bar shift=-3.0pt] table[x=x, y=cash]{data/clean/profiles/finh_002385.dat};
\addplot[fill=SoftRed!20, draw=SoftRed, pattern=north east lines, bar shift=3.0pt] table[x=x, y=ibd]{data/clean/profiles/finh_002385.dat};
\legend{货币资金（左轴）, 有息负债（左轴）}
\end{axis}
\begin{axis}[
  at={(finbot.south east)}, anchor=south east,
  width=13.75cm, height=3.10cm, scale only axis,
  axis y line*=right, axis x line*=none, xtick=\empty,
  ymin=-21.89, ymax=35.95,
  tick label style={font=\tiny},
  yticklabel style={font=\tiny}, scaled y ticks=false,
  legend style={font=\scriptsize, at={(0.5,-0.44)}, anchor=north,
                legend columns=1, draw=none, fill=none},
]
\addplot[OilBlue, mark=*, mark size=2.2pt, line width=1.3pt] table[x=i, y=ocf]{data/clean/profiles/fin_002385.dat};
\addplot[only marks, OilBlue, mark=*, mark size=3.2pt] table[x=x, y=ocf]{data/clean/profiles/finh_002385.dat};
\legend{经营活动现金流净额（右轴，亿元，蓝点）}
\end{axis}
\end{tikzpicture}

\captionof{figure}[大北农（002385）近年主要财务指标]{大北农（002385）近年主要财务指标（左上：营业收入与归母净利润；右上：ROE 与资产负债率；下：货币资金、有息负债为左轴柱状，经营活动现金流净额为其右轴的蓝点折线。
2021---2025 年为年报口径，2026H1 为 2026 年半年报/半年末，与其前一年报之间留出间隔、以斜纹或空心点区分；半年报数据未年化）}
\end{minipage}
\end{center}

\pfl{财务分析} 2026 年 6 月末资产负债率 69.9\%，较 2025 年末上升 2.9 个百分点（样本中位数 52.2\%，所处三级行业内按由低到高排第\zhnumber{8}位）。 2026 年 6 月末货币资金 32.81 亿元、短期债务（短期借款与一年内到期非流动负债）104.33 亿元、长期债务 31.12 亿元，有息负债合计 135.45 亿元，现金短债比 0.31 倍——覆盖偏紧。 净有息负债 102.64 亿元，相当于其 2026 年上半年营业收入的 74\%。 流动比率 0.73、速动比率 0.43，短期偿债指标偏紧。 2026 年上半年经营活动现金流净额 -12.25 亿元，经营现金流与利润同时为负，日常运营对债务与股东投入的依赖度较高。 2026 年 6 月末在建工程 0.84 亿元，较上年末减少 0.05 亿元，在建投入在收缩。 2026 年 6 月末商誉 6.40 亿元，占归母净资产的 7.5\%，是净资产中最脆弱的一块。 综合看，账面货币资金对有息负债与短期债务的覆盖不足，滚动偿付对新增融资的依赖度较高；资产负债率高于样本中位数 17.7 个百分点；有息负债中短期债务占比更高，期限结构仍以滚动续作为主。 公司披露的股东与质押风险为：2026 年 2 月，其名下约 9.28 亿股中有近 5 亿股处于质押状态；2026 年 9 月，控股股东莫云将所持股份质押给质权人以置换并解除邵根伙名下同等数量质押股份。

\pfl{重大战略转型} 公司的战略动作可归为以下几类：其他动作（3 项）：2026 年 7 月，公司参股公司增资暨关联交易。 公告口径上，近三年（2023 年 9 月---2026 年 9 月）公司披露重大重组与并购类公告 1 条、对外投资与转型类公告 6 条、回购与股权激励类公告 38 条，合计公告 604 条。 第一大产品仍是“饲料产品”，收入占比 67.7\%，与 2021 年年报相比没有方向性变化。 综合判断，报告期内公司有 7 条与战略相关的公告落地（重大重组与并购 1 条、对外投资与转型 6 条），资本运作与主业调整之间存在直接关联。

\begin{center}\small\setlength{\tabcolsep}{4pt}
\captionof{table}[大北农（002385）历史融资与资本运作]{大北农（002385）历史融资与资本运作（据公司公告与公开报道整理；金额为披露值，含首次公开发行、再融资、债券、银行借款与重整投资等）}
\begin{adjustbox}{max width=\textwidth}
\begin{tabular}{@{}p{1.45cm}p{2.55cm}p{4.05cm}p{4.95cm}@{}}
\toprule
时间 & 方式 & 规模 & 用途与说明 \\
\midrule
2010-03 & IPO（深交所） & 6,\allowbreak{}080 万股、\allowbreak{}35.00 元/\allowbreak{}股，\allowbreak{}募集资金净额 20.18 亿元（20.18 亿元） & 新型高效预混料等项目；2010-04-09 上市；截至 2023 年首发后累计派现金额亦在该文件中列示 \\
2014-06 & 公司债券 & 5.00 亿元，\allowbreak{}“14 北农债”（112203），\allowbreak{}票面利率 7.24\% & 2014-06-13 上市，\allowbreak{}发行完成后累计公司债券余额不超过 5.00 亿元 \\
2023—2024 & 向特定对象发行股票（定向增发） & 2.11 亿股、\allowbreak{}3.31 元/\allowbreak{}股，\allowbreak{}募集资金总额 7.00 亿元（7.00 亿元） & 饲料项目建设、\allowbreak{}补充流动资金等（原预案拟募集不超过 19.43 亿元 \\
2025-01 & 募投项目变更（募集资金内部调整） & 终止“大北农辽宁区核心科技园建设项目”，\allowbreak{}节余募集资金 1.51 亿元（1.51 亿元） & 全部改投饲料项目、\allowbreak{}技改项目和信息化建设项目；2025 年 1 月 24 日董事会、\allowbreak{}2025 年 2 月 10 日股东大会通过 \\
\bottomrule
\end{tabular}
\end{adjustbox}
\end{center}

\pfl{历史融投资情况} 从现金流量表看，2025 年公司取得借款收到的现金 101.26 亿元、偿还债务支付的现金 86.72 亿元、吸收权益性投资 0.12 亿元，筹资活动现金流净额 4.09 亿元（2023 年为取得借款 85.89 亿元、偿还债务 68.23 亿元）。 2026 年上半年筹资活动现金流净额为 -0.25 亿元（取得借款 50.29 亿元、偿还债务 49.38 亿元，未年化），已转为净流出，处于净还债状态。 2025 年分配股利、利润或偿付利息支付的现金合计 6.03 亿元。 公开披露的融资记录共 4 笔，工具类型以 IPO、公司债、定向增发为主；金额与用途见下表。 其中规模最大的两笔为：2025 年 1 月，募投项目变更，终止“大北农辽宁区核心科技园建设项目”，节余募集资金 1.51 亿元（1.51 亿元）（该事项后续已终止或解除，不构成已到位资金）；2023---2024 年，定向增发，2.11 亿股、3.31 元/股，募集资金总额 7.00 亿元（7.00 亿元）。 股本方面，近三年披露股本变动 8 次（股份回购 3 次、增发新股上市 1 次、其他 1 次）；总股本由 2021 年末的 41.41 亿股变为 43.24 亿股（+4.41\%）。 分红方面，近三年现金分红 4 次、合计每股派现 0.19 元，最近一次实施在 2025 年 9 月。 近三年“再融资”类公告 16 条，涉及定向增发、募集资金使用。 综合看，公司融资结构上以债务滚动为主、股权融资为补充；2025 年筹资活动现金流净额 4.09 亿元，年度内仍为净融入，与前述经营与投资现金流的方向基本一致。

\pfl{未来潜在融资需求分析} 2026 年 6 月末货币资金 32.81 亿元，而短期债务 104.33 亿元（现金短债比 0.31 倍），2026 年上半年经营现金流净额 -12.25 亿元。按“短期债务减去账面货币资金”的滚动偿付口径，静态缺口约 71.52 亿元；若再计入上半年亏损的年化额（13.53 亿元），维持一年正常经营所需的外部资金约 71.52---85.05 亿元。 按第 \ref{sec:financing-need} 节的三档划分，公司属于第一档“补血型”：2026 年上半年归母净利润 -6.77 亿元、6 月末资产负债率 69.9\%。 可行的融资路径按现实性排序为：债务重组或展期（与银行和债券持有人协商）、引入国资或产业投资人（部分公司以控制权让渡为对价）、出售非核心资产与参股股权。 公开披露的后续资金安排线索包括：股权融资线索：2022 年度，公司向特定对象发行股票原拟募集不超过 19.43 亿元；债务与授信安排：2026 年 9 月，发布公司及控股子公司担保进展公告，为子公司经营融资提供担保；资本开支安排：2025 年 10 月，董事会通过部分募集资金投资项目延期。 资金需求还要叠加在建工程：期末在建工程 0.84 亿元、2026 年上半年购建固定资产等支出 2.25 亿元（未年化），这部分支出在周期底部通常会被主动放缓。 综合看，公司的外部资金需求以债务接续为主，滚动偿付口径下的静态缺口约 71.52 亿元，融资节奏将受制于银行续贷意愿与股权融资窗口。




\subsubsection{天康生物（002100）}
\label{co:002100}

\pfl{公司基本情况} 公司前身可追溯至 2006 年，新疆兵团国资背景的农牧企业，主营饲料、生猪养殖与食品、动物疫苗；2026 年 6 月收购羌都畜牧 51\% 股权扩产生猪，2026 年上半年亏损 4.41 亿元。 公司于 2006-12-26 在深交所主板首发上市，注册资本 13.65 亿元，总股本 13.65 亿股。 截至 2026 年 9 月 24 日收盘价 7.58 元，流通市值 103.5 亿元，近一年-1.6\%，流通市值在三级行业“畜禽饲料”的 11 家公司中按由大到小排第\zhnumber{3}位，近一年涨跌幅高于全样本中位数（-10.1\%）8.5 个百分点。 按 2026 年半年报披露的行业构成，收入占比前三位为生猪养殖产业链行业 32.4\%；饲料行业 27.8\%；蛋白油脂加工行业 16.7\%。 与 2021 年年报相比，第一大行业口径由“饲料行业”（32.6\%）变为“生猪养殖产业链行业”（32.4\%）；两期披露的分类口径有调整，占比差异中同时包含分类因素。

\pfl{历史股价} 自 2021 年 1 月至 2026 年 9 月，股价由 10.76 元变为 7.58 元，累计 -29.6\%，区间最高 12.41 元（2021 年 2 月）、最低 6.14 元（2025 年 2 月）；当前价相当于区间最高价的 61.1\%，为区间最低价的 1.23 倍。 月度收益的年化波动率 33.8\%，低于全样本中位数 37.4\%，在三级行业“畜禽饲料”的 11 家公司中波动率由高到低排第\zhnumber{7}位。

\begin{center}
\begin{minipage}{\textwidth}\centering
\begin{tikzpicture}
\begin{axis}[
  width=12.6cm, height=3.9cm, scale only axis,
  xmin=2020.922, xmax=2026.888, ymin=5.71, ymax=13.28,
  xtick={2021,2022,2023,2024,2025,2026},
  yearticks,
  yticklabel style={font=\scriptsize},
  ylabel={元/股}, ylabel style={font=\scriptsize},
  grid=major, grid style={LightGray, line width=0.3pt},
  axis line style={InkGray, line width=0.5pt},
  every axis plot/.append style={line width=0.9pt},
]
\addplot[AgriGreen] table[x=x,y=y]{data/clean/profiles/px_002100.dat};
\addplot[SoftRed, dashed, line width=0.7pt] table[x=x,y=y]{data/clean/profiles/pxzz_002100.dat};
\addplot[only marks, mark=*, mark size=1.1pt, SoftRed] table[x=x,y=y]{data/clean/profiles/pxzz_002100.dat};
\end{axis}
\end{tikzpicture}

\captionof{figure}[天康生物（002100）股价与周期骨架]{天康生物（002100）月度收盘价（元/股）与 ZigZag 周期骨架（阈值 25\%，圆点为周期转折点）}
\end{minipage}
\end{center}

\pfl{过去周期分析} 以 25\% 的 ZigZag 阈值划分，样本区间内股价形成 3 个主升段与 3 个主跌段。 最大主升段出现在 2021 年 7 月 至 2022 年 3 月，8 个月+66.9\%；最大主跌段为 2022 年 7 月 至 2025 年 2 月，31 个月-49.4\%。 最近一次转折出现在 2026 年 4 月（上涨段自 2025 年 2 月起，14 个月+29.5\%），当前处于该段之后的震荡之中。 公司股价较申万农林牧渔指数提前约 15 个月见底，是板块中较早定价周期反转的公司之一。 饲料夹在原料成本与养殖需求之间（第 \ref{sec:feed} 节），量的部分取决于存栏、利的部分取决于原料与成品的价差，公司 2026 年 6 月末资产负债率 55.6\%、半年 ROE -6.6\%（未年化），在本轮下行中属于随行业同步调整的一类。 综合区间形态看，该股单段涨跌幅度偏大、以趋势段为主（最大主升 +66.9\%、最大主跌 -49.4\%），年化波动率 33.8\% 与全样本中位数 37.4\% 相比波动与样本中位数接近。

\pfl{盈利情况} 2026 年上半年营业收入 83.25 亿元（同比 -5.9\%），归母净利润 -4.41 亿元，由上年同期的盈利转为亏损。 以年报为趋势对照，2023---2025 年营业收入由 190.26 亿元变为 169.81 亿元（两年年均复合增速 -5.5\%），归母净利润由 -13.63 亿元变为 2.11 亿元。 按半年报口径，2026 年上半年的 ROE 为 -6.6\%（未年化）、销售净利率 -5.9\%、毛利率 1.4\%，ROE 在“畜禽饲料”11 家公司中按数值由高到低排第\zhnumber{7}位。 同期全样本半年 ROE 中位数 0.75\%，公司与之相差 7.32 个百分点（弱于中位数公司）。 从公司披露的原因看，业绩变动主要来自：2026 年上半年，兽药（动物疫苗）收入同比下降 43.95\%，为除猪价外的主要拖累项。 针对上述压力，公司披露的应对举措包括：2026 年 6 月，公司表示本次收购整合以现有存量产能优化配置为核心。 综合看，公司仍处亏损区间、由盈转亏，半年 ROE 在三级行业内按由高到低排第\zhnumber{7}位，单位盈利能力的修复依赖行业景气与自身成本端的改善，报表弹性主要体现为价格与销量的方向性变化。

\begin{center}\small\setlength{\tabcolsep}{3.4pt}
\captionof{table}[天康生物（002100）关键财务指标]{天康生物（002100）关键财务指标（合并报表口径；两个半年报期均未年化；现金短债比 $=$ 货币资金 $\div$ 短期债务）}
\begin{adjustbox}{max width=\textwidth}
\begin{tabular}{@{}lrrrrr@{}}
\toprule
指标 & 2023 年 & 2024 年 & 2025 年 & 2025 年半年报 & 2026 年半年报 \\
\midrule
营业收入（亿元） & 190.26 & 171.76 & 169.81 & 88.47 & 83.25 \\
归母净利润（亿元） & -13.63 & 6.05 & 2.11 & 3.38 & -4.41 \\
毛利率（\%） & 2.9 & 12.1 & 9.9 & 12.1 & 1.4 \\
ROE（\%） & -19.2 & 9.0 & 3.0 & 4.7 & -6.6 \\
资产负债率（\%） & 52.8 & 51.5 & 50.9 & 47.0 & 55.6 \\
经营活动现金流净额（亿元） & 17.52 & 11.62 & 3.81 & 16.19 & 14.96 \\
货币资金（亿元） & 28.22 & 31.02 & 20.81 & 29.73 & 25.35 \\
有息负债合计（亿元） & 63.96 & 64.79 & 63.68 & 52.26 & 65.33 \\
现金短债比（倍） & 0.59 & 0.76 & 0.50 & 0.90 & 0.83 \\
\bottomrule
\end{tabular}
\end{adjustbox}
\end{center}

\begin{center}
\begin{minipage}{\textwidth}\centering
\begin{tikzpicture}
\begin{groupplot}[
  group style={group size=2 by 1, horizontal sep=0.60cm},
  width=6.85cm, height=4.60cm,
  tick label style={font=\tiny},
  title style={font=\scriptsize\heiti},
  yticklabel style={font=\tiny},
  ylabel style={font=\tiny},
  grid=major, grid style={LightGray, line width=0.3pt},
  axis line style={InkGray, line width=0.5pt},
  enlarge x limits={abs=0.8},
  legend style={font=\scriptsize, at={(0.5,-0.20)}, anchor=north,
                legend columns=2, draw=none, fill=none, column sep=6pt},
]
\nextgroupplot[
  ybar, bar width=4pt,
  ymin=-34.02, ymax=214.73,
  xtick={0,1,2,3,4,5.7},
  xticklabels={2021,2022,2023,2024,2025,2026H1},
  ylabel={亿元},
]
\addplot[fill=AgriGreen!75, draw=AgriGreen, bar shift=-2.6pt] table[x=i, y=rev]{data/clean/profiles/fin_002100.dat};
\addplot[fill=SoftRed!60, draw=SoftRed, bar shift=2.6pt] table[x=i, y=ni]{data/clean/profiles/fin_002100.dat};
\addplot[fill=AgriGreen!30, draw=AgriGreen, pattern=north east lines, bar shift=-2.6pt] table[x=x, y=rev]{data/clean/profiles/finh_002100.dat};
\addplot[fill=SoftRed!20, draw=SoftRed, pattern=north east lines, bar shift=2.6pt] table[x=x, y=ni]{data/clean/profiles/finh_002100.dat};
\legend{营业收入（年/半年）, 归母净利润（年/半年）}
\nextgroupplot[
  ymin=-25.2, ymax=63.0,
  xtick={0,1,2,3,4,5.7},
  xticklabels={2021,2022,2023,2024,2025,2026H1},
  ylabel={\%},
]
\addplot[AgriGold, mark=*, mark size=1.3pt] table[x=i, y=roe]{data/clean/profiles/fin_002100.dat};
\addplot[OilBlue, mark=square*, mark size=1.3pt] table[x=i, y=debt]{data/clean/profiles/fin_002100.dat};
\addplot[only marks, AgriGold, mark=*, mark size=2.0pt] table[x=x, y=roe]{data/clean/profiles/finh_002100.dat};
\addplot[only marks, OilBlue, mark=square*, mark size=2.0pt] table[x=x, y=debt]{data/clean/profiles/finh_002100.dat};
\legend{ROE, 资产负债率}
\end{groupplot}
\end{tikzpicture}
\begin{tikzpicture}
\begin{axis}[
  name=finbot,
  width=13.75cm, height=3.10cm, scale only axis,
  ymin=-7.38, ymax=82.62,
  xtick={0,1,2,3,4,5.7},
  xticklabels={2021,2022,2023,2024,2025,2026H1},
  tick label style={font=\tiny},
  yticklabel style={font=\tiny},
  ylabel={亿元}, ylabel style={font=\tiny},
  ybar, bar width=4pt,
  grid=major, grid style={LightGray, line width=0.3pt},
  axis line style={InkGray, line width=0.5pt},
  enlarge x limits={abs=0.8},
  legend style={font=\scriptsize, at={(0.5,-0.26)}, anchor=north,
                legend columns=2, draw=none, fill=none, column sep=10pt},
]
\addplot[fill=AgriGreen!55, draw=AgriGreen, bar shift=-3.0pt] table[x=i, y=cash]{data/clean/profiles/fin_002100.dat};
\addplot[fill=SoftRed!45, draw=SoftRed, bar shift=3.0pt] table[x=i, y=ibd]{data/clean/profiles/fin_002100.dat};
\addplot[fill=AgriGreen!25, draw=AgriGreen, pattern=north east lines, bar shift=-3.0pt] table[x=x, y=cash]{data/clean/profiles/finh_002100.dat};
\addplot[fill=SoftRed!20, draw=SoftRed, pattern=north east lines, bar shift=3.0pt] table[x=x, y=ibd]{data/clean/profiles/finh_002100.dat};
\legend{货币资金（左轴）, 有息负债（左轴）}
\end{axis}
\begin{axis}[
  at={(finbot.south east)}, anchor=south east,
  width=13.75cm, height=3.10cm, scale only axis,
  axis y line*=right, axis x line*=none, xtick=\empty,
  ymin=-4.55, ymax=22.77,
  tick label style={font=\tiny},
  yticklabel style={font=\tiny}, scaled y ticks=false,
  legend style={font=\scriptsize, at={(0.5,-0.44)}, anchor=north,
                legend columns=1, draw=none, fill=none},
]
\addplot[OilBlue, mark=*, mark size=2.2pt, line width=1.3pt] table[x=i, y=ocf]{data/clean/profiles/fin_002100.dat};
\addplot[only marks, OilBlue, mark=*, mark size=3.2pt] table[x=x, y=ocf]{data/clean/profiles/finh_002100.dat};
\legend{经营活动现金流净额（右轴，亿元，蓝点）}
\end{axis}
\end{tikzpicture}

\captionof{figure}[天康生物（002100）近年主要财务指标]{天康生物（002100）近年主要财务指标（左上：营业收入与归母净利润；右上：ROE 与资产负债率；下：货币资金、有息负债为左轴柱状，经营活动现金流净额为其右轴的蓝点折线。
2021---2025 年为年报口径，2026H1 为 2026 年半年报/半年末，与其前一年报之间留出间隔、以斜纹或空心点区分；半年报数据未年化）}
\end{minipage}
\end{center}

\pfl{财务分析} 2026 年 6 月末资产负债率 55.6\%，较 2025 年末上升 4.6 个百分点（样本中位数 52.2\%，所处三级行业内按由低到高排第\zhnumber{3}位）。 2026 年 6 月末货币资金 25.35 亿元、短期债务（短期借款与一年内到期非流动负债）30.51 亿元、长期债务 34.83 亿元，有息负债合计 65.33 亿元，现金短债比 0.83 倍——覆盖偏紧。 净有息负债 39.99 亿元，相当于其 2026 年上半年营业收入的 48\%。 流动比率 1.31、速动比率 0.68，短期流动性无明显缺口。 2026 年上半年经营活动现金流净额 14.96 亿元，净利润为负而经营现金流为正，差额主要来自折旧摊销与营运资本变动，说明亏损尚未完全转化为现金流出。 2026 年 6 月末在建工程 5.64 亿元，较上年末增加 3.09 亿元，仍有在建投入。 2026 年 6 月末商誉 7.20 亿元，占归母净资产的 8.4\%，是净资产中最脆弱的一块。 综合看，现金短债比 0.83 倍，短期资金安排偏紧；资产负债率高于样本中位数 3.4 个百分点；有息负债中长期债务占比更高，期限结构对短债滚动的压力较小。

\pfl{重大战略转型} 从公开披露看，公司的战略推进集中在：并购与重组（3 项）：2026 年 1 月，国家市场监督管理总局公示天康生物收购新疆羌都畜牧科技有限公司股权案（经营者集中简易案件）；2025 年 12 月，与新疆七星羌都集团农牧有限公司等股东签署股权收购协议，以现金 12.75 亿元收购羌都畜牧 51\% 股权；融资与资本运作（1 项）：2025 年 6 月，子公司天康制药推进分拆上市：招股说明书显示天康生物直接持股 58.9388\%。 公告口径上，近三年（2023 年 9 月---2026 年 9 月）公司披露重大重组与并购类公告 5 条、对外投资与转型类公告 1 条、回购与股权激励类公告 0 条，合计公告 346 条。 主营结构的口径变化已经落到报表上：2021 年年报的第一大行业为“饲料行业”（32.6\%），2026 年半年报为“生猪养殖产业链行业”（32.4\%）；公司在这两期的分部划分方式有过调整。 综合判断，报告期内公司有 6 条与战略相关的公告落地（重大重组与并购 5 条、对外投资与转型 1 条），资本运作与主业调整之间存在直接关联。

\begin{center}\small\setlength{\tabcolsep}{4pt}
\captionof{table}[天康生物（002100）历史融资与资本运作]{天康生物（002100）历史融资与资本运作（据公司公告与公开报道整理；金额为披露值，含首次公开发行、再融资、债券、银行借款与重整投资等）}
\begin{adjustbox}{max width=\textwidth}
\begin{tabular}{@{}p{1.45cm}p{2.55cm}p{4.05cm}p{4.95cm}@{}}
\toprule
时间 & 方式 & 规模 & 用途与说明 \\
\midrule
2006-12 & IPO（深交所） & 1,\allowbreak{}600 万股、\allowbreak{}10.86 元/\allowbreak{}股 & 牛羊专用饲料生产线、\allowbreak{}专业浮性水产饲料厂、\allowbreak{}酶制剂和益绿素纯植物饲料添加剂等项目；上市日 2006-12-26 \\
2015-07 & 发行股份吸收合并（关联交易） & 发行 3.88 亿股，\allowbreak{}同时注销天康控股所持 2.93 亿股 & 吸收合并新疆天康（控股）集团 \\
2017-12 & 公开发行可转换公司债券 & 10.00 亿元（1,\allowbreak{}000 万张，\allowbreak{}面值 100.00 元，\allowbreak{}期限 6 年） & 联合信用评级给予主体与债项 AA；扣除承销费及保荐费 1，\allowbreak{}450.00 万元后实际收到认购资金 9.86 亿元 \\
2021-12 & 非公开发行股票（2020 年度方案） & 2.77 亿股、\allowbreak{}7.45 元/\allowbreak{}股，\allowbreak{}募集资金总额 20.67 亿元，\allowbreak{}净额 20.43 亿元 & 生猪养殖及饲料等项目；2021-12-24 新增股份上市；2025 年报显示该次募集资金累计已大量使用 \\
2026 & 综合授信额度申请 & 公司及下属公司 2026 年度拟向银行及其他类金融机构申请总额不超过 120.00 亿元综合授信 & 建设项目贷款、\allowbreak{}流动资金贷款及原材料采购等用于降低原材料价格波动风险、\allowbreak{}提高各业务板块盈利能力 \\
\bottomrule
\end{tabular}
\end{adjustbox}
\end{center}

\pfl{历史融投资情况} 从现金流量表看，2025 年公司取得借款收到的现金 53.64 亿元、偿还债务支付的现金 54.08 亿元、吸收权益性投资 0.00 亿元，筹资活动现金流净额 -6.36 亿元（2023 年为取得借款 49.15 亿元、偿还债务 58.82 亿元）。 2026 年上半年筹资活动现金流净额为 -3.13 亿元（取得借款 26.99 亿元、偿还债务 28.40 亿元，未年化），已转为净流出，处于净还债状态。 2025 年分配股利、利润或偿付利息支付的现金合计 4.43 亿元。 公开披露的融资记录共 5 笔，工具类型以 IPO、发行股份吸收合并、可转债、定向增发等为主；金额与用途见下表。 其中规模最大的两笔为：2026 年，银行授信，公司及下属公司 2026 年度拟向银行及其他类金融机构申请总额不超过 120.00 亿元综合授信；2021 年 12 月，定向增发，2.77 亿股、7.45 元/股，募集资金总额 20.67 亿元，净额 20.43 亿元。 股本方面，近三年披露股本变动 7 次（可转债转股 5 次）；总股本由 2021 年末的 13.54 亿股变为 13.65 亿股（+0.84\%）。 分红方面，近三年现金分红 3 次、合计每股派现 0.52 元，最近一次实施在 2025 年 12 月。 近三年“再融资”类公告 2 条，涉及可转债。 综合看，公司融资结构上以债务滚动为主、股权融资为补充；2025 年筹资活动现金流净额 -6.36 亿元，融资节奏仍以维持存量债务周转为主，与前述经营与投资现金流的方向基本一致。

\pfl{未来潜在融资需求分析} 2026 年 6 月末货币资金 25.35 亿元，而短期债务 30.51 亿元（现金短债比 0.83 倍），2026 年上半年经营现金流净额 14.96 亿元。按“短期债务减去账面货币资金”的滚动偿付口径，静态缺口约 5.16 亿元；若再计入上半年亏损的年化额（8.83 亿元），维持一年正常经营所需的外部资金约 5.16---13.99 亿元。 公司 2026 年上半年归母净利润为负（-4.41 亿元），但 6 月末资产负债率 55.6\% 尚未进入第一档（亏损且负债率 $>$65\%）的区间，当前融资需求以补充流动资金、支撑低谷期经营性支出为主。 公开披露的后续资金安排线索包括：股权融资线索：2025---2026 年，天康制药分拆上市持续推进，为子公司股权融资线索；债务与授信安排：2026 年 4 月，董事会/年报披露 2026 年度资金需求计划：拟向银行及其他类金融机构申请总额不超过 120.00 亿元综合授信。 资金需求还要叠加在建工程：期末在建工程 5.64 亿元、2026 年上半年购建固定资产等支出 1.52 亿元（未年化），这部分支出在周期底部通常会被主动放缓。 综合看，公司的外部资金需求以补充营运资金为主、项目投入为辅，滚动偿付口径下的静态缺口约 5.16 亿元，能否落地取决于授信续作与发行窗口的配合。




\subsubsection{天马科技（603668）}
\label{co:603668}

\pfl{公司基本情况} 2017 年设立至今，天马科技是福建特种水产饲料（鳗鲡料等）起家的农牧食品企业，形成“以鳗鱼为核心、以食品为新蓝海、以饲料为主基石”格局；2026 年上半年鳗价低位致净利降 79.53\%。 公司于 2017-01-17 在上交所首发上市，注册资本 5.06 亿元，总股本 5.02 亿股。 截至 2026 年 9 月 24 日收盘价 12.18 元，流通市值 61.6 亿元，近一年-22.0\%，流通市值在三级行业“水产饲料”的 4 家公司中按由大到小排第\zhnumber{2}位，近一年涨跌幅低于全样本中位数（-10.1\%）11.9 个百分点。 按 2025 年年报披露的行业构成，收入占比前三位为饲料行业 87.4\%；养殖及食品行业 31.1\%；其他(补充) 0.5\%。

\pfl{历史股价} 自 2021 年 1 月至 2026 年 9 月，股价由 6.91 元变为 12.18 元，累计 +76.3\%，区间最高 21.35 元（2022 年 7 月）、最低 6.91 元（2021 年 1 月）；当前价相当于区间最高价的 57.0\%，为区间最低价的 1.76 倍。 月度收益的年化波动率 44.9\%，高于全样本中位数 37.4\%，在三级行业“水产饲料”的 4 家公司中波动率由高到低排第\zhnumber{1}位。

\begin{center}
\begin{minipage}{\textwidth}\centering
\begin{tikzpicture}
\begin{axis}[
  width=12.6cm, height=3.9cm, scale only axis,
  xmin=2020.922, xmax=2026.888, ymin=6.43, ymax=22.84,
  xtick={2021,2022,2023,2024,2025,2026},
  yearticks,
  yticklabel style={font=\scriptsize},
  ylabel={元/股}, ylabel style={font=\scriptsize},
  grid=major, grid style={LightGray, line width=0.3pt},
  axis line style={InkGray, line width=0.5pt},
  every axis plot/.append style={line width=0.9pt},
]
\addplot[AgriGreen] table[x=x,y=y]{data/clean/profiles/px_603668.dat};
\addplot[SoftRed, dashed, line width=0.7pt] table[x=x,y=y]{data/clean/profiles/pxzz_603668.dat};
\addplot[only marks, mark=*, mark size=1.1pt, SoftRed] table[x=x,y=y]{data/clean/profiles/pxzz_603668.dat};
\end{axis}
\end{tikzpicture}

\captionof{figure}[天马科技（603668）股价与周期骨架]{天马科技（603668）月度收盘价（元/股）与 ZigZag 周期骨架（阈值 25\%，圆点为周期转折点）}
\end{minipage}
\end{center}

\pfl{过去周期分析} 以 25\% 的 ZigZag 阈值划分，样本区间内股价形成 2 个主升段与 2 个主跌段。 最大主升段出现在 2021 年 1 月 至 2022 年 7 月，18 个月+209.0\%；最大主跌段为 2022 年 7 月 至 2025 年 1 月，30 个月-48.5\%。 最近一次转折出现在 2026 年 9 月（下跌段自 2025 年 12 月起，9 个月-26.9\%），当前处于该段的延续之中。 公司股价较申万农林牧渔指数提前约 64 个月见底，是板块中较早定价周期反转的公司之一。 水产饲料的需求取决于投苗量与存塘鱼价，第 \ref{sec:other-livestock} 章显示 2026 年淡水鱼价格与投苗量已率先回升，其需求驱动与生猪存栏的联系弱于畜禽饲料，公司 2026 年 6 月末资产负债率 72.1\%，高于全样本中位数 52.2\%，其股价因此更多反映资金面与信用风险，而不是价格弹性本身。 综合区间形态看，该股单段涨跌幅度偏大、以趋势段为主（最大主升 +209.0\%、最大主跌 -48.5\%），年化波动率 44.9\% 与全样本中位数 37.4\% 相比波动明显大于样本中位数。

\pfl{盈利情况} 2026 年上半年营业收入 31.32 亿元（同比 +5.3\%），归母净利润 0.12 亿元，同比 -79.5\%。 以年报为趋势对照，2023---2025 年营业收入由 69.98 亿元变为 60.01 亿元（两年年均复合增速 -7.4\%），归母净利润由 -1.88 亿元变为 -1.79 亿元。 按半年报口径，2026 年上半年的 ROE 为 0.6\%（未年化）、销售净利率 0.5\%、毛利率 8.3\%，ROE 在“水产饲料”4 家公司中按数值由高到低排第\zhnumber{3}位。 按同一口径，三级行业“水产饲料”4 家公司的半年 ROE 中位数为 1.22\%，该公司低于其中位数 0.64 个百分点。 从公司披露的原因看，业绩变动主要来自：2026 年，活鳗及烤鳗价格持续低位。 综合看，公司在报告期维持盈利（低于全样本半年 ROE 中位数 0.75\%），利润的可持续性取决于水产饲料景气的延续与成本管控的执行。

\begin{center}\small\setlength{\tabcolsep}{3.4pt}
\captionof{table}[天马科技（603668）关键财务指标]{天马科技（603668）关键财务指标（合并报表口径；两个半年报期均未年化；现金短债比 $=$ 货币资金 $\div$ 短期债务）}
\begin{adjustbox}{max width=\textwidth}
\begin{tabular}{@{}lrrrrr@{}}
\toprule
指标 & 2023 年 & 2024 年 & 2025 年 & 2025 年半年报 & 2026 年半年报 \\
\midrule
营业收入（亿元） & 69.98 & 58.54 & 60.01 & 29.74 & 31.32 \\
归母净利润（亿元） & -1.88 & 0.26 & -1.79 & 0.60 & 0.12 \\
毛利率（\%） & 7.5 & 10.9 & 7.3 & 10.7 & 8.3 \\
ROE（\%） & -9.5 & 1.2 & -8.3 & 2.7 & 0.6 \\
资产负债率（\%） & 70.8 & 68.6 & 72.6 & 69.6 & 72.1 \\
经营活动现金流净额（亿元） & 0.71 & 4.83 & 1.01 & 0.20 & 1.87 \\
货币资金（亿元） & 7.96 & 5.11 & 5.36 & 5.45 & 4.61 \\
有息负债合计（亿元） & 36.26 & 37.19 & 42.42 & 43.64 & 41.55 \\
现金短债比（倍） & 0.30 & 0.16 & 0.15 & 0.14 & 0.13 \\
\bottomrule
\end{tabular}
\end{adjustbox}
\end{center}

\begin{center}
\begin{minipage}{\textwidth}\centering
\begin{tikzpicture}
\begin{groupplot}[
  group style={group size=2 by 1, horizontal sep=0.60cm},
  width=6.85cm, height=4.60cm,
  tick label style={font=\tiny},
  title style={font=\scriptsize\heiti},
  yticklabel style={font=\tiny},
  ylabel style={font=\tiny},
  grid=major, grid style={LightGray, line width=0.3pt},
  axis line style={InkGray, line width=0.5pt},
  enlarge x limits={abs=0.8},
  legend style={font=\scriptsize, at={(0.5,-0.20)}, anchor=north,
                legend columns=2, draw=none, fill=none, column sep=6pt},
]
\nextgroupplot[
  ybar, bar width=4pt,
  ymin=-9.07, ymax=78.71,
  xtick={0,1,2,3,4,5.7},
  xticklabels={2021,2022,2023,2024,2025,2026H1},
  ylabel={亿元},
]
\addplot[fill=AgriGreen!75, draw=AgriGreen, bar shift=-2.6pt] table[x=i, y=rev]{data/clean/profiles/fin_603668.dat};
\addplot[fill=SoftRed!60, draw=SoftRed, bar shift=2.6pt] table[x=i, y=ni]{data/clean/profiles/fin_603668.dat};
\addplot[fill=AgriGreen!30, draw=AgriGreen, pattern=north east lines, bar shift=-2.6pt] table[x=x, y=rev]{data/clean/profiles/finh_603668.dat};
\addplot[fill=SoftRed!20, draw=SoftRed, pattern=north east lines, bar shift=2.6pt] table[x=x, y=ni]{data/clean/profiles/finh_603668.dat};
\legend{营业收入（年/半年）, 归母净利润（年/半年）}
\nextgroupplot[
  ymin=-16.1, ymax=80.8,
  xtick={0,1,2,3,4,5.7},
  xticklabels={2021,2022,2023,2024,2025,2026H1},
  ylabel={\%},
]
\addplot[AgriGold, mark=*, mark size=1.3pt] table[x=i, y=roe]{data/clean/profiles/fin_603668.dat};
\addplot[OilBlue, mark=square*, mark size=1.3pt] table[x=i, y=debt]{data/clean/profiles/fin_603668.dat};
\addplot[only marks, AgriGold, mark=*, mark size=2.0pt] table[x=x, y=roe]{data/clean/profiles/finh_603668.dat};
\addplot[only marks, OilBlue, mark=square*, mark size=2.0pt] table[x=x, y=debt]{data/clean/profiles/finh_603668.dat};
\legend{ROE, 资产负债率}
\end{groupplot}
\end{tikzpicture}
\begin{tikzpicture}
\begin{axis}[
  name=finbot,
  width=13.75cm, height=3.10cm, scale only axis,
  ymin=-4.24, ymax=47.51,
  xtick={0,1,2,3,4,5.7},
  xticklabels={2021,2022,2023,2024,2025,2026H1},
  tick label style={font=\tiny},
  yticklabel style={font=\tiny},
  ylabel={亿元}, ylabel style={font=\tiny},
  ybar, bar width=4pt,
  grid=major, grid style={LightGray, line width=0.3pt},
  axis line style={InkGray, line width=0.5pt},
  enlarge x limits={abs=0.8},
  legend style={font=\scriptsize, at={(0.5,-0.26)}, anchor=north,
                legend columns=2, draw=none, fill=none, column sep=10pt},
]
\addplot[fill=AgriGreen!55, draw=AgriGreen, bar shift=-3.0pt] table[x=i, y=cash]{data/clean/profiles/fin_603668.dat};
\addplot[fill=SoftRed!45, draw=SoftRed, bar shift=3.0pt] table[x=i, y=ibd]{data/clean/profiles/fin_603668.dat};
\addplot[fill=AgriGreen!25, draw=AgriGreen, pattern=north east lines, bar shift=-3.0pt] table[x=x, y=cash]{data/clean/profiles/finh_603668.dat};
\addplot[fill=SoftRed!20, draw=SoftRed, pattern=north east lines, bar shift=3.0pt] table[x=x, y=ibd]{data/clean/profiles/finh_603668.dat};
\legend{货币资金（左轴）, 有息负债（左轴）}
\end{axis}
\begin{axis}[
  at={(finbot.south east)}, anchor=south east,
  width=13.75cm, height=3.10cm, scale only axis,
  axis y line*=right, axis x line*=none, xtick=\empty,
  ymin=-1.26, ymax=6.28,
  tick label style={font=\tiny},
  yticklabel style={font=\tiny}, scaled y ticks=false,
  legend style={font=\scriptsize, at={(0.5,-0.44)}, anchor=north,
                legend columns=1, draw=none, fill=none},
]
\addplot[OilBlue, mark=*, mark size=2.2pt, line width=1.3pt] table[x=i, y=ocf]{data/clean/profiles/fin_603668.dat};
\addplot[only marks, OilBlue, mark=*, mark size=3.2pt] table[x=x, y=ocf]{data/clean/profiles/finh_603668.dat};
\legend{经营活动现金流净额（右轴，亿元，蓝点）}
\end{axis}
\end{tikzpicture}

\captionof{figure}[天马科技（603668）近年主要财务指标]{天马科技（603668）近年主要财务指标（左上：营业收入与归母净利润；右上：ROE 与资产负债率；下：货币资金、有息负债为左轴柱状，经营活动现金流净额为其右轴的蓝点折线。
2021---2025 年为年报口径，2026H1 为 2026 年半年报/半年末，与其前一年报之间留出间隔、以斜纹或空心点区分；半年报数据未年化）}
\end{minipage}
\end{center}

\pfl{财务分析} 2026 年 6 月末资产负债率 72.1\%，较 2025 年末下降 0.6 个百分点（样本中位数 52.2\%，所处三级行业内按由低到高排第\zhnumber{4}位）。 2026 年 6 月末货币资金 4.61 亿元、短期债务（短期借款与一年内到期非流动负债）34.56 亿元、长期债务 7.00 亿元，有息负债合计 41.55 亿元，现金短债比 0.13 倍——缺口明显。 净有息负债 36.95 亿元，相当于其 2026 年上半年营业收入的 118\%。 流动比率 0.88、速动比率 0.24，短期偿债指标偏紧。 2026 年上半年经营活动现金流净额 1.87 亿元，经营现金流与利润同向为正，内生资金可以支撑运营。 2026 年 6 月末在建工程 1.99 亿元，较上年末减少 0.54 亿元，在建投入在收缩。 2026 年 6 月末商誉 1.17 亿元，占归母净资产的 4.3\%，是净资产中最脆弱的一块。 综合看，账面货币资金对有息负债与短期债务的覆盖不足，滚动偿付对新增融资的依赖度较高；资产负债率高于样本中位数 19.9 个百分点；有息负债中短期债务占比更高，期限结构仍以滚动续作为主。 公司披露的股东与质押风险为：2026 年 8 月，实际控制人陈庆堂减持：2026-08-05 减持 207.99 万股（均价 11 元）、2026-08-10 减持 452.66 万股（均价 10 元）；2026 年 9 月，控股股东、实际控制人之一致行动人部分股份解除质押及质押。

\pfl{重大战略转型} 公司的战略动作可归为以下几类：其他动作（2 项）：2026 年，食品（烤鳗）销量创新高；2025 年 6 月，股票期权激励计划：第一个行权期行权条件于 2025 年 6 月经董事会审议通过成就；融资与资本运作（1 项）：2023 年，公司以简易程序定增募资建设“鳗鲡生态养殖基地建设项目（二期）”。 公告口径上，近三年（2023 年 9 月---2026 年 9 月）公司披露重大重组与并购类公告 0 条、对外投资与转型类公告 1 条、回购与股权激励类公告 12 条，合计公告 583 条。 第一大行业“饲料行业”的收入占比由 2021 年年报的 99.8\% 变为 87.4\%（12.4 个百分点），收入结构出现再平衡。 综合判断，报告期内公司有 1 条与战略相关的公告落地（重大重组与并购 0 条、对外投资与转型 1 条），资本运作与主业调整之间存在直接关联。

\begin{center}\small\setlength{\tabcolsep}{4pt}
\captionof{table}[天马科技（603668）历史融资与资本运作]{天马科技（603668）历史融资与资本运作（据公司公告与公开报道整理；金额为披露值，含首次公开发行、再融资、债券、银行借款与重整投资等）}
\begin{adjustbox}{max width=\textwidth}
\begin{tabular}{@{}p{1.45cm}p{2.55cm}p{4.05cm}p{4.95cm}@{}}
\toprule
时间 & 方式 & 规模 & 用途与说明 \\
\midrule
2017-01 & IPO（上交所主板） & 募集资金总额 3.29 亿元，\allowbreak{}净额 2.91 亿元 & 特种水产配合饲料生产项目（二期）等；发行价 6.21 元/\allowbreak{}股；发行前陈庆堂直接持股 38.43\%、\allowbreak{}通过天马投资控制 10.54\% \\
2018-04 & 公开发行可转换公司债券 & 5.00 亿元（305 万张，\allowbreak{}面值 100.00 元） & 特种水产配合饲料生产项目（二期）；；采用股份质押+保证担保 \\
2021-07 & 非公开发行股票（2020 年度方案） & 9,\allowbreak{}655.17 万股、\allowbreak{}5.80 元/\allowbreak{}股，\allowbreak{}募集资金总额 5.60 亿元、\allowbreak{}净额 5.52 亿元 & 2021-07-21 资金到账，\allowbreak{}7 月 30 日完成登记，\allowbreak{}总股本由 3.40 亿股增至 4.36 亿股 \\
2023-12 & 以简易程序向特定对象发行股票 & 2，\allowbreak{}053.39 万股、\allowbreak{}14.61 元/\allowbreak{}股 & 鳗鲡生态养殖基地建设项目（二期）及补充流动资金；；2023-12-29 资金到位 \\
2024-10-15 & 接受关联方财务资助 & 控股股东、\allowbreak{}实际控制人陈庆堂之全资控股公司天马投资提供最高不超过 3.50 亿元财务资助 & 公司及合并范围内子公司经营资金；额度有效期至 2025 年 12 月 31 日 \\
2025-12-09 & 接受关联方财务资助 & 天马投资提供最高不超过 4.50 亿元财务资助（无抵押、\allowbreak{}质押等担保措施） & 公司及子公司经营资金；利率不高于 LPR；截至 2026-09-16 公告日，\allowbreak{}天马投资向公司及子公司提供的借款本金余额为 1.88 亿元 \\
\bottomrule
\end{tabular}
\end{adjustbox}
\end{center}

\pfl{历史融投资情况} 从现金流量表看，2025 年公司取得借款收到的现金 35.80 亿元、偿还债务支付的现金 34.74 亿元、吸收权益性投资 0.46 亿元，筹资活动现金流净额 0.38 亿元（2023 年为取得借款 31.76 亿元、偿还债务 29.00 亿元）。 2026 年上半年筹资活动现金流净额为 -1.59 亿元（取得借款 13.55 亿元、偿还债务 12.64 亿元，未年化），已转为净流出，处于净还债状态。 2025 年分配股利、利润或偿付利息支付的现金合计 1.54 亿元。 公开披露的融资记录共 6 笔，工具类型以 IPO、可转债、定向增发为主；金额与用途见下表。 其中规模最大的两笔为：2025 年 12 月，接受关联方财务资助，天马投资提供最高不超过 4.50 亿元财务资助（无抵押、质押等担保措施）；2024 年 10 月，接受关联方财务资助，控股股东、实际控制人陈庆堂之全资控股公司天马投资提供最高不超过 3.50 亿元财务资助。 股本方面，近三年披露股本变动 7 次（限售股份上市 2 次、增发新股上市 1 次、转增 1 次）；总股本由 2021 年末的 4.36 亿股变为 5.02 亿股（+15.17\%）。 分红方面，近三年现金分红 4 次、合计每股派现 0.07 元，最近一次实施在 2024 年 6 月。 近三年“再融资”类公告 26 条，涉及定向增发、募集资金使用。 综合看，公司融资结构上股权与债务融资并举；2025 年筹资活动现金流净额 0.38 亿元，年度内仍为净融入，与前述经营与投资现金流的方向基本一致。

\pfl{未来潜在融资需求分析} 2026 年 6 月末货币资金 4.61 亿元，而短期债务 34.56 亿元（现金短债比 0.13 倍），2026 年上半年经营现金流净额 1.87 亿元。按“短期债务减去账面货币资金”的滚动偿付口径，静态缺口约 29.95 亿元；上半年盈利或接近盈亏平衡，该缺口不随经营亏损进一步扩大。 公司 2026 年 6 月末资产负债率 72.1\%、2026 年上半年 ROE 0.6\%（未年化），资产负债率偏高（72.1\%），融资需求首先用于置换与滚动存量债务、补充营运资金，属于“补血型”而非扩张型。 公开披露的后续资金安排线索包括：债务与授信安排：2025 年 12 月，截至 2026-09-16 借款本金余额 1.88 亿元。 资金需求还要叠加在建工程：期末在建工程 1.99 亿元、2026 年上半年购建固定资产等支出 1.16 亿元（未年化），这部分支出在周期底部通常会被主动放缓。 综合看，公司的外部资金需求以补充流动性与项目投入为主，滚动偿付口径下的静态缺口约 29.95 亿元，融资的时点与规模更多取决于信用条件与资本市场窗口。




\subsubsection{傲农生物（603363）}
\label{co:603363}

\pfl{公司基本情况} 公司前身可追溯至 2011 年，福建饲料起家、后跨界养猪的农牧企业；2024 年末完成破产重整，控股股东变更为泉州国资牵头联合体且无实际控制人，2026 年上半年猪价低迷亏损 1.90 亿元。 公司于 2017-09-26 在上交所首发上市，注册资本 27.31 亿元，总股本 8.70 亿股。 截至 2026 年 9 月 24 日收盘价 2.80 元，流通市值 60.1 亿元，近一年-43.9\%，流通市值在三级行业“畜禽饲料”的 11 家公司中按由大到小排第\zhnumber{4}位，近一年涨跌幅低于全样本中位数（-10.1\%）33.8 个百分点。 按 2026 年半年报披露的行业构成，收入占比前三位为饲料行业 69.5\%；饲养行业 16.9\%；屠宰食品 13.3\%。 与 2021 年年报相比，第一大行业口径由“饲料及动保行业”（56.4\%）变为“饲料行业”（69.5\%）；两期披露的分类口径有调整，占比差异中同时包含分类因素。

\pfl{历史股价} 本报告样本区间（2021 年 1 月 至 2026 年 9 月）内，股价由 13.40 元变为 2.80 元，累计 -79.1\%，区间最高 23.80 元（2022 年 3 月）、最低 2.68 元（2026 年 6 月）；当前价相当于区间最高价的 11.8\%，已接近区间最低价（1.04 倍）。 月度收益的年化波动率 52.5\%，高于全样本中位数 37.4\%，在三级行业“畜禽饲料”的 11 家公司中波动率由高到低排第\zhnumber{1}位。 把股价阶段与公司事件对齐看，2022 年 7 月 至 2025 年 2 月股价下跌-86.0\%，同期（2025 年 1 月）法院裁定终结重整程序后，公司被撤销因重整实施的退市风险警示。 把这些阶段与公司事件对齐后可以看到，几轮大涨大跌多由公司自身的资本运作与风险事件触发，行业景气只解释了其中一部分。

\begin{center}
\begin{minipage}{\textwidth}\centering
\begin{tikzpicture}
\begin{axis}[
  width=12.6cm, height=3.9cm, scale only axis,
  xmin=2020.922, xmax=2026.888, ymin=2.49, ymax=25.47,
  xtick={2021,2022,2023,2024,2025,2026},
  yearticks,
  yticklabel style={font=\scriptsize},
  ylabel={元/股}, ylabel style={font=\scriptsize},
  grid=major, grid style={LightGray, line width=0.3pt},
  axis line style={InkGray, line width=0.5pt},
  every axis plot/.append style={line width=0.9pt},
]
\addplot[AgriGreen] table[x=x,y=y]{data/clean/profiles/px_603363.dat};
\addplot[SoftRed, dashed, line width=0.7pt] table[x=x,y=y]{data/clean/profiles/pxzz_603363.dat};
\addplot[only marks, mark=*, mark size=1.1pt, SoftRed] table[x=x,y=y]{data/clean/profiles/pxzz_603363.dat};
\end{axis}
\end{tikzpicture}

\captionof{figure}[傲农生物（603363）股价与周期骨架]{傲农生物（603363）月度收盘价（元/股）与 ZigZag 周期骨架（阈值 25\%，圆点为周期转折点）}
\end{minipage}
\end{center}

\pfl{过去周期分析} 以 25\% 的 ZigZag 阈值划分，样本区间内股价形成 3 个主升段与 4 个主跌段。 最大主升段出现在 2021 年 7 月 至 2022 年 3 月，8 个月+196.8\%；最大主跌段为 2022 年 7 月 至 2025 年 2 月，31 个月-86.0\%。 最近一次转折出现在 2026 年 6 月（下跌段自 2025 年 10 月起，8 个月-47.6\%），当前处于该段的延续之中。 公司股价与申万农林牧渔指数几乎同步见底，个股并未走出独立行情。 饲料夹在原料成本与养殖需求之间（第 \ref{sec:feed} 节），量的部分取决于存栏、利的部分取决于原料与成品的价差，公司 2026 年 6 月末资产负债率 62.1\%、半年 ROE -6.2\%（未年化），在本轮下行中属于随行业同步调整的一类。 综合区间形态看，该股单段涨跌幅度偏大、以趋势段为主（最大主升 +196.8\%、最大主跌 -86.0\%），年化波动率 52.5\% 与全样本中位数 37.4\% 相比波动明显大于样本中位数。

\pfl{盈利情况} 2026 年上半年营业收入 45.21 亿元（同比 +14.3\%），归母净利润 -1.90 亿元，由上年同期的盈利转为亏损。 以年报为趋势对照，2023---2025 年营业收入由 194.58 亿元变为 84.75 亿元（两年年均复合增速 -34.0\%），归母净利润由 -36.51 亿元变为 0.78 亿元。 按半年报口径，2026 年上半年的 ROE 为 -6.2\%（未年化）、销售净利率 -6.0\%、毛利率 1.7\%，ROE 在“畜禽饲料”11 家公司中按数值由高到低排第\zhnumber{6}位。 按同一口径，三级行业“畜禽饲料”11 家公司的半年 ROE 中位数为 -6.25\%，该公司低于其中位数 0.00 个百分点。 从公司披露的原因看，业绩变动主要来自：2026 年上半年，生猪出栏量同比增加 40.53\%（78.33 万头→110.08 万头）。 综合看，公司仍处亏损区间、由盈转亏，半年 ROE 在三级行业内按由高到低排第\zhnumber{6}位，单位盈利能力的修复依赖行业景气与自身成本端的改善，报表弹性主要体现为价格与销量的方向性变化。

\begin{center}\small\setlength{\tabcolsep}{3.4pt}
\captionof{table}[傲农生物（603363）关键财务指标]{傲农生物（603363）关键财务指标（合并报表口径；两个半年报期均未年化；现金短债比 $=$ 货币资金 $\div$ 短期债务）}
\begin{adjustbox}{max width=\textwidth}
\begin{tabular}{@{}lrrrrr@{}}
\toprule
指标 & 2023 年 & 2024 年 & 2025 年 & 2025 年半年报 & 2026 年半年报 \\
\midrule
营业收入（亿元） & 194.58 & 87.63 & 84.75 & 39.57 & 45.21 \\
归母净利润（亿元） & -36.51 & 5.79 & 0.78 & 3.61 & -1.90 \\
毛利率（\%） & -4.7 & 4.5 & 6.3 & 8.3 & 1.7 \\
ROE（\%） & 0.0 & --- & 2.8 & 12.5 & -6.2 \\
资产负债率（\%） & 103.7 & 67.8 & 59.0 & 60.5 & 62.1 \\
经营活动现金流净额（亿元） & 9.84 & -1.17 & 1.53 & -3.16 & -3.29 \\
货币资金（亿元） & 2.08 & 13.84 & 6.98 & 8.20 & 4.75 \\
有息负债合计（亿元） & 80.52 & 29.08 & 23.29 & 26.64 & 23.80 \\
现金短债比（倍） & 0.04 & 1.09 & 0.80 & 0.93 & 0.49 \\
\bottomrule
\end{tabular}
\end{adjustbox}
\end{center}

\begin{center}
\begin{minipage}{\textwidth}\centering
\begin{tikzpicture}
\begin{groupplot}[
  group style={group size=2 by 1, horizontal sep=0.60cm},
  width=6.85cm, height=4.60cm,
  tick label style={font=\tiny},
  title style={font=\scriptsize\heiti},
  yticklabel style={font=\tiny},
  ylabel style={font=\tiny},
  grid=major, grid style={LightGray, line width=0.3pt},
  axis line style={InkGray, line width=0.5pt},
  enlarge x limits={abs=0.8},
  legend style={font=\scriptsize, at={(0.5,-0.20)}, anchor=north,
                legend columns=2, draw=none, fill=none, column sep=6pt},
]
\nextgroupplot[
  ybar, bar width=4pt,
  ymin=-61.77, ymax=246.45,
  xtick={0,1,2,3,4,5.7},
  xticklabels={2021,2022,2023,2024,2025,2026H1},
  ylabel={亿元},
]
\addplot[fill=AgriGreen!75, draw=AgriGreen, bar shift=-2.6pt] table[x=i, y=rev]{data/clean/profiles/fin_603363.dat};
\addplot[fill=SoftRed!60, draw=SoftRed, bar shift=2.6pt] table[x=i, y=ni]{data/clean/profiles/fin_603363.dat};
\addplot[fill=AgriGreen!30, draw=AgriGreen, pattern=north east lines, bar shift=-2.6pt] table[x=x, y=rev]{data/clean/profiles/finh_603363.dat};
\addplot[fill=SoftRed!20, draw=SoftRed, pattern=north east lines, bar shift=2.6pt] table[x=x, y=ni]{data/clean/profiles/finh_603363.dat};
\legend{营业收入（年/半年）, 归母净利润（年/半年）}
\nextgroupplot[
  ymin=-88.3, ymax=121.5,
  xtick={0,1,2,3,4,5.7},
  xticklabels={2021,2022,2023,2024,2025,2026H1},
  ylabel={\%},
]
\addplot[AgriGold, mark=*, mark size=1.3pt] table[x=i, y=roe]{data/clean/profiles/fin_603363.dat};
\addplot[OilBlue, mark=square*, mark size=1.3pt] table[x=i, y=debt]{data/clean/profiles/fin_603363.dat};
\addplot[only marks, AgriGold, mark=*, mark size=2.0pt] table[x=x, y=roe]{data/clean/profiles/finh_603363.dat};
\addplot[only marks, OilBlue, mark=square*, mark size=2.0pt] table[x=x, y=debt]{data/clean/profiles/finh_603363.dat};
\legend{ROE, 资产负债率}
\end{groupplot}
\end{tikzpicture}
\begin{tikzpicture}
\begin{axis}[
  name=finbot,
  width=13.75cm, height=3.10cm, scale only axis,
  ymin=-9.21, ymax=103.16,
  xtick={0,1,2,3,4,5.7},
  xticklabels={2021,2022,2023,2024,2025,2026H1},
  tick label style={font=\tiny},
  yticklabel style={font=\tiny},
  ylabel={亿元}, ylabel style={font=\tiny},
  ybar, bar width=4pt,
  grid=major, grid style={LightGray, line width=0.3pt},
  axis line style={InkGray, line width=0.5pt},
  enlarge x limits={abs=0.8},
  legend style={font=\scriptsize, at={(0.5,-0.26)}, anchor=north,
                legend columns=2, draw=none, fill=none, column sep=10pt},
]
\addplot[fill=AgriGreen!55, draw=AgriGreen, bar shift=-3.0pt] table[x=i, y=cash]{data/clean/profiles/fin_603363.dat};
\addplot[fill=SoftRed!45, draw=SoftRed, bar shift=3.0pt] table[x=i, y=ibd]{data/clean/profiles/fin_603363.dat};
\addplot[fill=AgriGreen!25, draw=AgriGreen, pattern=north east lines, bar shift=-3.0pt] table[x=x, y=cash]{data/clean/profiles/finh_603363.dat};
\addplot[fill=SoftRed!20, draw=SoftRed, pattern=north east lines, bar shift=3.0pt] table[x=x, y=ibd]{data/clean/profiles/finh_603363.dat};
\legend{货币资金（左轴）, 有息负债（左轴）}
\end{axis}
\begin{axis}[
  at={(finbot.south east)}, anchor=south east,
  width=13.75cm, height=3.10cm, scale only axis,
  axis y line*=right, axis x line*=none, xtick=\empty,
  ymin=-6.70, ymax=13.78,
  tick label style={font=\tiny},
  yticklabel style={font=\tiny}, scaled y ticks=false,
  legend style={font=\scriptsize, at={(0.5,-0.44)}, anchor=north,
                legend columns=1, draw=none, fill=none},
]
\addplot[OilBlue, mark=*, mark size=2.2pt, line width=1.3pt] table[x=i, y=ocf]{data/clean/profiles/fin_603363.dat};
\addplot[only marks, OilBlue, mark=*, mark size=3.2pt] table[x=x, y=ocf]{data/clean/profiles/finh_603363.dat};
\legend{经营活动现金流净额（右轴，亿元，蓝点）}
\end{axis}
\end{tikzpicture}

\captionof{figure}[傲农生物（603363）近年主要财务指标]{傲农生物（603363）近年主要财务指标（左上：营业收入与归母净利润；右上：ROE 与资产负债率；下：货币资金、有息负债为左轴柱状，经营活动现金流净额为其右轴的蓝点折线。
2021---2025 年为年报口径，2026H1 为 2026 年半年报/半年末，与其前一年报之间留出间隔、以斜纹或空心点区分；半年报数据未年化）}
\end{minipage}
\end{center}

\pfl{财务分析} 2026 年 6 月末资产负债率 62.1\%，较 2025 年末上升 3.0 个百分点（样本中位数 52.2\%，所处三级行业内按由低到高排第\zhnumber{6}位）。 2026 年 6 月末货币资金 4.75 亿元、短期债务（短期借款与一年内到期非流动负债）9.72 亿元、长期债务 14.08 亿元，有息负债合计 23.80 亿元，现金短债比 0.49 倍——覆盖偏紧。 净有息负债 19.05 亿元，相当于其 2026 年上半年营业收入的 42\%。 流动比率 0.89、速动比率 0.62，短期偿付压力较大。 2026 年上半年经营活动现金流净额 -3.29 亿元，经营现金流与利润同时为负，日常运营对债务与股东投入的依赖度较高。 2026 年 6 月末在建工程 4.45 亿元，较上年末增加 0.23 亿元，仍有在建投入。 2026 年 6 月末商誉 3.38 亿元，占归母净资产的 10.1\%，是净资产中最脆弱的一块。 综合看，账面货币资金对有息负债与短期债务的覆盖不足，滚动偿付对新增融资的依赖度较高；资产负债率高于样本中位数 9.9 个百分点；有息负债中长期债务占比更高，期限结构对短债滚动的压力较小。 公司披露的重整与债务违约风险为：2026 年上半年，由盈转亏：归母净利润 -1.90 亿元、扣非 -2.52 亿元（上年同期含重整收益基数）；2024 年，共发布 26 份新增部分债务逾期公告。

\pfl{重大战略转型} 公司的战略动作可归为以下几类：并购与重组（3 项）：2025---2026 年，重整后剥离部分闲置资产、保留规模猪场有序复产；2025 年 1 月，法院裁定终结重整程序后，公司被撤销因重整实施的退市风险警示；其他动作（1 项）：2026 年，推出 2026 年限制性股票激励计划。 公告口径上，近三年（2023 年 9 月---2026 年 9 月）公司披露重大重组与并购类公告 4 条、对外投资与转型类公告 5 条、回购与股权激励类公告 31 条，合计公告 749 条。 主营结构的口径变化已经落到报表上：2021 年年报的第一大行业为“饲料及动保行业”（56.4\%），2026 年半年报为“饲料行业”（69.5\%）；公司在这两期的分部划分方式有过调整。 综合判断，公司当前的战略重心不是扩张而是化解债务与完成重整，业务层面的调整让位于司法重整程序。

\begin{center}\small\setlength{\tabcolsep}{4pt}
\captionof{table}[傲农生物（603363）历史融资与资本运作]{傲农生物（603363）历史融资与资本运作（据公司公告与公开报道整理；金额为披露值，含首次公开发行、再融资、债券、银行借款与重整投资等）}
\begin{adjustbox}{max width=\textwidth}
\begin{tabular}{@{}p{1.45cm}p{2.55cm}p{4.05cm}p{4.95cm}@{}}
\toprule
时间 & 方式 & 规模 & 用途与说明 \\
\midrule
2017-09 & IPO（上交所主板） & 6,\allowbreak{}000.00 万股、\allowbreak{}4.79 元/\allowbreak{}股，\allowbreak{}募集资金总额 2.87 亿元 & 饲料及养殖项目等；上市日 2017-09-26 \\
2020—2021 & 公开发行可转换公司债券（公司口径为“2020 年公开发行可转换公司债券”） & 10.00 亿元（1，\allowbreak{}000 万张，\allowbreak{}面值 100.00 元，\allowbreak{}期限 6 年） & 母猪场、\allowbreak{}饲料及生猪项目，\allowbreak{}其中 3.00 亿元偿还银行借款；；方案于 2020 年 12 月处于申报推进阶段 \\
2020-05 & 非公开发行股票 & 募集资金总额 13.90 亿元、\allowbreak{}净额 13.65 亿元 & 9 个生猪养殖项目和 2 个饲料生产项目，\allowbreak{}其中 4.10 亿元用于偿还银行贷款；原拟募集不超过 13.90 亿元 \\
2022-04 & 非公开发行股票 & 1.30 亿股，\allowbreak{}募集资金 14.33 亿元 & 宜丰傲农 15，\allowbreak{}构成关联交易；原拟募集不超过 15.00 亿元 \\
2024-12 & 重整计划出资人权益调整（转增股票） & 转增股票中的 10.05 亿股由重整投资人有条件受让 & 清偿债务、\allowbreak{}引入战略及财务投资人；重整投资转增股份于 2026-01-29、\allowbreak{}2026-03-27、\allowbreak{}2026-05-27、\allowbreak{}2026-06-25 分批上市流通，\allowbreak{}合计 1.02 亿股 \\
2024-12-25 & 重整投资款 & 管理人账户已收到全体重整投资人或其指定主体支付的全部重整投资款合计 17.09 亿元（约 17.09 亿元） & 执行重整计划、\allowbreak{}清偿债务及补充营运资金；2024-09-13 确定泉发外贸联合体为中选重整投资人 \\
\bottomrule
\end{tabular}
\end{adjustbox}
\end{center}

\pfl{历史融投资情况} 从现金流量表看，2025 年公司取得借款收到的现金 2.97 亿元、偿还债务支付的现金 5.20 亿元、吸收权益性投资 0.32 亿元，筹资活动现金流净额 -5.67 亿元（2023 年为取得借款 55.35 亿元、偿还债务 57.39 亿元）。 2026 年上半年筹资活动现金流净额为 1.68 亿元（取得借款 1.19 亿元、偿还债务 1.18 亿元，未年化），仍处于净融入状态。 2025 年分配股利、利润或偿付利息支付的现金合计 1.07 亿元。 公开披露的融资记录共 6 笔，工具类型以 IPO、可转债、定向增发、重整投资为主；金额与用途见下表。 其中规模最大的两笔为：2024 年 12 月，重整投资，管理人账户已收到全体重整投资人或其指定主体支付的全部重整投资款合计 17.09 亿元（约 17.09 亿元）；2022 年 4 月，定向增发，1.30 亿股，募集资金 14.33 亿元。 股本方面，近三年披露股本变动 8 次（激励股份解禁 4 次、股份回购 1 次）；总股本由 2021 年末的 6.84 亿股变为 8.70 亿股（+27.16\%）。 分红方面，近三年未实施现金分红，与重整期间的资金约束一致。 近三年“再融资”类公告 23 条，涉及定向增发、募集资金使用。 综合看，公司融资结构上以债务滚动为主、股权融资为补充；2025 年筹资活动现金流净额 -5.67 亿元，融资节奏仍以维持存量债务周转为主，与前述经营与投资现金流的方向基本一致。

\pfl{未来潜在融资需求分析} 2026 年 6 月末货币资金 4.75 亿元，而短期债务 9.72 亿元（现金短债比 0.49 倍），2026 年上半年经营现金流净额 -3.29 亿元。按“短期债务减去账面货币资金”的滚动偿付口径，静态缺口约 4.97 亿元；若再计入上半年亏损的年化额（3.80 亿元），维持一年正常经营所需的外部资金约 4.97---8.77 亿元。 公司 2026 年上半年归母净利润为负（-1.90 亿元），但 6 月末资产负债率 62.1\% 尚未进入第一档（亏损且负债率 $>$65\%）的区间，当前融资需求以补充流动资金、支撑低谷期经营性支出为主。 公开披露的后续资金安排线索包括：股权融资线索：2024 年，公司终止 2022 年度向特定对象发行股票，转而通过破产重整化解债务；2026 年，第二次临时股东会审议因 2026 年限制性股票激励计划导致的公司股份总数与注册资本变更。 资金需求还要叠加在建工程：期末在建工程 4.45 亿元、2026 年上半年购建固定资产等支出 0.68 亿元（未年化），这部分支出在周期底部通常会被主动放缓。 综合看，公司的外部资金需求以营运资金周转与在建项目投入为主，滚动偿付口径下的静态缺口约 4.97 亿元，融资的时点与规模更多取决于债权人与重整投资人的进度。




\subsubsection{禾丰股份（603609）}
\label{co:603609}

\pfl{公司基本情况} 2014 年设立至今，禾丰股份是以饲料为核心、肉禽（白羽肉鸡屠宰）与生猪业务协同的农牧食品企业，荷兰 De Heus 持股约 9\%；2026 年上半年饲料外销创新高，受猪价拖累归母净利微亏。 公司于 2014-08-08 在上交所首发上市，注册资本 9.04 亿元，总股本 9.19 亿股。 截至 2026 年 9 月 24 日收盘价 6.22 元，流通市值 56.2 亿元，近一年-22.7\%，流通市值在三级行业“畜禽饲料”的 11 家公司中按由大到小排第\zhnumber{5}位，近一年涨跌幅低于全样本中位数（-10.1\%）12.6 个百分点。 按 2025 年年报披露的行业构成，收入占比前三位为肉禽 45.3\%；饲料 37.6\%；生猪 9.9\%。 与 2021 年年报相比，第一大行业口径由“饲料业务”（52.1\%）变为“肉禽”（45.3\%）；两期披露的分类口径有调整，占比差异中同时包含分类因素。

\pfl{历史股价} 以 2021 年 1 月为起点、2026 年 9 月为终点，股价由 11.70 元变为 6.22 元，累计 -46.8\%，区间最高 13.85 元（2022 年 11 月）、最低 5.76 元（2026 年 6 月）；当前价相当于区间最高价的 44.9\%，为区间最低价的 1.08 倍。 月度收益的年化波动率 28.8\%，低于全样本中位数 37.4\%，在三级行业“畜禽饲料”的 11 家公司中波动率由高到低排第\zhnumber{11}位。

\begin{center}
\begin{minipage}{\textwidth}\centering
\begin{tikzpicture}
\begin{axis}[
  width=12.6cm, height=3.9cm, scale only axis,
  xmin=2020.922, xmax=2026.888, ymin=5.36, ymax=14.82,
  xtick={2021,2022,2023,2024,2025,2026},
  yearticks,
  yticklabel style={font=\scriptsize},
  ylabel={元/股}, ylabel style={font=\scriptsize},
  grid=major, grid style={LightGray, line width=0.3pt},
  axis line style={InkGray, line width=0.5pt},
  every axis plot/.append style={line width=0.9pt},
]
\addplot[AgriGreen] table[x=x,y=y]{data/clean/profiles/px_603609.dat};
\addplot[SoftRed, dashed, line width=0.7pt] table[x=x,y=y]{data/clean/profiles/pxzz_603609.dat};
\addplot[only marks, mark=*, mark size=1.1pt, SoftRed] table[x=x,y=y]{data/clean/profiles/pxzz_603609.dat};
\end{axis}
\end{tikzpicture}

\captionof{figure}[禾丰股份（603609）股价与周期骨架]{禾丰股份（603609）月度收盘价（元/股）与 ZigZag 周期骨架（阈值 25\%，圆点为周期转折点）}
\end{minipage}
\end{center}

\pfl{过去周期分析} 以 25\% 的 ZigZag 阈值划分，样本区间内股价形成 2 个主升段与 3 个主跌段。 最大主跌段为 2022 年 11 月 至 2024 年 6 月，19 个月-52.8\%；最大主升段出现在 2022 年 4 月 至 2022 年 11 月，7 个月+74.2\%。 最近一次转折出现在 2026 年 6 月（下跌段自 2025 年 5 月起，13 个月-34.9\%），当前处于该段的延续之中。 公司股价与申万农林牧渔指数几乎同步见底，个股并未走出独立行情。 饲料夹在原料成本与养殖需求之间（第 \ref{sec:feed} 节），量的部分取决于存栏、利的部分取决于原料与成品的价差，公司 2026 年 6 月末资产负债率 57.7\%、半年 ROE -0.1\%（未年化），在本轮下行中属于随行业同步调整的一类。 综合区间形态看，该股单段涨跌幅度偏大、以趋势段为主（最大主升 +74.2\%、最大主跌 -52.8\%），年化波动率 28.8\% 与全样本中位数 37.4\% 相比波动小于样本中位数。

\pfl{盈利情况} 2026 年上半年营业收入 189.63 亿元（同比 +8.9\%），归母净利润 -0.08 亿元，由上年同期的盈利转为亏损。 以年报为趋势对照，2023---2025 年营业收入由 359.70 亿元变为 357.62 亿元（两年年均复合增速 -0.3\%），归母净利润由 -4.57 亿元变为 0.53 亿元。 按半年报口径，2026 年上半年的 ROE 为 -0.1\%（未年化）、销售净利率 0.1\%、毛利率 6.3\%，ROE 在“畜禽饲料”11 家公司中按数值由高到低排第\zhnumber{4}位。 同期全样本半年 ROE 中位数 0.75\%，公司与之相差 0.88 个百分点（弱于中位数公司）。 综合看，公司仍处亏损区间、由盈转亏，半年 ROE 在三级行业内按由高到低排第\zhnumber{4}位，单位盈利能力的修复依赖行业景气与自身成本端的改善，报表弹性主要体现为价格与销量的方向性变化。

\begin{center}\small\setlength{\tabcolsep}{3.4pt}
\captionof{table}[禾丰股份（603609）关键财务指标]{禾丰股份（603609）关键财务指标（合并报表口径；两个半年报期均未年化；现金短债比 $=$ 货币资金 $\div$ 短期债务）}
\begin{adjustbox}{max width=\textwidth}
\begin{tabular}{@{}lrrrrr@{}}
\toprule
指标 & 2023 年 & 2024 年 & 2025 年 & 2025 年半年报 & 2026 年半年报 \\
\midrule
营业收入（亿元） & 359.70 & 325.45 & 357.62 & 174.07 & 189.63 \\
归母净利润（亿元） & -4.57 & 3.42 & 0.53 & 2.33 & -0.08 \\
毛利率（\%） & 4.0 & 6.4 & 5.8 & 6.4 & 6.3 \\
ROE（\%） & -6.6 & 5.1 & 0.8 & 3.4 & -0.1 \\
资产负债率（\%） & 49.0 & 49.1 & 56.3 & 56.4 & 57.7 \\
经营活动现金流净额（亿元） & 9.56 & 11.77 & 6.81 & -5.12 & -3.51 \\
货币资金（亿元） & 19.15 & 17.94 & 17.20 & 13.98 & 12.18 \\
有息负债合计（亿元） & 44.98 & 42.79 & 63.44 & 65.26 & 66.24 \\
现金短债比（倍） & 1.01 & 1.69 & 0.74 & 0.56 & 0.53 \\
\bottomrule
\end{tabular}
\end{adjustbox}
\end{center}

\begin{center}
\begin{minipage}{\textwidth}\centering
\begin{tikzpicture}
\begin{groupplot}[
  group style={group size=2 by 1, horizontal sep=0.60cm},
  width=6.85cm, height=4.60cm,
  tick label style={font=\tiny},
  title style={font=\scriptsize\heiti},
  yticklabel style={font=\tiny},
  ylabel style={font=\tiny},
  grid=major, grid style={LightGray, line width=0.3pt},
  axis line style={InkGray, line width=0.5pt},
  enlarge x limits={abs=0.8},
  legend style={font=\scriptsize, at={(0.5,-0.20)}, anchor=north,
                legend columns=2, draw=none, fill=none, column sep=6pt},
]
\nextgroupplot[
  ybar, bar width=4pt,
  ymin=-41.00, ymax=403.42,
  xtick={0,1,2,3,4,5.7},
  xticklabels={2021,2022,2023,2024,2025,2026H1},
  ylabel={亿元},
]
\addplot[fill=AgriGreen!75, draw=AgriGreen, bar shift=-2.6pt] table[x=i, y=rev]{data/clean/profiles/fin_603609.dat};
\addplot[fill=SoftRed!60, draw=SoftRed, bar shift=2.6pt] table[x=i, y=ni]{data/clean/profiles/fin_603609.dat};
\addplot[fill=AgriGreen!30, draw=AgriGreen, pattern=north east lines, bar shift=-2.6pt] table[x=x, y=rev]{data/clean/profiles/finh_603609.dat};
\addplot[fill=SoftRed!20, draw=SoftRed, pattern=north east lines, bar shift=2.6pt] table[x=x, y=ni]{data/clean/profiles/finh_603609.dat};
\legend{营业收入（年/半年）, 归母净利润（年/半年）}
\nextgroupplot[
  ymin=-11.7, ymax=64.1,
  xtick={0,1,2,3,4,5.7},
  xticklabels={2021,2022,2023,2024,2025,2026H1},
  ylabel={\%},
]
\addplot[AgriGold, mark=*, mark size=1.3pt] table[x=i, y=roe]{data/clean/profiles/fin_603609.dat};
\addplot[OilBlue, mark=square*, mark size=1.3pt] table[x=i, y=debt]{data/clean/profiles/fin_603609.dat};
\addplot[only marks, AgriGold, mark=*, mark size=2.0pt] table[x=x, y=roe]{data/clean/profiles/finh_603609.dat};
\addplot[only marks, OilBlue, mark=square*, mark size=2.0pt] table[x=x, y=debt]{data/clean/profiles/finh_603609.dat};
\legend{ROE, 资产负债率}
\end{groupplot}
\end{tikzpicture}
\begin{tikzpicture}
\begin{axis}[
  name=finbot,
  width=13.75cm, height=3.10cm, scale only axis,
  ymin=-6.62, ymax=74.18,
  xtick={0,1,2,3,4,5.7},
  xticklabels={2021,2022,2023,2024,2025,2026H1},
  tick label style={font=\tiny},
  yticklabel style={font=\tiny},
  ylabel={亿元}, ylabel style={font=\tiny},
  ybar, bar width=4pt,
  grid=major, grid style={LightGray, line width=0.3pt},
  axis line style={InkGray, line width=0.5pt},
  enlarge x limits={abs=0.8},
  legend style={font=\scriptsize, at={(0.5,-0.26)}, anchor=north,
                legend columns=2, draw=none, fill=none, column sep=10pt},
]
\addplot[fill=AgriGreen!55, draw=AgriGreen, bar shift=-3.0pt] table[x=i, y=cash]{data/clean/profiles/fin_603609.dat};
\addplot[fill=SoftRed!45, draw=SoftRed, bar shift=3.0pt] table[x=i, y=ibd]{data/clean/profiles/fin_603609.dat};
\addplot[fill=AgriGreen!25, draw=AgriGreen, pattern=north east lines, bar shift=-3.0pt] table[x=x, y=cash]{data/clean/profiles/finh_603609.dat};
\addplot[fill=SoftRed!20, draw=SoftRed, pattern=north east lines, bar shift=3.0pt] table[x=x, y=ibd]{data/clean/profiles/finh_603609.dat};
\legend{货币资金（左轴）, 有息负债（左轴）}
\end{axis}
\begin{axis}[
  at={(finbot.south east)}, anchor=south east,
  width=13.75cm, height=3.10cm, scale only axis,
  axis y line*=right, axis x line*=none, xtick=\empty,
  ymin=-7.49, ymax=16.35,
  tick label style={font=\tiny},
  yticklabel style={font=\tiny}, scaled y ticks=false,
  legend style={font=\scriptsize, at={(0.5,-0.44)}, anchor=north,
                legend columns=1, draw=none, fill=none},
]
\addplot[OilBlue, mark=*, mark size=2.2pt, line width=1.3pt] table[x=i, y=ocf]{data/clean/profiles/fin_603609.dat};
\addplot[only marks, OilBlue, mark=*, mark size=3.2pt] table[x=x, y=ocf]{data/clean/profiles/finh_603609.dat};
\legend{经营活动现金流净额（右轴，亿元，蓝点）}
\end{axis}
\end{tikzpicture}

\captionof{figure}[禾丰股份（603609）近年主要财务指标]{禾丰股份（603609）近年主要财务指标（左上：营业收入与归母净利润；右上：ROE 与资产负债率；下：货币资金、有息负债为左轴柱状，经营活动现金流净额为其右轴的蓝点折线。
2021---2025 年为年报口径，2026H1 为 2026 年半年报/半年末，与其前一年报之间留出间隔、以斜纹或空心点区分；半年报数据未年化）}
\end{minipage}
\end{center}

\pfl{财务分析} 2026 年 6 月末资产负债率 57.7\%，较 2025 年末上升 1.4 个百分点（样本中位数 52.2\%，所处三级行业内按由低到高排第\zhnumber{5}位）。 2026 年 6 月末货币资金 12.18 亿元、短期债务（短期借款与一年内到期非流动负债）22.89 亿元、长期债务 43.35 亿元，有息负债合计 66.24 亿元，现金短债比 0.53 倍——覆盖偏紧。 净有息负债 54.06 亿元，相当于其 2026 年上半年营业收入的 29\%。 流动比率 1.22、速动比率 0.58，短期偿债能力处于正常区间。 2026 年上半年经营活动现金流净额 -3.51 亿元，经营现金流与利润同时为负，现金消耗较快。 2026 年 6 月末在建工程 1.07 亿元，较上年末减少 1.77 亿元，在建投入在收缩。 综合看，现金短债比 0.53 倍，短期资金安排偏紧；资产负债率高于样本中位数 5.5 个百分点；有息负债中长期债务占比更高，期限结构对短债滚动的压力较小。

\pfl{重大战略转型} 从公开披露看，公司的战略推进集中在：其他动作（2 项）：2026 年，经营活动现金流同比改善 31.40\%；2025 年 4 月，公司推动 2026 年上半年肉禽业务收入同比增长 32.5\%；融资与资本运作（1 项）：2026 年 4 月，可转债募集资金结项及终止项目形成结余募集资金本金 7.30 亿元、利息收益 313.79 万元。 公告口径上，近三年（2023 年 9 月---2026 年 9 月）公司披露重大重组与并购类公告 5 条、对外投资与转型类公告 0 条、回购与股权激励类公告 52 条，合计公告 463 条。 主营结构的口径变化已经落到报表上：2021 年年报的第一大行业为“饲料业务”（52.1\%），2025 年年报为“肉禽”（45.3\%）；公司在这两期的分部划分方式有过调整。 综合判断，报告期内公司有 5 条与战略相关的公告落地（重大重组与并购 5 条、对外投资与转型 0 条），资本运作与主业调整之间存在直接关联。

\begin{center}\small\setlength{\tabcolsep}{4pt}
\captionof{table}[禾丰股份（603609）历史融资与资本运作]{禾丰股份（603609）历史融资与资本运作（据公司公告与公开报道整理；金额为披露值，含首次公开发行、再融资、债券、银行借款与重整投资等）}
\begin{adjustbox}{max width=\textwidth}
\begin{tabular}{@{}p{1.45cm}p{2.55cm}p{4.05cm}p{4.95cm}@{}}
\toprule
时间 & 方式 & 规模 & 用途与说明 \\
\midrule
2014-08 & IPO（上交所） & 募集资金总额 4.70 亿元 & 上市日 2014-08-08；2014 年年度报告披露“首次发行”募集资金总额 4.70 亿元 \\
2022-04 & 公开发行可转换公司债券 & 15.00 亿元（1，\allowbreak{}500 万张，\allowbreak{}100.00 元/\allowbreak{}张），\allowbreak{}扣除发行费用 1 & 募投项目；；2022-04-28 资金到位，\allowbreak{}初始转股价 10.22 元/\allowbreak{}股 \\
2024 & 股份回购（视同现金分红） & 两轮回购合计 2.87 亿元、\allowbreak{}回购 3 & 维护股价与优化股本结构、\allowbreak{}股东回报；2024 年度现金分红与回购金额合计约 3.38 亿元；截至 2025-04-20 总股本 9.19 亿股 \\
2025 & 股份回购（视同现金分红） & 以现金为对价集中竞价回购金额 5，\allowbreak{}646.65 万元 & 股东回报；2025 年度不派发现金红利、\allowbreak{}不送红股、\allowbreak{}不以公积金转增股本 \\
2026-01 & 回购股份注销（不调整可转债转股价） & --- & 注销回购股份；2026-01-13 公告回购股份注销不调整“禾丰转债”转股价格 \\
2026-09 & 综合授信额度调整 & 2026 年度拟向金融机构申请综合授信总额由“不超过 88.00 亿元”调整为“不超过 89.70 亿元” & 流动资金贷款、\allowbreak{}项目贷款、\allowbreak{}银行承兑汇票、\allowbreak{}保函、\allowbreak{}信用证、\allowbreak{}融资租赁等；2026 年第一次临时股东会会议材料（2026-09-01 披露） \\
\bottomrule
\end{tabular}
\end{adjustbox}
\end{center}

\pfl{历史融投资情况} 从现金流量表看，2025 年公司取得借款收到的现金 32.51 亿元、偿还债务支付的现金 25.24 亿元、吸收权益性投资 0.97 亿元，筹资活动现金流净额 3.00 亿元（2023 年为取得借款 19.59 亿元、偿还债务 19.13 亿元）。 2026 年上半年筹资活动现金流净额为 0.48 亿元（取得借款 19.36 亿元、偿还债务 16.96 亿元，未年化），仍处于净融入状态。 2025 年分配股利、利润或偿付利息支付的现金合计 2.05 亿元。 公开披露的融资记录共 6 笔，工具类型以 IPO、可转债、银行授信为主；金额与用途见下表。 其中规模最大的两笔为：2026 年 9 月，银行授信，2026 年度拟向金融机构申请综合授信总额由“不超过 88.00 亿元”调整为“不超过 89.70 亿元”；2022 年 4 月，可转债，15.00 亿元（1，500 万张，100.00 元/张），扣除发行费用 1。 股本方面，近三年披露股本变动 13 次（可转债转股 9 次、其他 2 次）；总股本由 2021 年末的 9.22 亿股变为 9.19 亿股（-0.27\%）。 分红方面，近三年现金分红 2 次、合计每股派现 0.18 元，最近一次实施在 2024 年 12 月。 近三年“再融资”类公告 26 条，涉及可转债、募集资金使用。 综合看，公司融资结构上以债务滚动为主、股权融资为补充；2025 年筹资活动现金流净额 3.00 亿元，年度内仍为净融入，与前述经营与投资现金流的方向基本一致。

\pfl{未来潜在融资需求分析} 2026 年 6 月末货币资金 12.18 亿元，而短期债务 22.89 亿元（现金短债比 0.53 倍），2026 年上半年经营现金流净额 -3.51 亿元。按“短期债务减去账面货币资金”的滚动偿付口径，静态缺口约 10.71 亿元；若再计入上半年亏损的年化额（0.17 亿元），维持一年正常经营所需的外部资金约 10.71---10.88 亿元。 公司 2026 年上半年归母净利润为负（-0.08 亿元），但 6 月末资产负债率 57.7\% 尚未进入第一档（亏损且负债率 $>$65\%）的区间，当前融资需求以补充流动资金、支撑低谷期经营性支出为主。 公开披露的后续资金安排线索包括：债务与授信安排：2026 年度，综合授信总额由不超过 88.00 亿元上调至不超过 89.70 亿元；2026 年 8 月，为子公司提供担保情况公告，公司以担保方式支持子公司融资。 资金需求还要叠加在建工程：期末在建工程 1.07 亿元、2026 年上半年购建固定资产等支出 1.80 亿元（未年化），这部分支出在周期底部通常会被主动放缓。 综合看，公司的外部资金需求以补充流动性与项目投入为主，滚动偿付口径下的静态缺口约 10.71 亿元，能否落地取决于授信续作与发行窗口的配合。




\subsubsection{唐人神（002567）}
\label{co:002567}

\pfl{公司基本情况} 公司前身可追溯至 2011 年，湖南株洲的饲料—生猪养殖—肉品一体化企业；2025 年亏损 11.87 亿元，2026 年上半年亏损扩大至 10.46 亿元，以低成本与降负债为穿越周期主线。 公司于 2011-03-25 在深交所主板首发上市，注册资本 14.33 亿元，总股本 14.33 亿股。 截至 2026 年 9 月 24 日收盘价 3.58 元，流通市值 51.3 亿元，近一年-26.2\%，流通市值在三级行业“畜禽饲料”的 11 家公司中按由大到小排第\zhnumber{6}位，近一年涨跌幅低于全样本中位数（-10.1\%）16.1 个百分点。 按 2026 年半年报披露的行业构成，收入占比前三位为饲料 69.0\%；牲猪养殖等 24.9\%；屠宰及肉类 6.0\%。

\pfl{历史股价} 自 2021 年 1 月至 2026 年 9 月，股价由 7.93 元变为 3.58 元，累计 -54.9\%，区间最高 10.50 元（2022 年 3 月）、最低 3.35 元（2026 年 6 月）；当前价相当于区间最高价的 34.1\%，为区间最低价的 1.07 倍。 月度收益的年化波动率 33.5\%，低于全样本中位数 37.4\%，在三级行业“畜禽饲料”的 11 家公司中波动率由高到低排第\zhnumber{8}位。

\begin{center}
\begin{minipage}{\textwidth}\centering
\begin{tikzpicture}
\begin{axis}[
  width=12.6cm, height=3.9cm, scale only axis,
  xmin=2020.922, xmax=2026.888, ymin=3.12, ymax=11.24,
  xtick={2021,2022,2023,2024,2025,2026},
  yearticks,
  yticklabel style={font=\scriptsize},
  ylabel={元/股}, ylabel style={font=\scriptsize},
  grid=major, grid style={LightGray, line width=0.3pt},
  axis line style={InkGray, line width=0.5pt},
  every axis plot/.append style={line width=0.9pt},
]
\addplot[AgriGreen] table[x=x,y=y]{data/clean/profiles/px_002567.dat};
\addplot[SoftRed, dashed, line width=0.7pt] table[x=x,y=y]{data/clean/profiles/pxzz_002567.dat};
\addplot[only marks, mark=*, mark size=1.1pt, SoftRed] table[x=x,y=y]{data/clean/profiles/pxzz_002567.dat};
\end{axis}
\end{tikzpicture}

\captionof{figure}[唐人神（002567）股价与周期骨架]{唐人神（002567）月度收盘价（元/股）与 ZigZag 周期骨架（阈值 25\%，圆点为周期转折点）}
\end{minipage}
\end{center}

\pfl{过去周期分析} 以 25\% 的 ZigZag 阈值划分，样本区间内股价形成 2 个主升段与 3 个主跌段。 最大主跌段为 2022 年 7 月 至 2026 年 6 月，47 个月-66.7\%；最大主升段出现在 2021 年 8 月 至 2022 年 3 月，7 个月+79.5\%。 最近一次转折出现在 2026 年 6 月（下跌段自 2022 年 7 月起，47 个月-66.7\%），当前处于该段的延续之中。 公司股价与申万农林牧渔指数几乎同步见底，个股并未走出独立行情。 饲料夹在原料成本与养殖需求之间（第 \ref{sec:feed} 节），量的部分取决于存栏、利的部分取决于原料与成品的价差，公司 2026 年 6 月末资产负债率 71.1\%，高于全样本中位数 52.2\%，其股价因此更多反映资金面与信用风险，而不是价格弹性本身。 综合区间形态看，该股单段涨跌幅度偏大、以趋势段为主（最大主升 +79.5\%、最大主跌 -66.7\%），年化波动率 33.5\% 与全样本中位数 37.4\% 相比波动小于样本中位数。

\pfl{盈利情况} 2026 年上半年营业收入 108.83 亿元（同比 -12.7\%），归母净利润 -10.46 亿元，亏损同比扩大 9.86 亿元。 以年报为趋势对照，2023---2025 年营业收入由 269.49 亿元变为 242.15 亿元（两年年均复合增速 -5.2\%），归母净利润由 -15.26 亿元变为 -11.87 亿元。 按半年报口径，2026 年上半年的 ROE 为 -24.9\%（未年化）、销售净利率 -9.8\%、毛利率 2.1\%，ROE 在“畜禽饲料”11 家公司中按数值由高到低排第\zhnumber{10}位。 按同一口径，三级行业“畜禽饲料”11 家公司的半年 ROE 中位数为 -6.25\%，该公司低于其中位数 18.60 个百分点。 从公司披露的原因看，业绩变动主要来自：2025 年，公司对存栏生猪计提存货跌价准备 3.76 亿元；2025---2026 年，坚持“养猪以低成本为中心”，通过项目制推进成本改善。 公司给出的应对方向是：具有年出栏生猪 180 万头产能。 综合看，公司仍处亏损区间、亏损同比扩大，半年 ROE 在三级行业内按由高到低排第\zhnumber{10}位，单位盈利能力的修复依赖行业景气与自身成本端的改善，报表弹性主要体现为价格与销量的方向性变化。

\begin{center}\small\setlength{\tabcolsep}{3.4pt}
\captionof{table}[唐人神（002567）关键财务指标]{唐人神（002567）关键财务指标（合并报表口径；两个半年报期均未年化；现金短债比 $=$ 货币资金 $\div$ 短期债务）}
\begin{adjustbox}{max width=\textwidth}
\begin{tabular}{@{}lrrrrr@{}}
\toprule
指标 & 2023 年 & 2024 年 & 2025 年 & 2025 年半年报 & 2026 年半年报 \\
\midrule
营业收入（亿元） & 269.49 & 243.43 & 242.15 & 124.68 & 108.83 \\
归母净利润（亿元） & -15.26 & 3.55 & -11.87 & -0.60 & -10.46 \\
毛利率（\%） & 2.7 & 8.8 & 4.1 & 6.6 & 2.1 \\
ROE（\%） & -25.0 & 6.3 & -22.4 & -1.0 & -24.9 \\
资产负债率（\%） & 64.8 & 63.4 & 66.4 & 64.5 & 71.1 \\
经营活动现金流净额（亿元） & 5.36 & 9.36 & 9.15 & 6.60 & -0.66 \\
货币资金（亿元） & 20.91 & 23.05 & 19.47 & 24.86 & 21.10 \\
有息负债合计（亿元） & 76.23 & 67.63 & 58.64 & 67.03 & 61.18 \\
现金短债比（倍） & 0.83 & 1.01 & 0.76 & 0.97 & 0.85 \\
\bottomrule
\end{tabular}
\end{adjustbox}
\end{center}

\begin{center}
\begin{minipage}{\textwidth}\centering
\begin{tikzpicture}
\begin{groupplot}[
  group style={group size=2 by 1, horizontal sep=0.60cm},
  width=6.85cm, height=4.60cm,
  tick label style={font=\tiny},
  title style={font=\scriptsize\heiti},
  yticklabel style={font=\tiny},
  ylabel style={font=\tiny},
  grid=major, grid style={LightGray, line width=0.3pt},
  axis line style={InkGray, line width=0.5pt},
  enlarge x limits={abs=0.8},
  legend style={font=\scriptsize, at={(0.5,-0.20)}, anchor=north,
                legend columns=2, draw=none, fill=none, column sep=6pt},
]
\nextgroupplot[
  ybar, bar width=4pt,
  ymin=-43.73, ymax=303.66,
  xtick={0,1,2,3,4,5.7},
  xticklabels={2021,2022,2023,2024,2025,2026H1},
  ylabel={亿元},
]
\addplot[fill=AgriGreen!75, draw=AgriGreen, bar shift=-2.6pt] table[x=i, y=rev]{data/clean/profiles/fin_002567.dat};
\addplot[fill=SoftRed!60, draw=SoftRed, bar shift=2.6pt] table[x=i, y=ni]{data/clean/profiles/fin_002567.dat};
\addplot[fill=AgriGreen!30, draw=AgriGreen, pattern=north east lines, bar shift=-2.6pt] table[x=x, y=rev]{data/clean/profiles/finh_002567.dat};
\addplot[fill=SoftRed!20, draw=SoftRed, pattern=north east lines, bar shift=2.6pt] table[x=x, y=ni]{data/clean/profiles/finh_002567.dat};
\legend{营业收入（年/半年）, 归母净利润（年/半年）}
\nextgroupplot[
  ymin=-32.7, ymax=80.7,
  xtick={0,1,2,3,4,5.7},
  xticklabels={2021,2022,2023,2024,2025,2026H1},
  ylabel={\%},
]
\addplot[AgriGold, mark=*, mark size=1.3pt] table[x=i, y=roe]{data/clean/profiles/fin_002567.dat};
\addplot[OilBlue, mark=square*, mark size=1.3pt] table[x=i, y=debt]{data/clean/profiles/fin_002567.dat};
\addplot[only marks, AgriGold, mark=*, mark size=2.0pt] table[x=x, y=roe]{data/clean/profiles/finh_002567.dat};
\addplot[only marks, OilBlue, mark=square*, mark size=2.0pt] table[x=x, y=debt]{data/clean/profiles/finh_002567.dat};
\legend{ROE, 资产负债率}
\end{groupplot}
\end{tikzpicture}
\begin{tikzpicture}
\begin{axis}[
  name=finbot,
  width=13.75cm, height=3.10cm, scale only axis,
  ymin=-7.62, ymax=85.37,
  xtick={0,1,2,3,4,5.7},
  xticklabels={2021,2022,2023,2024,2025,2026H1},
  tick label style={font=\tiny},
  yticklabel style={font=\tiny},
  ylabel={亿元}, ylabel style={font=\tiny},
  ybar, bar width=4pt,
  grid=major, grid style={LightGray, line width=0.3pt},
  axis line style={InkGray, line width=0.5pt},
  enlarge x limits={abs=0.8},
  legend style={font=\scriptsize, at={(0.5,-0.26)}, anchor=north,
                legend columns=2, draw=none, fill=none, column sep=10pt},
]
\addplot[fill=AgriGreen!55, draw=AgriGreen, bar shift=-3.0pt] table[x=i, y=cash]{data/clean/profiles/fin_002567.dat};
\addplot[fill=SoftRed!45, draw=SoftRed, bar shift=3.0pt] table[x=i, y=ibd]{data/clean/profiles/fin_002567.dat};
\addplot[fill=AgriGreen!25, draw=AgriGreen, pattern=north east lines, bar shift=-3.0pt] table[x=x, y=cash]{data/clean/profiles/finh_002567.dat};
\addplot[fill=SoftRed!20, draw=SoftRed, pattern=north east lines, bar shift=3.0pt] table[x=x, y=ibd]{data/clean/profiles/finh_002567.dat};
\legend{货币资金（左轴）, 有息负债（左轴）}
\end{axis}
\begin{axis}[
  at={(finbot.south east)}, anchor=south east,
  width=13.75cm, height=3.10cm, scale only axis,
  axis y line*=right, axis x line*=none, xtick=\empty,
  ymin=-7.21, ymax=13.31,
  tick label style={font=\tiny},
  yticklabel style={font=\tiny}, scaled y ticks=false,
  legend style={font=\scriptsize, at={(0.5,-0.44)}, anchor=north,
                legend columns=1, draw=none, fill=none},
]
\addplot[OilBlue, mark=*, mark size=2.2pt, line width=1.3pt] table[x=i, y=ocf]{data/clean/profiles/fin_002567.dat};
\addplot[only marks, OilBlue, mark=*, mark size=3.2pt] table[x=x, y=ocf]{data/clean/profiles/finh_002567.dat};
\legend{经营活动现金流净额（右轴，亿元，蓝点）}
\end{axis}
\end{tikzpicture}

\captionof{figure}[唐人神（002567）近年主要财务指标]{唐人神（002567）近年主要财务指标（左上：营业收入与归母净利润；右上：ROE 与资产负债率；下：货币资金、有息负债为左轴柱状，经营活动现金流净额为其右轴的蓝点折线。
2021---2025 年为年报口径，2026H1 为 2026 年半年报/半年末，与其前一年报之间留出间隔、以斜纹或空心点区分；半年报数据未年化）}
\end{minipage}
\end{center}

\pfl{财务分析} 2026 年 6 月末资产负债率 71.1\%，较 2025 年末上升 4.7 个百分点（样本中位数 52.2\%，所处三级行业内按由低到高排第\zhnumber{10}位）。 2026 年 6 月末货币资金 21.10 亿元、短期债务（短期借款与一年内到期非流动负债）24.84 亿元、长期债务 36.34 亿元，有息负债合计 61.18 亿元，现金短债比 0.85 倍——覆盖偏紧。 净有息负债 40.08 亿元，相当于其 2026 年上半年营业收入的 37\%。 流动比率 0.82、速动比率 0.51，短期偿付压力较大。 2026 年上半年经营活动现金流净额 -0.66 亿元，经营现金流与利润同时为负，日常运营对债务与股东投入的依赖度较高。 2026 年 6 月末在建工程 1.85 亿元，较上年末增加 0.25 亿元，仍有在建投入。 2026 年 6 月末商誉 1.97 亿元，占归母净资产的 4.2\%，是净资产中最脆弱的一块。 综合看，现金短债比 0.85 倍，短期资金安排偏紧；资产负债率高于样本中位数 18.9 个百分点；有息负债中长期债务占比更高，期限结构对短债滚动的压力较小。

\pfl{重大战略转型} 从公开披露看，公司的战略推进集中在：产能与项目（2 项）：2026 年 4 月，公司推进肉品转型：重点投资建设的茶陵香乡猪生产基地具备年出栏 180 万头产能；融资与资本运作（1 项）：2024 年 12 月，公司变更部分募集资金用途：将“东冲三期生猪养殖基地建设项目”剩余募集资金中的 5，800.00 万元变更至“雅安美神养殖”项目。 公告口径上，近三年（2023 年 9 月---2026 年 9 月）公司披露重大重组与并购类公告 4 条、对外投资与转型类公告 2 条、回购与股权激励类公告 42 条，合计公告 577 条。 第一大行业“饲料”的收入占比由 2021 年年报的 80.5\% 变为 69.0\%（11.5 个百分点），收入结构出现再平衡。 综合判断，报告期内公司有 6 条与战略相关的公告落地（重大重组与并购 4 条、对外投资与转型 2 条），资本运作与主业调整之间存在直接关联。

\begin{center}\small\setlength{\tabcolsep}{4pt}
\captionof{table}[唐人神（002567）历史融资与资本运作]{唐人神（002567）历史融资与资本运作（据公司公告与公开报道整理；金额为披露值，含首次公开发行、再融资、债券、银行借款与重整投资等）}
\begin{adjustbox}{max width=\textwidth}
\begin{tabular}{@{}p{1.45cm}p{2.55cm}p{4.05cm}p{4.95cm}@{}}
\toprule
时间 & 方式 & 规模 & 用途与说明 \\
\midrule
2011-03 & IPO（深交所） & 3，\allowbreak{}500 万股、\allowbreak{}27.00 元/\allowbreak{}股（网下 700 万股、\allowbreak{}网上 2，\allowbreak{}800 万股） & 上市日 2011-03-16 \\
2015-11 & 非公开发行股票 & 7,\allowbreak{}240.95 万股、\allowbreak{}8.01 元/\allowbreak{}股 & 实施 2014 年度权益分配方案后发行数量相应调整；2015-11-12 完成登记 \\
2019-12 & 公开发行可转换公司债券（“唐人转债”128092） & 428.00 亿元（1,\allowbreak{}242.8 万张） & 2020-01-22 上市；截至 2020-10-12（停止转股日）累计转股 1.43 亿股 \\
2021-02 & 非公开发行股票（2020 年度方案） & 2.27 亿股、\allowbreak{}6.83 元/\allowbreak{}股，\allowbreak{}募集资金总额 15.50 亿元、\allowbreak{}净额 15.14 亿元 & 239 股）；2021-02-06 资金到账 \\
2022-12 & 非公开发行股票 & 1.75 亿股、\allowbreak{}6.50 元/\allowbreak{}股，\allowbreak{}募集资金总额 11.40 亿元（净额 11.12 亿元） & 2022-12-22 新增股份上市 \\
2023 & 以简易程序向特定对象发行股票（小额快速定增） & 募集资金 3.00 亿元（净额 2.93 亿元） & 数字化信息系统与智能养殖体系建设；后部分募投项目延期至 2028 年 6 月并调整；2025 年年度报告披露累计使用募集资金 8，\allowbreak{}488.62 万元 \\
2026 & 担保额度 & 2026 年对合并报表内单位担保总额不超过 86.17 亿元（约 86.17 亿元） & 支持子公司保理业务、\allowbreak{}融资租赁、\allowbreak{}采购及履约；额度在 2026-01-01 至 2026-12-31 期间滚动使用 \\
2026-01 & 中期票据（科技创新债券） & 2025 年度第二期科技创新债券：注册金额 5.00 亿元，\allowbreak{}本期发行金额 2.00 亿元 & 募集说明书签署日期 2026 年 1 月 \\
\bottomrule
\end{tabular}
\end{adjustbox}
\end{center}

\pfl{历史融投资情况} 从现金流量表看，2025 年公司取得借款收到的现金 21.96 亿元、偿还债务支付的现金 30.63 亿元、吸收权益性投资 1.92 亿元，筹资活动现金流净额 -12.77 亿元（2023 年为取得借款 44.82 亿元、偿还债务 34.54 亿元）。 2026 年上半年筹资活动现金流净额为 2.91 亿元（取得借款 13.72 亿元、偿还债务 12.04 亿元，未年化），仍处于净融入状态。 2025 年分配股利、利润或偿付利息支付的现金合计 2.35 亿元。 公开披露的融资记录共 8 笔，工具类型以 IPO、定向增发、可转债、中期票据为主；金额与用途见下表。 其中规模最大的两笔为：2026 年 1 月，中期票据，2025 年度第二期科技创新债券：注册金额 5.00 亿元，本期发行金额 2.00 亿元；2023 年，定向增发，募集资金 3.00 亿元（净额 2.93 亿元）。 股本方面，近三年披露股本变动 12 次（限售股份上市 3 次、其他 3 次、期权行权 2 次、增发新股上市 1 次）；总股本由 2021 年末的 12.06 亿股变为 14.33 亿股（+18.83\%）。 分红方面，近三年现金分红 1 次、合计每股派现 0.04 元，最近一次实施在 2022 年 12 月。 近三年“再融资”类公告 12 条，涉及定向增发、募集资金使用。 综合看，公司融资结构上以债务滚动为主、股权融资为补充；2025 年筹资活动现金流净额 -12.77 亿元，融资节奏仍以维持存量债务周转为主，与前述经营与投资现金流的方向基本一致。

\pfl{未来潜在融资需求分析} 2026 年 6 月末货币资金 21.10 亿元，而短期债务 24.84 亿元（现金短债比 0.85 倍），2026 年上半年经营现金流净额 -0.66 亿元。按“短期债务减去账面货币资金”的滚动偿付口径，静态缺口约 3.73 亿元；若再计入上半年亏损的年化额（20.93 亿元），维持一年正常经营所需的外部资金约 3.73---24.66 亿元。 按第 \ref{sec:financing-need} 节的三档划分，公司属于第一档“补血型”：2026 年上半年归母净利润 -10.46 亿元、6 月末资产负债率 71.1\%。 可行的融资路径按现实性排序为：定向增发（需股价配合，且控股股东需放弃优先认购或同步增资）、出售非核心资产与参股股权、债务重组或展期（与银行和债券持有人协商）。 公开披露的后续资金安排线索包括：股权融资线索：2025 年度，第二期科技创新债券（中期票据）2.00 亿元发行；债务与授信安排：2026 年，公司担保总额度 86.17 亿元（报表内）+22.40 亿元（报表外），为子公司融资提供支持。 资金需求还要叠加在建工程：期末在建工程 1.85 亿元、2026 年上半年购建固定资产等支出 1.14 亿元（未年化），这部分支出在周期底部通常会被主动放缓。 综合看，公司的外部资金需求以债务接续为主，滚动偿付口径下的静态缺口约 3.73 亿元，融资节奏将受制于银行续贷意愿与股权融资窗口。




\subsubsection{粤海饲料（001313）}
\label{co:001313}

\pfl{公司基本情况} 1994 年设立至今，粤海饲料是总部位于广东湛江的特种水产饲料龙头，以虾料、海水鱼料为核心，2022 年 2 月登陆深交所主板；2024 年因水产行情低迷陷入亏损，2025 年起业绩修复并推进海外与上游纵向并购。 公司于 2022-02-16 在深交所主板首发上市，注册资本 7.00 亿元，总股本 7.00 亿股。 截至 2026 年 9 月 24 日收盘价 7.31 元，流通市值 51.1 亿元，近一年-0.4\%，流通市值在三级行业“水产饲料”的 4 家公司中按由大到小排第\zhnumber{3}位，近一年涨跌幅高于全样本中位数（-10.1\%）9.7 个百分点。 按 2026 年半年报披露的行业构成，收入几乎全部来自饲料行业 100.0\%。

\pfl{历史股价} 自 2022 年 2 月至 2026 年 9 月，股价由 13.60 元变为 7.31 元，累计 -46.2\%，区间最高 15.15 元（2022 年 3 月）、最低 6.92 元（2024 年 7 月）；当前价相当于区间最高价的 48.3\%，为区间最低价的 1.06 倍。 月度收益的年化波动率 30.2\%，低于全样本中位数 37.4\%，在三级行业“水产饲料”的 4 家公司中波动率由高到低排第\zhnumber{4}位。

\begin{center}
\begin{minipage}{\textwidth}\centering
\begin{tikzpicture}
\begin{axis}[
  width=12.6cm, height=3.9cm, scale only axis,
  xmin=2022.005, xmax=2026.888, ymin=6.44, ymax=16.21,
  xtick={2022,2023,2024,2025,2026},
  yearticks,
  yticklabel style={font=\scriptsize},
  ylabel={元/股}, ylabel style={font=\scriptsize},
  grid=major, grid style={LightGray, line width=0.3pt},
  axis line style={InkGray, line width=0.5pt},
  every axis plot/.append style={line width=0.9pt},
]
\addplot[AgriGreen] table[x=x,y=y]{data/clean/profiles/px_001313.dat};
\addplot[SoftRed, dashed, line width=0.7pt] table[x=x,y=y]{data/clean/profiles/pxzz_001313.dat};
\addplot[only marks, mark=*, mark size=1.1pt, SoftRed] table[x=x,y=y]{data/clean/profiles/pxzz_001313.dat};
\end{axis}
\end{tikzpicture}

\captionof{figure}[粤海饲料（001313）股价与周期骨架]{粤海饲料（001313）月度收盘价（元/股）与 ZigZag 周期骨架（阈值 25\%，圆点为周期转折点）}
\end{minipage}
\end{center}

\pfl{过去周期分析} 以 25\% 的 ZigZag 阈值划分，样本区间内股价形成 1 个主升段与 1 个主跌段。 最大主升段出现在 2024 年 7 月 至 2025 年 2 月，7 个月+31.9\%；最大主跌段为 2022 年 2 月 至 2024 年 7 月，29 个月-49.1\%。 最近一次转折出现在 2025 年 2 月（上涨段自 2024 年 7 月起，7 个月+31.9\%），当前处于该段之后的震荡之中。 公司股价较申万农林牧渔指数提前约 22 个月见底，是板块中较早定价周期反转的公司之一。 水产饲料的需求取决于投苗量与存塘鱼价，第 \ref{sec:other-livestock} 章显示 2026 年淡水鱼价格与投苗量已率先回升，其需求驱动与生猪存栏的联系弱于畜禽饲料，公司 2026 年上半年实现盈利、半年 ROE 1.9\%（未年化），好于全样本中位数（0.75\%），下行周期对其报表的冲击小于同业。 综合区间形态看，该股单段涨跌幅度偏大、以趋势段为主（最大主升 +31.9\%、最大主跌 -49.1\%），年化波动率 30.2\% 与全样本中位数 37.4\% 相比波动小于样本中位数。

\pfl{盈利情况} 2026 年上半年营业收入 37.79 亿元（同比 +41.7\%），归母净利润 0.48 亿元，同比 +1240.5\%。 以年报为趋势对照，2023---2025 年营业收入由 68.72 亿元变为 64.89 亿元（两年年均复合增速 -2.8\%），归母净利润由 0.41 亿元变为 0.18 亿元。 按半年报口径，2026 年上半年的 ROE 为 1.9\%（未年化）、销售净利率 1.4\%、毛利率 11.0\%，ROE 在“水产饲料”4 家公司中按数值由高到低排第\zhnumber{2}位。 同期全样本半年 ROE 中位数 0.75\%，公司与之相差 1.11 个百分点（略好于中位数公司）。 面对这一局面，公司披露的举措包括：公司整体设计产能 250 万吨以上。 综合看，公司在报告期维持盈利（高于全样本半年 ROE 中位数 0.75\%），利润的可持续性取决于水产饲料景气的延续与成本管控的执行。

\begin{center}\small\setlength{\tabcolsep}{3.4pt}
\captionof{table}[粤海饲料（001313）关键财务指标]{粤海饲料（001313）关键财务指标（合并报表口径；两个半年报期均未年化；现金短债比 $=$ 货币资金 $\div$ 短期债务）}
\begin{adjustbox}{max width=\textwidth}
\begin{tabular}{@{}lrrrrr@{}}
\toprule
指标 & 2023 年 & 2024 年 & 2025 年 & 2025 年半年报 & 2026 年半年报 \\
\midrule
营业收入（亿元） & 68.72 & 59.12 & 64.89 & 26.68 & 37.79 \\
归母净利润（亿元） & 0.41 & -0.85 & 0.18 & 0.04 & 0.48 \\
毛利率（\%） & 10.7 & 10.7 & 9.5 & 9.6 & 11.0 \\
ROE（\%） & 1.5 & -3.2 & 0.7 & 0.1 & 1.9 \\
资产负债率（\%） & 41.2 & 46.3 & 46.6 & 45.3 & 52.1 \\
经营活动现金流净额（亿元） & 4.79 & 5.19 & 2.66 & -2.89 & -4.82 \\
货币资金（亿元） & 10.82 & 12.28 & 11.80 & 8.52 & 6.78 \\
有息负债合计（亿元） & 6.51 & 6.58 & 6.17 & 6.81 & 8.40 \\
现金短债比（倍） & 1.87 & 2.38 & 2.55 & 1.62 & 1.23 \\
\bottomrule
\end{tabular}
\end{adjustbox}
\end{center}

\begin{center}
\begin{minipage}{\textwidth}\centering
\begin{tikzpicture}
\begin{groupplot}[
  group style={group size=2 by 1, horizontal sep=0.60cm},
  width=6.85cm, height=4.60cm,
  tick label style={font=\tiny},
  title style={font=\scriptsize\heiti},
  yticklabel style={font=\tiny},
  ylabel style={font=\tiny},
  grid=major, grid style={LightGray, line width=0.3pt},
  axis line style={InkGray, line width=0.5pt},
  enlarge x limits={abs=0.8},
  legend style={font=\scriptsize, at={(0.5,-0.20)}, anchor=north,
                legend columns=2, draw=none, fill=none, column sep=6pt},
]
\nextgroupplot[
  ybar, bar width=4pt,
  ymin=-8.03, ymax=79.53,
  xtick={0,1,2,3,4,5.7},
  xticklabels={2021,2022,2023,2024,2025,2026H1},
  ylabel={亿元},
]
\addplot[fill=AgriGreen!75, draw=AgriGreen, bar shift=-2.6pt] table[x=i, y=rev]{data/clean/profiles/fin_001313.dat};
\addplot[fill=SoftRed!60, draw=SoftRed, bar shift=2.6pt] table[x=i, y=ni]{data/clean/profiles/fin_001313.dat};
\addplot[fill=AgriGreen!30, draw=AgriGreen, pattern=north east lines, bar shift=-2.6pt] table[x=x, y=rev]{data/clean/profiles/finh_001313.dat};
\addplot[fill=SoftRed!20, draw=SoftRed, pattern=north east lines, bar shift=2.6pt] table[x=x, y=ni]{data/clean/profiles/finh_001313.dat};
\legend{营业收入（年/半年）, 归母净利润（年/半年）}
\nextgroupplot[
  ymin=-7.7, ymax=57.6,
  xtick={0,1,2,3,4,5.7},
  xticklabels={2021,2022,2023,2024,2025,2026H1},
  ylabel={\%},
]
\addplot[AgriGold, mark=*, mark size=1.3pt] table[x=i, y=roe]{data/clean/profiles/fin_001313.dat};
\addplot[OilBlue, mark=square*, mark size=1.3pt] table[x=i, y=debt]{data/clean/profiles/fin_001313.dat};
\addplot[only marks, AgriGold, mark=*, mark size=2.0pt] table[x=x, y=roe]{data/clean/profiles/finh_001313.dat};
\addplot[only marks, OilBlue, mark=square*, mark size=2.0pt] table[x=x, y=debt]{data/clean/profiles/finh_001313.dat};
\legend{ROE, 资产负债率}
\end{groupplot}
\end{tikzpicture}
\begin{tikzpicture}
\begin{axis}[
  name=finbot,
  width=13.75cm, height=3.10cm, scale only axis,
  ymin=-1.23, ymax=13.75,
  xtick={0,1,2,3,4,5.7},
  xticklabels={2021,2022,2023,2024,2025,2026H1},
  tick label style={font=\tiny},
  yticklabel style={font=\tiny},
  ylabel={亿元}, ylabel style={font=\tiny},
  ybar, bar width=4pt,
  grid=major, grid style={LightGray, line width=0.3pt},
  axis line style={InkGray, line width=0.5pt},
  enlarge x limits={abs=0.8},
  legend style={font=\scriptsize, at={(0.5,-0.26)}, anchor=north,
                legend columns=2, draw=none, fill=none, column sep=10pt},
]
\addplot[fill=AgriGreen!55, draw=AgriGreen, bar shift=-3.0pt] table[x=i, y=cash]{data/clean/profiles/fin_001313.dat};
\addplot[fill=SoftRed!45, draw=SoftRed, bar shift=3.0pt] table[x=i, y=ibd]{data/clean/profiles/fin_001313.dat};
\addplot[fill=AgriGreen!25, draw=AgriGreen, pattern=north east lines, bar shift=-3.0pt] table[x=x, y=cash]{data/clean/profiles/finh_001313.dat};
\addplot[fill=SoftRed!20, draw=SoftRed, pattern=north east lines, bar shift=3.0pt] table[x=x, y=ibd]{data/clean/profiles/finh_001313.dat};
\legend{货币资金（左轴）, 有息负债（左轴）}
\end{axis}
\begin{axis}[
  at={(finbot.south east)}, anchor=south east,
  width=13.75cm, height=3.10cm, scale only axis,
  axis y line*=right, axis x line*=none, xtick=\empty,
  ymin=-7.71, ymax=9.64,
  tick label style={font=\tiny},
  yticklabel style={font=\tiny}, scaled y ticks=false,
  legend style={font=\scriptsize, at={(0.5,-0.44)}, anchor=north,
                legend columns=1, draw=none, fill=none},
]
\addplot[OilBlue, mark=*, mark size=2.2pt, line width=1.3pt] table[x=i, y=ocf]{data/clean/profiles/fin_001313.dat};
\addplot[only marks, OilBlue, mark=*, mark size=3.2pt] table[x=x, y=ocf]{data/clean/profiles/finh_001313.dat};
\legend{经营活动现金流净额（右轴，亿元，蓝点）}
\end{axis}
\end{tikzpicture}

\captionof{figure}[粤海饲料（001313）近年主要财务指标]{粤海饲料（001313）近年主要财务指标（左上：营业收入与归母净利润；右上：ROE 与资产负债率；下：货币资金、有息负债为左轴柱状，经营活动现金流净额为其右轴的蓝点折线。
2021---2025 年为年报口径，2026H1 为 2026 年半年报/半年末，与其前一年报之间留出间隔、以斜纹或空心点区分；半年报数据未年化）}
\end{minipage}
\end{center}

\pfl{财务分析} 2026 年 6 月末资产负债率 52.1\%，较 2025 年末上升 5.5 个百分点（样本中位数 52.2\%，所处三级行业内按由低到高排第\zhnumber{1}位）。 2026 年 6 月末货币资金 6.78 亿元、短期债务（短期借款与一年内到期非流动负债）5.52 亿元、长期债务 2.89 亿元，有息负债合计 8.40 亿元，现金短债比 1.23 倍——覆盖充分。 净有息负债 1.62 亿元，相当于其 2026 年上半年营业收入的 4\%。 流动比率 1.46、速动比率 1.12，短期偿债能力处于正常区间。 2026 年上半年经营活动现金流净额 -4.82 亿元，净利润为正但经营现金流为负，利润的现金含量偏低。 综合看，现金短债比 1.23 倍，短期偿付压力可控；资产负债率低于样本中位数 0.1 个百分点；有息负债中短期债务占比更高，期限结构仍以滚动续作为主。

\pfl{重大战略转型} 公司的战略动作可归为以下几类：海外与国际化（2 项）：2026 年上半年，公司以越南为支点辐射东南亚市场，并规划在西南、华北、西北及海外其他养殖区域继续布局产能；2025 年 10 月，总体规划 20 万吨产能，一期工程产能 10 万吨；其他动作（1 项）：2025 年 7 月，第四届董事会第二次会议审议通过以固定资产对全资子公司进行增资，以固定资产对全资子公司增资。 公告口径上，近三年（2023 年 9 月---2026 年 9 月）公司披露重大重组与并购类公告 6 条、对外投资与转型类公告 3 条、回购与股权激励类公告 48 条，合计公告 464 条。 第一大行业仍是“饲料行业”，收入占比 100.0\%，与 2021 年年报相比没有方向性变化。 综合判断，报告期内公司有 9 条与战略相关的公告落地（重大重组与并购 6 条、对外投资与转型 3 条），资本运作与主业调整之间存在直接关联。

\pfl{历史融投资情况} 从现金流量表看，2025 年公司取得借款收到的现金 4.73 亿元、偿还债务支付的现金 5.36 亿元、吸收权益性投资 0.03 亿元，筹资活动现金流净额 -0.99 亿元（2023 年为取得借款 8.08 亿元、偿还债务 11.11 亿元）。 2026 年上半年筹资活动现金流净额为 0.24 亿元（取得借款 5.42 亿元、偿还债务 4.41 亿元，未年化），仍处于净融入状态。 2025 年分配股利、利润或偿付利息支付的现金合计 1.06 亿元。 公开披露的融资记录共 2 笔，工具类型以 IPO 为主；金额与用途见下表。 其中规模最大的两笔为：2022 年 2 月，IPO，公开发行 10，000 万股，发行价 5.38 元/股，募集资金总额 5.38 亿元；2025 年 11 月，现金收购股权，32.00 亿元收购宜兴市天石饲料有限公司 60\% 股权。 股本方面，近三年披露股本变动 5 次（限售股份上市 1 次、其他 1 次）；总股本由 2021 年末的 6.00 亿股变为 7.00 亿股（+16.67\%）。 分红方面，近三年现金分红 3 次、合计每股派现 0.21 元，最近一次实施在 2025 年 6 月。 综合看，公司融资结构上以债务滚动为主、股权融资为补充；2025 年筹资活动现金流净额 -0.99 亿元，融资节奏仍以维持存量债务周转为主，与前述经营与投资现金流的方向基本一致。

\pfl{未来潜在融资需求分析} 2026 年 6 月末货币资金 6.78 亿元，而短期债务 5.52 亿元（现金短债比 1.23 倍），2026 年上半年经营现金流净额 -4.82 亿元。账面货币资金可以覆盖短期债务，缺口不体现在静态偿付层面。 公司 2026 年 6 月末资产负债率 52.1\%、2026 年上半年 ROE 1.9\%（未年化），财务结构相对宽松而盈利弹性有限，融资属于可选项而非必选项，触发条件多与产能建设或外延并购相关。 静态偿付不存在缺口，后续融资更可能服务于产能或并购安排，而不是填补流动性窟窿。




\subsubsection{邦基科技（603151）}
\label{co:603151}

\pfl{公司基本情况} 公司前身可追溯至 2007 年，山东淄博的猪饲料专业化企业，以猪预混料、浓缩料、配合料为主业，2022 年 10 月登陆上交所主板；2025—2026 年因向生猪养殖下游延伸并受猪价下行冲击，业绩由盈转亏。 公司于 2022-10-19 在上交所首发上市，注册资本 1.71 亿元，总股本 1.68 亿股。 截至 2026 年 9 月 24 日收盘价 18.32 元，流通市值 31.4 亿元，近一年-28.9\%，流通市值在三级行业“畜禽饲料”的 11 家公司中按由大到小排第\zhnumber{7}位，近一年涨跌幅低于全样本中位数（-10.1\%）18.8 个百分点。 按 2025 年年报披露的行业构成，收入占比前三位为饲料行业 96.1\%；畜牧养殖 2.3\%；其他(补充) 1.6\%。

\pfl{历史股价} 自 2022 年 10 月至 2026 年 9 月，股价由 17.82 元变为 18.32 元，累计 +2.8\%，区间最高 25.74 元（2025 年 10 月）、最低 9.08 元（2024 年 8 月）；当前价相当于区间最高价的 71.2\%，为区间最低价的 2.02 倍。 月度收益的年化波动率 45.6\%，高于全样本中位数 37.4\%，在三级行业“畜禽饲料”的 11 家公司中波动率由高到低排第\zhnumber{3}位。 把股价阶段与公司事件对齐看，2024 年 8 月 至 2025 年 10 月股价上涨+183.5\%，同期（2025 年 6 月）发行股份及支付现金购买资产暨关联交易预案；2026 年 6 月 至 2026 年 9 月股价上涨+66.7\%，同期（2026 年）归属于上市公司股东的净利润 -47.97 万元（同比 -100.74\%）。 把这些阶段与公司事件对齐后可以看到，几轮大涨大跌多由公司自身的资本运作与风险事件触发，行业景气只解释了其中一部分。

\begin{center}
\begin{minipage}{\textwidth}\centering
\begin{tikzpicture}
\begin{axis}[
  width=12.6cm, height=3.9cm, scale only axis,
  xmin=2022.672, xmax=2026.888, ymin=8.44, ymax=27.54,
  xtick={2022,2023,2024,2025,2026},
  yearticks,
  yticklabel style={font=\scriptsize},
  ylabel={元/股}, ylabel style={font=\scriptsize},
  grid=major, grid style={LightGray, line width=0.3pt},
  axis line style={InkGray, line width=0.5pt},
  every axis plot/.append style={line width=0.9pt},
]
\addplot[AgriGreen] table[x=x,y=y]{data/clean/profiles/px_603151.dat};
\addplot[SoftRed, dashed, line width=0.7pt] table[x=x,y=y]{data/clean/profiles/pxzz_603151.dat};
\addplot[only marks, mark=*, mark size=1.1pt, SoftRed] table[x=x,y=y]{data/clean/profiles/pxzz_603151.dat};
\end{axis}
\end{tikzpicture}

\captionof{figure}[邦基科技（603151）股价与周期骨架]{邦基科技（603151）月度收盘价（元/股）与 ZigZag 周期骨架（阈值 25\%，圆点为周期转折点）}
\end{minipage}
\end{center}

\pfl{过去周期分析} 以 25\% 的 ZigZag 阈值划分，样本区间内股价形成 3 个主升段与 3 个主跌段。 最大主升段出现在 2024 年 8 月 至 2025 年 10 月，14 个月+183.5\%；最大主跌段为 2025 年 10 月 至 2026 年 6 月，8 个月-57.3\%。 最近一次转折出现在 2026 年 9 月（上涨段自 2026 年 6 月起，3 个月+66.7\%），当前处于该段之后的震荡之中。 公司股价较申万农林牧渔指数提前约 21 个月见底，是板块中较早定价周期反转的公司之一。 饲料夹在原料成本与养殖需求之间（第 \ref{sec:feed} 节），量的部分取决于存栏、利的部分取决于原料与成品的价差，公司 2026 年 6 月末资产负债率 63.4\%、半年 ROE -0.0\%（未年化），在本轮下行中属于随行业同步调整的一类。 综合区间形态看，该股单段涨跌幅度偏大、以趋势段为主（最大主升 +183.5\%、最大主跌 -57.3\%），年化波动率 45.6\% 与全样本中位数 37.4\% 相比波动明显大于样本中位数。

\pfl{盈利情况} 2026 年上半年营业收入 26.75 亿元（同比 +11.8\%），归母净利润 -0.00 亿元，由上年同期的盈利转为亏损。 以年报为趋势对照，2023---2025 年营业收入由 16.47 亿元变为 56.99 亿元（两年年均复合增速 +86.0\%），归母净利润由 0.84 亿元变为 1.11 亿元。 按半年报口径，2026 年上半年的 ROE 为 -0.0\%（未年化）、销售净利率 -1.3\%、毛利率 5.9\%，ROE 在“畜禽饲料”11 家公司中按数值由高到低排第\zhnumber{3}位。 按同一口径，三级行业“畜禽饲料”11 家公司的半年 ROE 中位数为 -6.25\%，该公司高于其中位数 6.22 个百分点。 综合看，公司仍处亏损区间、由盈转亏，半年 ROE 在三级行业内按由高到低排第\zhnumber{3}位，单位盈利能力的修复依赖行业景气与自身成本端的改善，报表弹性主要体现为价格与销量的方向性变化。

\begin{center}\small\setlength{\tabcolsep}{3.4pt}
\captionof{table}[邦基科技（603151）关键财务指标]{邦基科技（603151）关键财务指标（合并报表口径；两个半年报期均未年化；现金短债比 $=$ 货币资金 $\div$ 短期债务）}
\begin{adjustbox}{max width=\textwidth}
\begin{tabular}{@{}lrrrrr@{}}
\toprule
指标 & 2023 年 & 2024 年 & 2025 年 & 2025 年半年报 & 2026 年半年报 \\
\midrule
营业收入（亿元） & 16.47 & 25.42 & 56.99 & 23.93 & 26.75 \\
归母净利润（亿元） & 0.84 & 0.50 & 1.11 & 0.64 & -0.00 \\
毛利率（\%） & 13.0 & 9.2 & 7.6 & 8.7 & 5.9 \\
ROE（\%） & 6.9 & 4.1 & 8.5 & 4.9 & -0.0 \\
资产负债率（\%） & 15.9 & 28.8 & 55.9 & 50.7 & 63.4 \\
经营活动现金流净额（亿元） & -0.54 & -1.10 & -6.01 & -5.48 & 1.44 \\
货币资金（亿元） & 3.91 & 1.68 & 3.40 & 2.53 & 7.22 \\
有息负债合计（亿元） & 1.03 & 2.46 & 11.91 & 8.42 & 15.02 \\
现金短债比（倍） & 4.26 & 0.72 & 0.32 & 0.31 & 0.55 \\
\bottomrule
\end{tabular}
\end{adjustbox}
\end{center}

\begin{center}
\begin{minipage}{\textwidth}\centering
\begin{tikzpicture}
\begin{groupplot}[
  group style={group size=2 by 1, horizontal sep=0.60cm},
  width=6.85cm, height=4.60cm,
  tick label style={font=\tiny},
  title style={font=\scriptsize\heiti},
  yticklabel style={font=\tiny},
  ylabel style={font=\tiny},
  grid=major, grid style={LightGray, line width=0.3pt},
  axis line style={InkGray, line width=0.5pt},
  enlarge x limits={abs=0.8},
  legend style={font=\scriptsize, at={(0.5,-0.20)}, anchor=north,
                legend columns=2, draw=none, fill=none, column sep=6pt},
]
\nextgroupplot[
  ybar, bar width=4pt,
  ymin=-5.70, ymax=63.83,
  xtick={0,1,2,3,4,5.7},
  xticklabels={2021,2022,2023,2024,2025,2026H1},
  ylabel={亿元},
]
\addplot[fill=AgriGreen!75, draw=AgriGreen, bar shift=-2.6pt] table[x=i, y=rev]{data/clean/profiles/fin_603151.dat};
\addplot[fill=SoftRed!60, draw=SoftRed, bar shift=2.6pt] table[x=i, y=ni]{data/clean/profiles/fin_603151.dat};
\addplot[fill=AgriGreen!30, draw=AgriGreen, pattern=north east lines, bar shift=-2.6pt] table[x=x, y=rev]{data/clean/profiles/finh_603151.dat};
\addplot[fill=SoftRed!20, draw=SoftRed, pattern=north east lines, bar shift=2.6pt] table[x=x, y=ni]{data/clean/profiles/finh_603151.dat};
\legend{营业收入（年/半年）, 归母净利润（年/半年）}
\nextgroupplot[
  ymin=-5.1, ymax=69.7,
  xtick={0,1,2,3,4,5.7},
  xticklabels={2021,2022,2023,2024,2025,2026H1},
  ylabel={\%},
]
\addplot[AgriGold, mark=*, mark size=1.3pt] table[x=i, y=roe]{data/clean/profiles/fin_603151.dat};
\addplot[OilBlue, mark=square*, mark size=1.3pt] table[x=i, y=debt]{data/clean/profiles/fin_603151.dat};
\addplot[only marks, AgriGold, mark=*, mark size=2.0pt] table[x=x, y=roe]{data/clean/profiles/finh_603151.dat};
\addplot[only marks, OilBlue, mark=square*, mark size=2.0pt] table[x=x, y=debt]{data/clean/profiles/finh_603151.dat};
\legend{ROE, 资产负债率}
\end{groupplot}
\end{tikzpicture}
\begin{tikzpicture}
\begin{axis}[
  name=finbot,
  width=13.75cm, height=3.10cm, scale only axis,
  ymin=-1.50, ymax=16.82,
  xtick={0,1,2,3,4,5.7},
  xticklabels={2021,2022,2023,2024,2025,2026H1},
  tick label style={font=\tiny},
  yticklabel style={font=\tiny},
  ylabel={亿元}, ylabel style={font=\tiny},
  ybar, bar width=4pt,
  grid=major, grid style={LightGray, line width=0.3pt},
  axis line style={InkGray, line width=0.5pt},
  enlarge x limits={abs=0.8},
  legend style={font=\scriptsize, at={(0.5,-0.26)}, anchor=north,
                legend columns=2, draw=none, fill=none, column sep=10pt},
]
\addplot[fill=AgriGreen!55, draw=AgriGreen, bar shift=-3.0pt] table[x=i, y=cash]{data/clean/profiles/fin_603151.dat};
\addplot[fill=SoftRed!45, draw=SoftRed, bar shift=3.0pt] table[x=i, y=ibd]{data/clean/profiles/fin_603151.dat};
\addplot[fill=AgriGreen!25, draw=AgriGreen, pattern=north east lines, bar shift=-3.0pt] table[x=x, y=cash]{data/clean/profiles/finh_603151.dat};
\addplot[fill=SoftRed!20, draw=SoftRed, pattern=north east lines, bar shift=3.0pt] table[x=x, y=ibd]{data/clean/profiles/finh_603151.dat};
\legend{货币资金（左轴）, 有息负债（左轴）}
\end{axis}
\begin{axis}[
  at={(finbot.south east)}, anchor=south east,
  width=13.75cm, height=3.10cm, scale only axis,
  axis y line*=right, axis x line*=none, xtick=\empty,
  ymin=-7.94, ymax=3.67,
  tick label style={font=\tiny},
  yticklabel style={font=\tiny}, scaled y ticks=false,
  legend style={font=\scriptsize, at={(0.5,-0.44)}, anchor=north,
                legend columns=1, draw=none, fill=none},
]
\addplot[OilBlue, mark=*, mark size=2.2pt, line width=1.3pt] table[x=i, y=ocf]{data/clean/profiles/fin_603151.dat};
\addplot[only marks, OilBlue, mark=*, mark size=3.2pt] table[x=x, y=ocf]{data/clean/profiles/finh_603151.dat};
\legend{经营活动现金流净额（右轴，亿元，蓝点）}
\end{axis}
\end{tikzpicture}

\captionof{figure}[邦基科技（603151）近年主要财务指标]{邦基科技（603151）近年主要财务指标（左上：营业收入与归母净利润；右上：ROE 与资产负债率；下：货币资金、有息负债为左轴柱状，经营活动现金流净额为其右轴的蓝点折线。
2021---2025 年为年报口径，2026H1 为 2026 年半年报/半年末，与其前一年报之间留出间隔、以斜纹或空心点区分；半年报数据未年化）}
\end{minipage}
\end{center}

\pfl{财务分析} 2026 年 6 月末资产负债率 63.4\%，较 2025 年末上升 7.5 个百分点（样本中位数 52.2\%，所处三级行业内按由低到高排第\zhnumber{7}位）。 2026 年 6 月末货币资金 7.22 亿元、短期债务（短期借款与一年内到期非流动负债）13.06 亿元、长期债务 1.96 亿元，有息负债合计 15.02 亿元，现金短债比 0.55 倍——覆盖偏紧。 净有息负债 7.79 亿元，相当于其 2026 年上半年营业收入的 29\%。 流动比率 1.28、速动比率 1.01，短期指标尚可。 2026 年上半年经营活动现金流净额 1.44 亿元，净利润为负而经营现金流为正，差额主要来自折旧摊销与营运资本变动，说明亏损尚未完全转化为现金流出。 综合看，现金短债比 0.55 倍，短期资金安排偏紧；资产负债率高于样本中位数 11.2 个百分点；有息负债中短期债务占比更高，期限结构仍以滚动续作为主。

\pfl{重大战略转型} 从公开披露看，公司的战略推进集中在：其他动作（2 项）：2024 年 6 月，推出 2024 年股票期权激励计划，以 12.85 元/股向 154 名激励对象授予 875.00 万份股票期权；并购与重组（2 项）：2025 年 11 月，公司终止上述重大资产重组；2025 年 6 月，发行股份及支付现金购买资产暨关联交易预案。 公告口径上，近三年（2023 年 9 月---2026 年 9 月）公司披露重大重组与并购类公告 16 条、对外投资与转型类公告 4 条、回购与股权激励类公告 0 条，合计公告 405 条。 第一大行业仍是“饲料行业”，收入占比 96.1\%，与 2021 年年报相比没有方向性变化。 综合判断，报告期内公司有 20 条与战略相关的公告落地（重大重组与并购 16 条、对外投资与转型 4 条），资本运作与主业调整之间存在直接关联。

\begin{center}\small\setlength{\tabcolsep}{4pt}
\captionof{table}[邦基科技（603151）历史融资与资本运作]{邦基科技（603151）历史融资与资本运作（据公司公告与公开报道整理；金额为披露值，含首次公开发行、再融资、债券、银行借款与重整投资等）}
\begin{adjustbox}{max width=\textwidth}
\begin{tabular}{@{}p{1.45cm}p{2.55cm}p{4.05cm}p{4.95cm}@{}}
\toprule
时间 & 方式 & 规模 & 用途与说明 \\
\midrule
2022-10 & IPO（上交所主板） & 公开发行 4，\allowbreak{}200.00 万股，\allowbreak{}发行价 17.95 元/\allowbreak{}股，\allowbreak{}募集资金总额 7.54 亿元 & 投向公司及辽宁、\allowbreak{}长春、\allowbreak{}云南、\allowbreak{}张家口、\allowbreak{}山西等地子公司的高档配合猪饲料、\allowbreak{}浓缩饲料智能生产车间等项目；募集资金 2022 年 10 月 14 日到账；项目总投资 6.90 亿元 \\
2024-06-06 & 股票期权激励（授予） & 以 12.85 元/\allowbreak{}份授予 154 名激励对象 875.00 万份股票期权 & 员工长期激励；2025 年 5 月 13 日首次授予第一个行权期可行权数量 332.00 万股（140 人） \\
2025-06-16 & 发行股份及支付现金购买资产（重大资产重组，后终止） & 拟收购山东北溪农牧、\allowbreak{}山东瑞东伟力农牧、\allowbreak{}山东鑫牧农牧、\allowbreak{}瑞东农牧（利津）、\allowbreak{}瑞东农牧（山东）、\allowbreak{}瑞东威力牧业（滨州）各 100\% 股权及派斯东畜牧技术咨询（上海）80\% 股权 & 由单一猪饲料业务向下游养殖业延伸；交易对方为 Riverstone Farm Pte. Ltd \\
\bottomrule
\end{tabular}
\end{adjustbox}
\end{center}

\pfl{历史融投资情况} 从现金流量表看，2025 年公司取得借款收到的现金 11.95 亿元、偿还债务支付的现金 4.19 亿元、吸收权益性投资 1.42 亿元，筹资活动现金流净额 8.98 亿元（2023 年为取得借款 0.94 亿元、偿还债务 0.50 亿元）。 2026 年上半年筹资活动现金流净额为 2.83 亿元（取得借款 8.68 亿元、偿还债务 6.63 亿元，未年化），仍处于净融入状态。 2025 年分配股利、利润或偿付利息支付的现金合计 0.58 亿元。 公开披露的融资记录共 3 笔，工具类型以 IPO、发行股份及支付现金购为主；金额与用途见下表。 股本方面，近三年披露股本变动 4 次（限售股份上市 1 次）；总股本自 2021 年末以来未发生变化（1.68 亿股）。 分红方面，近三年现金分红 6 次、合计每股派现 1.05 元，最近一次实施在 2026 年 6 月。 近三年“再融资”类公告 4 条，涉及定向增发。 综合看，公司融资结构上以债务滚动为主、股权融资为补充；2025 年筹资活动现金流净额 8.98 亿元，年度内仍为净融入，与前述经营与投资现金流的方向基本一致。

\pfl{未来潜在融资需求分析} 2026 年 6 月末货币资金 7.22 亿元，而短期债务 13.06 亿元（现金短债比 0.55 倍），2026 年上半年经营现金流净额 1.44 亿元。按“短期债务减去账面货币资金”的滚动偿付口径，静态缺口约 5.83 亿元；上半年盈利或接近盈亏平衡，该缺口不随经营亏损进一步扩大。 按第 \ref{sec:financing-need} 节的三档划分，公司属于第二档“保壳型”：近三年重大重组与并购类公告 16 条，2026 年上半年归母净利润 -0.00 亿元，融资需求更多体现为“资产置换 + 维持上市地位”，而不是扩产。 公开披露的后续资金安排线索包括：债务与授信安排：2025 年年度报告，公司披露公司债券（含企业债券）和非金融企业债务融资工具、可转换公司债券均不适用。 综合看，公司的外部资金需求以补充营运资金为主、项目投入为辅，滚动偿付口径下的静态缺口约 5.83 亿元，融资节奏将受制于银行续贷意愿与股权融资窗口。




\subsubsection{金新农（002548）}
\label{co:002548}

\pfl{公司基本情况} 公司前身可追溯至 2008 年，深圳的饲料与生猪养殖双主业企业，2011 年上市，历经 2019 年、2020 年两轮控制权与股东结构变动后由广州金农产投控股；2026 年上半年猪价创近十年新低，归母净利润亏损 3.88 亿元，净资产大幅收缩。 公司于 2011-02-18 在深交所主板首发上市，注册资本 8.28 亿元，总股本 8.05 亿股。 截至 2026 年 9 月 24 日收盘价 3.82 元，流通市值 30.7 亿元，近一年-29.1\%，流通市值在三级行业“畜禽饲料”的 11 家公司中按由大到小排第\zhnumber{8}位，近一年涨跌幅低于全样本中位数（-10.1\%）19.0 个百分点。 按 2026 年半年报披露的行业构成，收入占比前三位为饲料加工 69.7\%；畜牧养殖 28.8\%；其他 1.5\%。

\pfl{历史股价} 本报告样本区间（2021 年 1 月 至 2026 年 9 月）内，股价由 8.03 元变为 3.82 元，累计 -52.4\%，区间最高 8.16 元（2022 年 3 月）、最低 3.69 元（2026 年 6 月）；当前价相当于区间最高价的 46.8\%，已接近区间最低价（1.04 倍）。 月度收益的年化波动率 40.3\%，高于全样本中位数 37.4\%，在三级行业“畜禽饲料”的 11 家公司中波动率由高到低排第\zhnumber{5}位。 把股价阶段与公司事件对齐看，2024 年 8 月 至 2026 年 2 月股价上涨+93.0\%，同期（2025 年）半年度末资产权利受限合计账面余额 10.90 亿元。 从时间对应关系看，公司层面的资本运作与股权、风险事件解释了区间内多数急涨急跌，与行业指数的同步性相对有限。

\begin{center}
\begin{minipage}{\textwidth}\centering
\begin{tikzpicture}
\begin{axis}[
  width=12.6cm, height=3.9cm, scale only axis,
  xmin=2020.922, xmax=2026.888, ymin=3.43, ymax=8.73,
  xtick={2021,2022,2023,2024,2025,2026},
  yearticks,
  yticklabel style={font=\scriptsize},
  ylabel={元/股}, ylabel style={font=\scriptsize},
  grid=major, grid style={LightGray, line width=0.3pt},
  axis line style={InkGray, line width=0.5pt},
  every axis plot/.append style={line width=0.9pt},
]
\addplot[AgriGreen] table[x=x,y=y]{data/clean/profiles/px_002548.dat};
\addplot[SoftRed, dashed, line width=0.7pt] table[x=x,y=y]{data/clean/profiles/pxzz_002548.dat};
\addplot[only marks, mark=*, mark size=1.1pt, SoftRed] table[x=x,y=y]{data/clean/profiles/pxzz_002548.dat};
\end{axis}
\end{tikzpicture}

\captionof{figure}[金新农（002548）股价与周期骨架]{金新农（002548）月度收盘价（元/股）与 ZigZag 周期骨架（阈值 25\%，圆点为周期转折点）}
\end{minipage}
\end{center}

\pfl{过去周期分析} 以 25\% 的 ZigZag 阈值划分，样本区间内股价形成 3 个主升段与 4 个主跌段。 最大主跌段为 2023 年 6 月 至 2024 年 8 月，14 个月-51.4\%；最大主升段出现在 2024 年 8 月 至 2026 年 2 月，18 个月+93.0\%。 最近一次转折出现在 2026 年 6 月（下跌段自 2026 年 2 月起，4 个月-48.6\%），当前处于该段的延续之中。 公司股价与申万农林牧渔指数几乎同步见底，个股并未走出独立行情。 饲料夹在原料成本与养殖需求之间（第 \ref{sec:feed} 节），量的部分取决于存栏、利的部分取决于原料与成品的价差，公司 2026 年 6 月末资产负债率 81.5\%，高于全样本中位数 52.2\%，其股价因此更多反映资金面与信用风险，而不是价格弹性本身。 综合区间形态看，该股单段涨跌幅度偏大、以趋势段为主（最大主升 +93.0\%、最大主跌 -51.4\%），年化波动率 40.3\% 与全样本中位数 37.4\% 相比波动与样本中位数接近。

\pfl{盈利情况} 2026 年上半年营业收入 23.20 亿元（同比 -2.4\%），归母净利润 -3.88 亿元，亏损同比扩大 3.65 亿元。 以年报为趋势对照，2023---2025 年营业收入由 40.40 亿元变为 47.46 亿元（两年年均复合增速 +8.4\%），归母净利润由 -6.60 亿元变为 -2.93 亿元。 按半年报口径，2026 年上半年的 ROE 为 -37.4\%（未年化）、销售净利率 -16.8\%、毛利率 -2.5\%，ROE 在“畜禽饲料”11 家公司中按数值由高到低排第\zhnumber{11}位。 同期全样本半年 ROE 中位数 0.75\%，公司与之相差 38.19 个百分点（弱于中位数公司）。 综合看，公司仍处亏损区间、亏损同比扩大，半年 ROE 在三级行业内按由高到低排第\zhnumber{11}位，单位盈利能力的修复依赖行业景气与自身成本端的改善，报表弹性主要体现为价格与销量的方向性变化。

\begin{center}\small\setlength{\tabcolsep}{3.4pt}
\captionof{table}[金新农（002548）关键财务指标]{金新农（002548）关键财务指标（合并报表口径；两个半年报期均未年化；现金短债比 $=$ 货币资金 $\div$ 短期债务）}
\begin{adjustbox}{max width=\textwidth}
\begin{tabular}{@{}lrrrrr@{}}
\toprule
指标 & 2023 年 & 2024 年 & 2025 年 & 2025 年半年报 & 2026 年半年报 \\
\midrule
营业收入（亿元） & 40.40 & 45.62 & 47.46 & 23.76 & 23.20 \\
归母净利润（亿元） & -6.60 & 0.30 & -2.93 & -0.24 & -3.88 \\
毛利率（\%） & 3.3 & 10.6 & 4.8 & 7.5 & -2.5 \\
ROE（\%） & -36.5 & 2.0 & -21.4 & -1.6 & -37.4 \\
资产负债率（\%） & 74.1 & 70.4 & 75.7 & 72.1 & 81.5 \\
经营活动现金流净额（亿元） & 1.25 & 4.23 & 3.32 & 1.45 & 0.50 \\
货币资金（亿元） & 7.54 & 3.38 & 3.05 & 5.11 & 2.36 \\
有息负债合计（亿元） & 29.75 & 27.57 & 28.60 & 30.67 & 28.92 \\
现金短债比（倍） & 0.36 & 0.17 & 0.14 & 0.22 & 0.10 \\
\bottomrule
\end{tabular}
\end{adjustbox}
\end{center}

\begin{center}
\begin{minipage}{\textwidth}\centering
\begin{tikzpicture}
\begin{groupplot}[
  group style={group size=2 by 1, horizontal sep=0.60cm},
  width=6.85cm, height=4.60cm,
  tick label style={font=\tiny},
  title style={font=\scriptsize\heiti},
  yticklabel style={font=\tiny},
  ylabel style={font=\tiny},
  grid=major, grid style={LightGray, line width=0.3pt},
  axis line style={InkGray, line width=0.5pt},
  enlarge x limits={abs=0.8},
  legend style={font=\scriptsize, at={(0.5,-0.20)}, anchor=north,
                legend columns=2, draw=none, fill=none, column sep=6pt},
]
\nextgroupplot[
  ybar, bar width=4pt,
  ymin=-15.78, ymax=55.70,
  xtick={0,1,2,3,4,5.7},
  xticklabels={2021,2022,2023,2024,2025,2026H1},
  ylabel={亿元},
]
\addplot[fill=AgriGreen!75, draw=AgriGreen, bar shift=-2.6pt] table[x=i, y=rev]{data/clean/profiles/fin_002548.dat};
\addplot[fill=SoftRed!60, draw=SoftRed, bar shift=2.6pt] table[x=i, y=ni]{data/clean/profiles/fin_002548.dat};
\addplot[fill=AgriGreen!30, draw=AgriGreen, pattern=north east lines, bar shift=-2.6pt] table[x=x, y=rev]{data/clean/profiles/finh_002548.dat};
\addplot[fill=SoftRed!20, draw=SoftRed, pattern=north east lines, bar shift=2.6pt] table[x=x, y=ni]{data/clean/profiles/finh_002548.dat};
\legend{营业收入（年/半年）, 归母净利润（年/半年）}
\nextgroupplot[
  ymin=-52.1, ymax=93.9,
  xtick={0,1,2,3,4,5.7},
  xticklabels={2021,2022,2023,2024,2025,2026H1},
  ylabel={\%},
]
\addplot[AgriGold, mark=*, mark size=1.3pt] table[x=i, y=roe]{data/clean/profiles/fin_002548.dat};
\addplot[OilBlue, mark=square*, mark size=1.3pt] table[x=i, y=debt]{data/clean/profiles/fin_002548.dat};
\addplot[only marks, AgriGold, mark=*, mark size=2.0pt] table[x=x, y=roe]{data/clean/profiles/finh_002548.dat};
\addplot[only marks, OilBlue, mark=square*, mark size=2.0pt] table[x=x, y=debt]{data/clean/profiles/finh_002548.dat};
\legend{ROE, 资产负债率}
\end{groupplot}
\end{tikzpicture}
\begin{tikzpicture}
\begin{axis}[
  name=finbot,
  width=13.75cm, height=3.10cm, scale only axis,
  ymin=-3.12, ymax=34.93,
  xtick={0,1,2,3,4,5.7},
  xticklabels={2021,2022,2023,2024,2025,2026H1},
  tick label style={font=\tiny},
  yticklabel style={font=\tiny},
  ylabel={亿元}, ylabel style={font=\tiny},
  ybar, bar width=4pt,
  grid=major, grid style={LightGray, line width=0.3pt},
  axis line style={InkGray, line width=0.5pt},
  enlarge x limits={abs=0.8},
  legend style={font=\scriptsize, at={(0.5,-0.26)}, anchor=north,
                legend columns=2, draw=none, fill=none, column sep=10pt},
]
\addplot[fill=AgriGreen!55, draw=AgriGreen, bar shift=-3.0pt] table[x=i, y=cash]{data/clean/profiles/fin_002548.dat};
\addplot[fill=SoftRed!45, draw=SoftRed, bar shift=3.0pt] table[x=i, y=ibd]{data/clean/profiles/fin_002548.dat};
\addplot[fill=AgriGreen!25, draw=AgriGreen, pattern=north east lines, bar shift=-3.0pt] table[x=x, y=cash]{data/clean/profiles/finh_002548.dat};
\addplot[fill=SoftRed!20, draw=SoftRed, pattern=north east lines, bar shift=3.0pt] table[x=x, y=ibd]{data/clean/profiles/finh_002548.dat};
\legend{货币资金（左轴）, 有息负债（左轴）}
\end{axis}
\begin{axis}[
  at={(finbot.south east)}, anchor=south east,
  width=13.75cm, height=3.10cm, scale only axis,
  axis y line*=right, axis x line*=none, xtick=\empty,
  ymin=-1.10, ymax=5.50,
  tick label style={font=\tiny},
  yticklabel style={font=\tiny}, scaled y ticks=false,
  legend style={font=\scriptsize, at={(0.5,-0.44)}, anchor=north,
                legend columns=1, draw=none, fill=none},
]
\addplot[OilBlue, mark=*, mark size=2.2pt, line width=1.3pt] table[x=i, y=ocf]{data/clean/profiles/fin_002548.dat};
\addplot[only marks, OilBlue, mark=*, mark size=3.2pt] table[x=x, y=ocf]{data/clean/profiles/finh_002548.dat};
\legend{经营活动现金流净额（右轴，亿元，蓝点）}
\end{axis}
\end{tikzpicture}

\captionof{figure}[金新农（002548）近年主要财务指标]{金新农（002548）近年主要财务指标（左上：营业收入与归母净利润；右上：ROE 与资产负债率；下：货币资金、有息负债为左轴柱状，经营活动现金流净额为其右轴的蓝点折线。
2021---2025 年为年报口径，2026H1 为 2026 年半年报/半年末，与其前一年报之间留出间隔、以斜纹或空心点区分；半年报数据未年化）}
\end{minipage}
\end{center}

\pfl{财务分析} 2026 年 6 月末资产负债率 81.5\%，较 2025 年末上升 5.8 个百分点（样本中位数 52.2\%，所处三级行业内按由低到高排第\zhnumber{11}位）。 2026 年 6 月末货币资金 2.36 亿元、短期债务（短期借款与一年内到期非流动负债）22.72 亿元、长期债务 6.20 亿元，有息负债合计 28.92 亿元，现金短债比 0.10 倍——缺口明显。 净有息负债 26.57 亿元，相当于其 2026 年上半年营业收入的 115\%。 流动比率 0.32、速动比率 0.18，短期偿付压力较大。 2026 年上半年经营活动现金流净额 0.50 亿元，净利润为负而经营现金流为正，差额主要来自折旧摊销与营运资本变动，说明亏损尚未完全转化为现金流出。 2026 年 6 月末在建工程 0.76 亿元，较上年末增加 0.10 亿元，仍有在建投入。 综合看，账面货币资金对有息负债与短期债务的覆盖不足，滚动偿付对新增融资的依赖度较高；资产负债率高于样本中位数 29.3 个百分点；有息负债中短期债务占比更高，期限结构仍以滚动续作为主。 公司披露的股东与质押风险为：2025 年，半年度末资产权利受限合计账面余额 10.90 亿元；2024 年度报告，长期股权投资 3.37 亿元质押用于融资借款。

\pfl{重大战略转型} 近三年公司的战略动作集中在：其他动作（2 项）：2025 年，推出 2025 年股票期权与限制性股票激励计划；2020 年，公司拟向 154 名激励对象授予 1，354.9932 万股限制性股票；并购与重组（1 项）：2016 年参股东进农牧 4\% 股权、收购福建一春 60\% 股权、增资赣州东进 20\% 股权。 公告口径上，近三年（2023 年 9 月---2026 年 9 月）公司披露重大重组与并购类公告 0 条、对外投资与转型类公告 4 条、回购与股权激励类公告 25 条，合计公告 432 条。 第一大行业“饲料加工”的收入占比由 2021 年年报的 47.6\% 变为 69.7\%（22.1 个百分点），收入结构出现再平衡。 综合判断，报告期内公司有 4 条与战略相关的公告落地（重大重组与并购 0 条、对外投资与转型 4 条），资本运作与主业调整之间存在直接关联。

\begin{center}\small\setlength{\tabcolsep}{4pt}
\captionof{table}[金新农（002548）历史融资与资本运作]{金新农（002548）历史融资与资本运作（据公司公告与公开报道整理；金额为披露值，含首次公开发行、再融资、债券、银行借款与重整投资等）}
\begin{adjustbox}{max width=\textwidth}
\begin{tabular}{@{}p{1.45cm}p{2.55cm}p{4.05cm}p{4.95cm}@{}}
\toprule
时间 & 方式 & 规模 & 用途与说明 \\
\midrule
2011-02 & IPO（深交所） & 发行 2，\allowbreak{}400 万股，\allowbreak{}发行价 24.00 元/\allowbreak{}股，\allowbreak{}募集资金总额 5.76 亿元 & 饲料产能建设 \\
2015 & 发行股份及支付现金购买资产并募集配套资金 & 以 5.25 亿元价格收购盈华讯方 80\% 股权；募集配套资金总额 3.67 亿元 & 支付现金对价（1.84 亿元）、\allowbreak{}补充流动资金、\allowbreak{}支付本次交易税费等；发行对象为金新农 2015 年第一期员工持股计划和陈俊海等 6 人；；2018 年 10 月经股东大会审议通过以 1.20 亿元有条件分阶段支付现金方式继续收购盈华讯方剩余股权 \\
2018-03 & 公开发行可转换公司债券 & 发行 650 万张，\allowbreak{}每张面值 100.00 元，\allowbreak{}发行总额 6.50 亿元（6.50 亿元） & 补充流动资金及项目建设；债券简称“金农转债”，\allowbreak{}代码 128036 \\
2020-12-30 & 非公开发行（定增） & 向广州湾区金农投资合伙企业（有限合伙）非公开发行 1.28 亿股，\allowbreak{}发行价 5.08 元/\allowbreak{}股 & 生猪养殖项目及补充流动资金；新增股份 2020 年 12 月 30 日在深交所上市；本次认购后收购人持有公司权益超过 30\% \\
2022 & 非公开发行（定增） & 向特定对象非公开发行 1.17 亿股，\allowbreak{}发行价 5.98 元/\allowbreak{}股，\allowbreak{}募集资金总额 7.00 亿元 & 募投项目建设及补充流动资金 \\
2023 & 子公司增资扩股引入投资者 & 广州乡村振兴投资合伙企业（有限合伙）对子公司广州金农增资 8，\allowbreak{}000.00 万元 & 发展生猪养殖业务；投资期为 3+2 年，\allowbreak{}投资期内每年需按投资金额的 6\% 支付投资收益；投资期 3 年结束后投资方有权要求公司回购股权 \\
\bottomrule
\end{tabular}
\end{adjustbox}
\end{center}

\pfl{历史融投资情况} 从现金流量表看，2025 年公司取得借款收到的现金 19.11 亿元、偿还债务支付的现金 18.77 亿元、吸收权益性投资 0.49 亿元，筹资活动现金流净额 -2.41 亿元（2023 年为取得借款 18.74 亿元、偿还债务 20.03 亿元）。 2026 年上半年筹资活动现金流净额为 -1.73 亿元（取得借款 6.61 亿元、偿还债务 7.64 亿元，未年化），已转为净流出，处于净还债状态。 2025 年分配股利、利润或偿付利息支付的现金合计 1.17 亿元。 公开披露的融资记录共 6 笔，工具类型以 IPO、发行股份及支付现金购、可转债、定向增发为主；金额与用途见下表。 其中规模最大的两笔为：2018 年 3 月，可转债，发行 650 万张，每张面值 100.00 元，发行总额 6.50 亿元（6.50 亿元）；2015 年，发行股份及支付现金购，以 5.25 亿元价格收购盈华讯方 80\% 股权，募集配套资金总额 3.67 亿元。 股本方面，近三年披露股本变动 12 次（可转债转股 9 次、限售股份上市 2 次、激励股份解禁 2 次、股份回购 1 次）；总股本由 2021 年末的 6.91 亿股变为 8.05 亿股（+16.56\%）。 分红方面，近三年未实施现金分红，在全部样本中属于分红能力偏弱的一端。 近三年“再融资”类公告 9 条，涉及定向增发、可转债。 综合看，公司融资结构上以债务滚动为主、股权融资为补充；2025 年筹资活动现金流净额 -2.41 亿元，融资节奏仍以维持存量债务周转为主，与前述经营与投资现金流的方向基本一致。

\pfl{未来潜在融资需求分析} 2026 年 6 月末货币资金 2.36 亿元，而短期债务 22.72 亿元（现金短债比 0.10 倍），2026 年上半年经营现金流净额 0.50 亿元。按“短期债务减去账面货币资金”的滚动偿付口径，静态缺口约 20.37 亿元；若再计入上半年亏损的年化额（7.77 亿元），维持一年正常经营所需的外部资金约 20.37---28.14 亿元。 按第 \ref{sec:financing-need} 节的三档划分，公司属于第一档“补血型”：2026 年上半年归母净利润 -3.88 亿元、6 月末资产负债率 81.5\%。负债率已超过 80\%，净资产被亏损侵蚀的程度较深，属于全部样本中资产负债表最紧张的一组。 可行的融资路径按现实性排序为：债务重组或展期（与银行和债券持有人协商）、定向增发（需股价配合，且控股股东需放弃优先认购或同步增资）、出售非核心资产与参股股权。 资金需求还要叠加在建工程：期末在建工程 0.76 亿元、2026 年上半年购建固定资产等支出 1.06 亿元（未年化），这部分支出在周期底部通常会被主动放缓。 综合看，公司的外部资金需求以债务接续为主，滚动偿付口径下的静态缺口约 20.37 亿元，能否落地取决于授信续作与发行窗口的配合。




\subsubsection{路德科技（688156）}
\label{co:688156}

\pfl{公司基本情况} 公司前身可追溯至 2020 年，科创板上市的生物发酵饲料企业，前身为以无机固废处理为主的路德环境，2023 年起白酒糟资源化利用成为核心主业并于 2025 年 11 月更名为“路德科技”；连续两年亏损后 2026 年上半年扭亏，正推进简易程序定增扩充产能。 公司于 2020-09-22 在上交所科创板首发上市，注册资本 1.01 亿元，总股本 1.01 亿股。 截至 2026 年 9 月 24 日收盘价 22.87 元，流通市值 23.0 亿元，近一年+14.5\%，流通市值在三级行业“畜禽饲料”的 11 家公司中按由大到小排第\zhnumber{9}位，近一年涨跌幅高于全样本中位数（-10.1\%）24.6 个百分点。 按 2025 年年报披露的行业构成，收入占比前两位为生态保护和环境治理业 95.7\%；其他(补充) 4.3\%。

\pfl{历史股价} 自 2021 年 1 月至 2026 年 9 月，股价由 17.19 元变为 22.87 元，累计 +33.0\%，区间最高 37.60 元（2023 年 2 月）、最低 10.78 元（2024 年 8 月）；当前价相当于区间最高价的 60.8\%，为区间最低价的 2.12 倍。 月度收益的年化波动率 48.0\%，高于全样本中位数 37.4\%，在三级行业“畜禽饲料”的 11 家公司中波动率由高到低排第\zhnumber{2}位。

\begin{center}
\begin{minipage}{\textwidth}\centering
\begin{tikzpicture}
\begin{axis}[
  width=12.6cm, height=3.9cm, scale only axis,
  xmin=2020.922, xmax=2026.888, ymin=10.03, ymax=40.23,
  xtick={2021,2022,2023,2024,2025,2026},
  yearticks,
  yticklabel style={font=\scriptsize},
  ylabel={元/股}, ylabel style={font=\scriptsize},
  grid=major, grid style={LightGray, line width=0.3pt},
  axis line style={InkGray, line width=0.5pt},
  every axis plot/.append style={line width=0.9pt},
]
\addplot[AgriGreen] table[x=x,y=y]{data/clean/profiles/px_688156.dat};
\addplot[SoftRed, dashed, line width=0.7pt] table[x=x,y=y]{data/clean/profiles/pxzz_688156.dat};
\addplot[only marks, mark=*, mark size=1.1pt, SoftRed] table[x=x,y=y]{data/clean/profiles/pxzz_688156.dat};
\end{axis}
\end{tikzpicture}

\captionof{figure}[路德科技（688156）股价与周期骨架]{路德科技（688156）月度收盘价（元/股）与 ZigZag 周期骨架（阈值 25\%，圆点为周期转折点）}
\end{minipage}
\end{center}

\pfl{过去周期分析} 以 25\% 的 ZigZag 阈值划分，样本区间内股价形成 4 个主升段与 4 个主跌段。 最大主升段出现在 2022 年 4 月 至 2023 年 2 月，10 个月+145.8\%；最大主跌段为 2023 年 2 月 至 2024 年 8 月，18 个月-71.3\%。 最近一次转折出现在 2026 年 7 月（下跌段自 2026 年 6 月起，1 个月-39.2\%），当前处于该段的延续之中。 公司股价较申万农林牧渔指数提前约 21 个月见底，是板块中较早定价周期反转的公司之一。 饲料夹在原料成本与养殖需求之间（第 \ref{sec:feed} 节），量的部分取决于存栏、利的部分取决于原料与成品的价差，公司 2026 年上半年实现盈利、半年 ROE 1.2\%（未年化），好于全样本中位数（0.75\%），下行周期对其报表的冲击小于同业。 综合区间形态看，该股单段涨跌幅度偏大、以趋势段为主（最大主升 +145.8\%、最大主跌 -71.3\%），年化波动率 48.0\% 与全样本中位数 37.4\% 相比波动明显大于样本中位数。

\pfl{盈利情况} 2026 年上半年营业收入 2.10 亿元（同比 +42.9\%），归母净利润 0.09 亿元，实现扭亏。 以年报为趋势对照，2023---2025 年营业收入由 3.51 亿元变为 3.80 亿元（两年年均复合增速 +4.1\%），归母净利润由 0.27 亿元变为 -0.70 亿元。 按半年报口径，2026 年上半年的 ROE 为 1.2\%（未年化）、销售净利率 4.7\%、毛利率 30.7\%，ROE 在“畜禽饲料”11 家公司中按数值由高到低排第\zhnumber{2}位。 同期全样本半年 ROE 中位数 0.75\%，公司与之相差 0.41 个百分点（略好于中位数公司）。 综合看，公司在报告期维持盈利（高于全样本半年 ROE 中位数 0.75\%），利润的可持续性取决于畜禽饲料景气的延续与成本管控的执行。

\begin{center}\small\setlength{\tabcolsep}{3.4pt}
\captionof{table}[路德科技（688156）关键财务指标]{路德科技（688156）关键财务指标（合并报表口径；两个半年报期均未年化；现金短债比 $=$ 货币资金 $\div$ 短期债务）}
\begin{adjustbox}{max width=\textwidth}
\begin{tabular}{@{}lrrrrr@{}}
\toprule
指标 & 2023 年 & 2024 年 & 2025 年 & 2025 年半年报 & 2026 年半年报 \\
\midrule
营业收入（亿元） & 3.51 & 2.78 & 3.80 & 1.47 & 2.10 \\
归母净利润（亿元） & 0.27 & -0.57 & -0.70 & -0.13 & 0.09 \\
毛利率（\%） & 36.2 & 21.9 & 19.9 & 18.9 & 30.7 \\
ROE（\%） & 3.2 & -6.5 & -8.8 & -1.6 & 1.2 \\
资产负债率（\%） & 40.5 & 49.4 & 58.3 & 47.9 & 55.9 \\
经营活动现金流净额（亿元） & -0.19 & 0.53 & 0.33 & 0.12 & 0.33 \\
货币资金（亿元） & 1.99 & 2.26 & 3.12 & 1.55 & 2.43 \\
有息负债合计（亿元） & 2.90 & 5.02 & 5.85 & 5.22 & 5.50 \\
现金短债比（倍） & 1.27 & 0.86 & 0.83 & 0.51 & 0.72 \\
\bottomrule
\end{tabular}
\end{adjustbox}
\end{center}

\begin{center}
\begin{minipage}{\textwidth}\centering
\begin{tikzpicture}
\begin{groupplot}[
  group style={group size=2 by 1, horizontal sep=0.60cm},
  width=6.85cm, height=4.60cm,
  tick label style={font=\tiny},
  title style={font=\scriptsize\heiti},
  yticklabel style={font=\tiny},
  ylabel style={font=\tiny},
  grid=major, grid style={LightGray, line width=0.3pt},
  axis line style={InkGray, line width=0.5pt},
  enlarge x limits={abs=0.8},
  legend style={font=\scriptsize, at={(0.5,-0.20)}, anchor=north,
                legend columns=2, draw=none, fill=none, column sep=6pt},
]
\nextgroupplot[
  ybar, bar width=4pt,
  ymin=-1.15, ymax=4.36,
  xtick={0,1,2,3,4,5.7},
  xticklabels={2021,2022,2023,2024,2025,2026H1},
  ylabel={亿元},
]
\addplot[fill=AgriGreen!75, draw=AgriGreen, bar shift=-2.6pt] table[x=i, y=rev]{data/clean/profiles/fin_688156.dat};
\addplot[fill=SoftRed!60, draw=SoftRed, bar shift=2.6pt] table[x=i, y=ni]{data/clean/profiles/fin_688156.dat};
\addplot[fill=AgriGreen!30, draw=AgriGreen, pattern=north east lines, bar shift=-2.6pt] table[x=x, y=rev]{data/clean/profiles/finh_688156.dat};
\addplot[fill=SoftRed!20, draw=SoftRed, pattern=north east lines, bar shift=2.6pt] table[x=x, y=ni]{data/clean/profiles/finh_688156.dat};
\legend{营业收入（年/半年）, 归母净利润（年/半年）}
\nextgroupplot[
  ymin=-14.1, ymax=65.0,
  xtick={0,1,2,3,4,5.7},
  xticklabels={2021,2022,2023,2024,2025,2026H1},
  ylabel={\%},
]
\addplot[AgriGold, mark=*, mark size=1.3pt] table[x=i, y=roe]{data/clean/profiles/fin_688156.dat};
\addplot[OilBlue, mark=square*, mark size=1.3pt] table[x=i, y=debt]{data/clean/profiles/fin_688156.dat};
\addplot[only marks, AgriGold, mark=*, mark size=2.0pt] table[x=x, y=roe]{data/clean/profiles/finh_688156.dat};
\addplot[only marks, OilBlue, mark=square*, mark size=2.0pt] table[x=x, y=debt]{data/clean/profiles/finh_688156.dat};
\legend{ROE, 资产负债率}
\end{groupplot}
\end{tikzpicture}
\begin{tikzpicture}
\begin{axis}[
  name=finbot,
  width=13.75cm, height=3.10cm, scale only axis,
  ymin=-0.59, ymax=6.55,
  xtick={0,1,2,3,4,5.7},
  xticklabels={2021,2022,2023,2024,2025,2026H1},
  tick label style={font=\tiny},
  yticklabel style={font=\tiny},
  ylabel={亿元}, ylabel style={font=\tiny},
  ybar, bar width=4pt,
  grid=major, grid style={LightGray, line width=0.3pt},
  axis line style={InkGray, line width=0.5pt},
  enlarge x limits={abs=0.8},
  legend style={font=\scriptsize, at={(0.5,-0.26)}, anchor=north,
                legend columns=2, draw=none, fill=none, column sep=10pt},
]
\addplot[fill=AgriGreen!55, draw=AgriGreen, bar shift=-3.0pt] table[x=i, y=cash]{data/clean/profiles/fin_688156.dat};
\addplot[fill=SoftRed!45, draw=SoftRed, bar shift=3.0pt] table[x=i, y=ibd]{data/clean/profiles/fin_688156.dat};
\addplot[fill=AgriGreen!25, draw=AgriGreen, pattern=north east lines, bar shift=-3.0pt] table[x=x, y=cash]{data/clean/profiles/finh_688156.dat};
\addplot[fill=SoftRed!20, draw=SoftRed, pattern=north east lines, bar shift=3.0pt] table[x=x, y=ibd]{data/clean/profiles/finh_688156.dat};
\legend{货币资金（左轴）, 有息负债（左轴）}
\end{axis}
\begin{axis}[
  at={(finbot.south east)}, anchor=south east,
  width=13.75cm, height=3.10cm, scale only axis,
  axis y line*=right, axis x line*=none, xtick=\empty,
  ymin=-0.38, ymax=0.75,
  tick label style={font=\tiny},
  yticklabel style={font=\tiny}, scaled y ticks=false,
  legend style={font=\scriptsize, at={(0.5,-0.44)}, anchor=north,
                legend columns=1, draw=none, fill=none},
]
\addplot[OilBlue, mark=*, mark size=2.2pt, line width=1.3pt] table[x=i, y=ocf]{data/clean/profiles/fin_688156.dat};
\addplot[only marks, OilBlue, mark=*, mark size=3.2pt] table[x=x, y=ocf]{data/clean/profiles/finh_688156.dat};
\legend{经营活动现金流净额（右轴，亿元，蓝点）}
\end{axis}
\end{tikzpicture}

\captionof{figure}[路德科技（688156）近年主要财务指标]{路德科技（688156）近年主要财务指标（左上：营业收入与归母净利润；右上：ROE 与资产负债率；下：货币资金、有息负债为左轴柱状，经营活动现金流净额为其右轴的蓝点折线。
2021---2025 年为年报口径，2026H1 为 2026 年半年报/半年末，与其前一年报之间留出间隔、以斜纹或空心点区分；半年报数据未年化）}
\end{minipage}
\end{center}

\pfl{财务分析} 2026 年 6 月末资产负债率 55.9\%，较 2025 年末下降 2.4 个百分点（样本中位数 52.2\%，所处三级行业内按由低到高排第\zhnumber{4}位）。 2026 年 6 月末货币资金 2.43 亿元、短期债务（短期借款与一年内到期非流动负债）3.38 亿元、长期债务 2.12 亿元，有息负债合计 5.50 亿元，现金短债比 0.72 倍——覆盖偏紧。 净有息负债 3.06 亿元，相当于其 2026 年上半年营业收入的 146\%。 流动比率 0.82、速动比率 0.69，短期偿付压力较大。 2026 年上半年经营活动现金流净额 0.33 亿元，经营现金流与利润同向为正，内生资金可以支撑运营。 2026 年 6 月末在建工程 1.09 亿元，较上年末增加 0.62 亿元，仍有在建投入。 综合看，现金短债比 0.72 倍，短期资金安排偏紧；资产负债率高于样本中位数 3.7 个百分点；有息负债中短期债务占比更高，期限结构仍以滚动续作为主。 公司披露的股东与质押风险为：2026 年 4 月，持股 1.4894\% 的股东武汉德天众享企业管理合伙企业（有限合伙）计划减持不超 1.4894\% 公司股份。

\pfl{重大战略转型} 公司的战略动作可归为以下几类：融资与资本运作（2 项）：2026 年 9 月，公司以简易程序定增募集资金主要用于洋河酒糟资源化利用项目建设；2024 年 4 月，提请 2025 年年度股东会授权董事会办理以简易程序向特定对象发行融资总额不超过人民币三亿元且不超过最近一年末净资产 20\% 的股票；主业转型（1 项）：2025 年 11 月，证券简称由“路德环境”变更为“路德科技”。 公告口径上，近三年（2023 年 9 月---2026 年 9 月）公司披露重大重组与并购类公告 0 条、对外投资与转型类公告 3 条、回购与股权激励类公告 20 条，合计公告 411 条。 第一大行业仍是“生态保护和环境治理业”，收入占比 95.7\%，与 2021 年年报相比没有方向性变化。 综合判断，报告期内公司有 3 条与战略相关的公告落地（重大重组与并购 0 条、对外投资与转型 3 条），资本运作与主业调整之间存在直接关联。

\begin{center}\small\setlength{\tabcolsep}{4pt}
\captionof{table}[路德科技（688156）历史融资与资本运作]{路德科技（688156）历史融资与资本运作（据公司公告与公开报道整理；金额为披露值，含首次公开发行、再融资、债券、银行借款与重整投资等）}
\begin{adjustbox}{max width=\textwidth}
\begin{tabular}{@{}p{1.45cm}p{2.55cm}p{4.05cm}p{4.95cm}@{}}
\toprule
时间 & 方式 & 规模 & 用途与说明 \\
\midrule
2020-09 & IPO（上交所科创板） & 首次公开发行 2，\allowbreak{}296 万股，\allowbreak{}发行价 15.91 元/\allowbreak{}股，\allowbreak{}募集资金总额 3.65 亿元 & 研发中心及固废处理相关项目建设；；2020 年 9 月 10 日发行 \\
2023-05 & 向特定对象发行股票（定增） & 发行 834.04 万股，\allowbreak{}发行价 13.57 元/\allowbreak{}股，\allowbreak{}募集资金总额 1.13 亿元 & 补充流动资金及项目建设；发行对象为公司实际控制人季光明 \\
2025-12-22 & 子公司引入战略投资者增资 & 国投聚力旗下聚力杭实基金对永乐路德、\allowbreak{}亳州路德、\allowbreak{}宿迁路德三家子公司战略性增资，\allowbreak{}总投资额 1.17 亿元（亳州路德 3 & 加码白酒糟生物发酵饲料核心主业；公司与国投聚力投资管理有限公司签署战略合作协议；聚力杭实基金股东包括国投聚力、\allowbreak{}杭州汇实投资及萧山经开区产业基金 \\
2026-07-27 & 以简易程序向特定对象发行股票（定增，推进中） & 拟募集资金总额不超过 1.50 亿元（含本数） & 洋河酒糟资源化利用项目（拟使用 1.10 亿元）及补充流动资金（4 \\
\bottomrule
\end{tabular}
\end{adjustbox}
\end{center}

\pfl{历史融投资情况} 从现金流量表看，2025 年公司取得借款收到的现金 3.29 亿元、偿还债务支付的现金 2.77 亿元，筹资活动现金流净额 2.15 亿元（2023 年为取得借款 2.14 亿元、偿还债务 1.00 亿元）。 2026 年上半年筹资活动现金流净额为 -0.36 亿元（取得借款 1.36 亿元、偿还债务 1.39 亿元，未年化），已转为净流出，处于净还债状态。 2025 年分配股利、利润或偿付利息支付的现金合计 0.20 亿元。 公开披露的融资记录共 4 笔，工具类型以 IPO、定向增发、战略投资为主；金额与用途见下表。 其中规模最大的两笔为：2026 年 7 月，定向增发，拟募集资金总额不超过 1.50 亿元（含本数）；2025 年 12 月，战略投资，国投聚力旗下聚力杭实基金对永乐路德、亳州路德、宿迁路德三家子公司战略性增资，总投资额 1.17 亿元（亳州路德 3。 股本方面，近三年披露股本变动 5 次（增发新股上市 1 次、限售股份上市 1 次）；总股本由 2021 年末的 0.92 亿股变为 1.01 亿股（+9.66\%）。 分红方面，近三年现金分红 2 次、合计每股派现 0.15 元，最近一次实施在 2023 年 12 月。 近三年“再融资”类公告 45 条，涉及定向增发、可转债、募集资金使用。 综合看，公司融资结构上股权与债务融资并举；2025 年筹资活动现金流净额 2.15 亿元，年度内仍为净融入，与前述经营与投资现金流的方向基本一致。

\pfl{未来潜在融资需求分析} 2026 年 6 月末货币资金 2.43 亿元，而短期债务 3.38 亿元（现金短债比 0.72 倍），2026 年上半年经营现金流净额 0.33 亿元。按“短期债务减去账面货币资金”的滚动偿付口径，静态缺口约 0.94 亿元；上半年盈利或接近盈亏平衡，该缺口不随经营亏损进一步扩大。 公司 2026 年 6 月末资产负债率 55.9\%、2026 年上半年 ROE 1.2\%（未年化），现金短债比 0.72 倍、短债覆盖不足，融资需求首先用于置换与滚动存量债务、补充营运资金，属于“补血型”而非扩张型。 公开披露的后续资金安排线索包括：股权融资线索：2026 年 7 月，公司以简易程序向特定对象发行股票募资不超 1.52 亿元，用途为洋河酒糟资源化利用项目与补充流动资金；资本开支安排：2020 年，首次公开发行股票募集资金使用情况。 资金需求还要叠加在建工程：期末在建工程 1.09 亿元、2026 年上半年购建固定资产等支出 0.70 亿元（未年化），这部分支出在周期底部通常会被主动放缓。 综合看，公司的外部资金需求以补充营运资金为主、项目投入为辅，滚动偿付口径下的静态缺口约 0.94 亿元，融资节奏将受制于银行续贷意愿与股权融资窗口。




\subsubsection{百洋股份（002696）}
\label{co:002696}

\pfl{公司基本情况} 2000 年设立至今，百洋股份是广西南宁的水产企业，以罗非鱼等水产品加工出口起家，2012 年上市；2017 年以 9.74 亿元收购火星时代跨界数字艺术教育、2019 年原价回购剥离，2020 年 5 月控制权转让给青岛国资，现回归水产与饲料主业，2026 年上半年重新陷入亏损。 公司于 2012-09-05 在深交所主板首发上市，注册资本 3.46 亿元，总股本 3.46 亿股。 截至 2026 年 9 月 24 日收盘价 6.65 元，流通市值 22.2 亿元，近一年+8.0\%，流通市值在三级行业“水产饲料”的 4 家公司中按由大到小排第\zhnumber{4}位，近一年涨跌幅高于全样本中位数（-10.1\%）18.0 个百分点。 按 2026 年半年报披露的行业构成，收入占比前三位为饲料行业 51.8\%；食品加工行业 38.2\%；其他 9.9\%。 股权结构方面，实际控制人为青岛市国资委。

\pfl{历史股价} 以 2021 年 1 月为起点、2026 年 9 月为终点，股价由 5.44 元变为 6.65 元，累计 +22.2\%，区间最高 7.86 元（2026 年 2 月）、最低 3.88 元（2024 年 8 月）；当前价相当于区间最高价的 84.6\%，为区间最低价的 1.71 倍。 月度收益的年化波动率 32.2\%，低于全样本中位数 37.4\%，在三级行业“水产饲料”的 4 家公司中波动率由高到低排第\zhnumber{2}位。

\begin{center}
\begin{minipage}{\textwidth}\centering
\begin{tikzpicture}
\begin{axis}[
  width=12.6cm, height=3.9cm, scale only axis,
  xmin=2020.922, xmax=2026.888, ymin=3.61, ymax=8.41,
  xtick={2021,2022,2023,2024,2025,2026},
  yearticks,
  yticklabel style={font=\scriptsize},
  ylabel={元/股}, ylabel style={font=\scriptsize},
  grid=major, grid style={LightGray, line width=0.3pt},
  axis line style={InkGray, line width=0.5pt},
  every axis plot/.append style={line width=0.9pt},
]
\addplot[AgriGreen] table[x=x,y=y]{data/clean/profiles/px_002696.dat};
\addplot[SoftRed, dashed, line width=0.7pt] table[x=x,y=y]{data/clean/profiles/pxzz_002696.dat};
\addplot[only marks, mark=*, mark size=1.1pt, SoftRed] table[x=x,y=y]{data/clean/profiles/pxzz_002696.dat};
\end{axis}
\end{tikzpicture}

\captionof{figure}[百洋股份（002696）股价与周期骨架]{百洋股份（002696）月度收盘价（元/股）与 ZigZag 周期骨架（阈值 25\%，圆点为周期转折点）}
\end{minipage}
\end{center}

\pfl{过去周期分析} 以 25\% 的 ZigZag 阈值划分，样本区间内股价形成 3 个主升段与 3 个主跌段。 最大主跌段为 2022 年 12 月 至 2024 年 8 月，20 个月-48.5\%；最大主升段出现在 2024 年 8 月 至 2026 年 2 月，18 个月+102.6\%。 最近一次转折出现在 2026 年 9 月（上涨段自 2026 年 6 月起，3 个月+33.8\%），当前处于该段之后的震荡之中。 公司股价较申万农林牧渔指数提前约 21 个月见底，是板块中较早定价周期反转的公司之一。 水产饲料的需求取决于投苗量与存塘鱼价，第 \ref{sec:other-livestock} 章显示 2026 年淡水鱼价格与投苗量已率先回升，其需求驱动与生猪存栏的联系弱于畜禽饲料，公司 2026 年 6 月末资产负债率 66.4\%，高于全样本中位数 52.2\%，其股价因此更多反映资金面与信用风险，而不是价格弹性本身。 综合区间形态看，该股单段涨跌幅度偏大、以趋势段为主（最大主升 +102.6\%、最大主跌 -48.5\%），年化波动率 32.2\% 与全样本中位数 37.4\% 相比波动小于样本中位数。

\pfl{盈利情况} 2026 年上半年营业收入 20.51 亿元（同比 +21.1\%），归母净利润 -0.83 亿元，由上年同期的盈利转为亏损。 以年报为趋势对照，2023---2025 年营业收入由 27.20 亿元变为 40.55 亿元（两年年均复合增速 +22.1\%），归母净利润由 -0.28 亿元变为 0.41 亿元。 按半年报口径，2026 年上半年的 ROE 为 -6.2\%（未年化）、销售净利率 -4.2\%、毛利率 6.2\%，ROE 在“水产饲料”4 家公司中按数值由高到低排第\zhnumber{4}位。 按同一口径，三级行业“水产饲料”4 家公司的半年 ROE 中位数为 1.22\%，该公司低于其中位数 7.42 个百分点。 综合看，公司仍处亏损区间、由盈转亏，半年 ROE 在三级行业内按由高到低排第\zhnumber{4}位，单位盈利能力的修复依赖行业景气与自身成本端的改善，报表弹性主要体现为价格与销量的方向性变化。

\begin{center}\small\setlength{\tabcolsep}{3.4pt}
\captionof{table}[百洋股份（002696）关键财务指标]{百洋股份（002696）关键财务指标（合并报表口径；两个半年报期均未年化；现金短债比 $=$ 货币资金 $\div$ 短期债务）}
\begin{adjustbox}{max width=\textwidth}
\begin{tabular}{@{}lrrrrr@{}}
\toprule
指标 & 2023 年 & 2024 年 & 2025 年 & 2025 年半年报 & 2026 年半年报 \\
\midrule
营业收入（亿元） & 27.20 & 30.78 & 40.55 & 16.94 & 20.51 \\
归母净利润（亿元） & -0.28 & -0.15 & 0.41 & 0.12 & -0.83 \\
毛利率（\%） & 8.3 & 9.0 & 8.8 & 8.8 & 6.2 \\
ROE（\%） & -1.9 & -1.1 & 2.9 & 0.9 & -6.2 \\
资产负债率（\%） & 49.8 & 53.5 & 60.9 & 60.0 & 66.4 \\
经营活动现金流净额（亿元） & 2.40 & -0.60 & -3.15 & -1.99 & -3.30 \\
货币资金（亿元） & 6.20 & 4.55 & 4.09 & 6.65 & 5.32 \\
有息负债合计（亿元） & 8.10 & 6.35 & 12.19 & 11.48 & 16.70 \\
现金短债比（倍） & 1.08 & 0.84 & 0.46 & 0.78 & 0.35 \\
\bottomrule
\end{tabular}
\end{adjustbox}
\end{center}

\begin{center}
\begin{minipage}{\textwidth}\centering
\begin{tikzpicture}
\begin{groupplot}[
  group style={group size=2 by 1, horizontal sep=0.60cm},
  width=6.85cm, height=4.60cm,
  tick label style={font=\tiny},
  title style={font=\scriptsize\heiti},
  yticklabel style={font=\tiny},
  ylabel style={font=\tiny},
  grid=major, grid style={LightGray, line width=0.3pt},
  axis line style={InkGray, line width=0.5pt},
  enlarge x limits={abs=0.8},
  legend style={font=\scriptsize, at={(0.5,-0.20)}, anchor=north,
                legend columns=2, draw=none, fill=none, column sep=6pt},
]
\nextgroupplot[
  ybar, bar width=4pt,
  ymin=-4.96, ymax=45.52,
  xtick={0,1,2,3,4,5.7},
  xticklabels={2021,2022,2023,2024,2025,2026H1},
  ylabel={亿元},
]
\addplot[fill=AgriGreen!75, draw=AgriGreen, bar shift=-2.6pt] table[x=i, y=rev]{data/clean/profiles/fin_002696.dat};
\addplot[fill=SoftRed!60, draw=SoftRed, bar shift=2.6pt] table[x=i, y=ni]{data/clean/profiles/fin_002696.dat};
\addplot[fill=AgriGreen!30, draw=AgriGreen, pattern=north east lines, bar shift=-2.6pt] table[x=x, y=rev]{data/clean/profiles/finh_002696.dat};
\addplot[fill=SoftRed!20, draw=SoftRed, pattern=north east lines, bar shift=2.6pt] table[x=x, y=ni]{data/clean/profiles/finh_002696.dat};
\legend{营业收入（年/半年）, 归母净利润（年/半年）}
\nextgroupplot[
  ymin=-12.0, ymax=73.7,
  xtick={0,1,2,3,4,5.7},
  xticklabels={2021,2022,2023,2024,2025,2026H1},
  ylabel={\%},
]
\addplot[AgriGold, mark=*, mark size=1.3pt] table[x=i, y=roe]{data/clean/profiles/fin_002696.dat};
\addplot[OilBlue, mark=square*, mark size=1.3pt] table[x=i, y=debt]{data/clean/profiles/fin_002696.dat};
\addplot[only marks, AgriGold, mark=*, mark size=2.0pt] table[x=x, y=roe]{data/clean/profiles/finh_002696.dat};
\addplot[only marks, OilBlue, mark=square*, mark size=2.0pt] table[x=x, y=debt]{data/clean/profiles/finh_002696.dat};
\legend{ROE, 资产负债率}
\end{groupplot}
\end{tikzpicture}
\begin{tikzpicture}
\begin{axis}[
  name=finbot,
  width=13.75cm, height=3.10cm, scale only axis,
  ymin=-1.67, ymax=18.71,
  xtick={0,1,2,3,4,5.7},
  xticklabels={2021,2022,2023,2024,2025,2026H1},
  tick label style={font=\tiny},
  yticklabel style={font=\tiny},
  ylabel={亿元}, ylabel style={font=\tiny},
  ybar, bar width=4pt,
  grid=major, grid style={LightGray, line width=0.3pt},
  axis line style={InkGray, line width=0.5pt},
  enlarge x limits={abs=0.8},
  legend style={font=\scriptsize, at={(0.5,-0.26)}, anchor=north,
                legend columns=2, draw=none, fill=none, column sep=10pt},
]
\addplot[fill=AgriGreen!55, draw=AgriGreen, bar shift=-3.0pt] table[x=i, y=cash]{data/clean/profiles/fin_002696.dat};
\addplot[fill=SoftRed!45, draw=SoftRed, bar shift=3.0pt] table[x=i, y=ibd]{data/clean/profiles/fin_002696.dat};
\addplot[fill=AgriGreen!25, draw=AgriGreen, pattern=north east lines, bar shift=-3.0pt] table[x=x, y=cash]{data/clean/profiles/finh_002696.dat};
\addplot[fill=SoftRed!20, draw=SoftRed, pattern=north east lines, bar shift=3.0pt] table[x=x, y=ibd]{data/clean/profiles/finh_002696.dat};
\legend{货币资金（左轴）, 有息负债（左轴）}
\end{axis}
\begin{axis}[
  at={(finbot.south east)}, anchor=south east,
  width=13.75cm, height=3.10cm, scale only axis,
  axis y line*=right, axis x line*=none, xtick=\empty,
  ymin=-4.78, ymax=4.11,
  tick label style={font=\tiny},
  yticklabel style={font=\tiny}, scaled y ticks=false,
  legend style={font=\scriptsize, at={(0.5,-0.44)}, anchor=north,
                legend columns=1, draw=none, fill=none},
]
\addplot[OilBlue, mark=*, mark size=2.2pt, line width=1.3pt] table[x=i, y=ocf]{data/clean/profiles/fin_002696.dat};
\addplot[only marks, OilBlue, mark=*, mark size=3.2pt] table[x=x, y=ocf]{data/clean/profiles/finh_002696.dat};
\legend{经营活动现金流净额（右轴，亿元，蓝点）}
\end{axis}
\end{tikzpicture}

\captionof{figure}[百洋股份（002696）近年主要财务指标]{百洋股份（002696）近年主要财务指标（左上：营业收入与归母净利润；右上：ROE 与资产负债率；下：货币资金、有息负债为左轴柱状，经营活动现金流净额为其右轴的蓝点折线。
2021---2025 年为年报口径，2026H1 为 2026 年半年报/半年末，与其前一年报之间留出间隔、以斜纹或空心点区分；半年报数据未年化）}
\end{minipage}
\end{center}

\pfl{财务分析} 2026 年 6 月末资产负债率 66.4\%，较 2025 年末上升 5.5 个百分点（样本中位数 52.2\%，所处三级行业内按由低到高排第\zhnumber{3}位）。 2026 年 6 月末货币资金 5.32 亿元、短期债务（短期借款与一年内到期非流动负债）15.11 亿元、长期债务 1.59 亿元，有息负债合计 16.70 亿元，现金短债比 0.35 倍——覆盖偏紧。 净有息负债 11.38 亿元，相当于其 2026 年上半年营业收入的 55\%。 流动比率 1.28、速动比率 0.71，短期流动性无明显缺口。 2026 年上半年经营活动现金流净额 -3.30 亿元，经营现金流与利润同时为负，日常运营对债务与股东投入的依赖度较高。 2026 年 6 月末在建工程 1.18 亿元，较上年末增加 0.64 亿元，仍有在建投入。 综合看，账面货币资金对有息负债与短期债务的覆盖不足，滚动偿付对新增融资的依赖度较高；资产负债率高于样本中位数 14.2 个百分点；有息负债中短期债务占比更高，期限结构仍以滚动续作为主。

\pfl{重大战略转型} 近三年公司的战略动作集中在：并购与重组（2 项）：2019 年，教育文化业务大幅下滑后剥离火星时代；2017 年，当年以 9.74 亿元收购火星时代 100\% 股权。 公告口径上，近三年（2023 年 9 月---2026 年 9 月）公司披露重大重组与并购类公告 0 条、对外投资与转型类公告 1 条、回购与股权激励类公告 39 条，合计公告 337 条。 第一大行业仍是“饲料行业”，收入占比 51.8\%，与 2021 年年报相比没有方向性变化。 综合判断，报告期内公司有 1 条与战略相关的公告落地（重大重组与并购 0 条、对外投资与转型 1 条），资本运作与主业调整之间存在直接关联。

\begin{center}\small\setlength{\tabcolsep}{4pt}
\captionof{table}[百洋股份（002696）历史融资与资本运作]{百洋股份（002696）历史融资与资本运作（据公司公告与公开报道整理；金额为披露值，含首次公开发行、再融资、债券、银行借款与重整投资等）}
\begin{adjustbox}{max width=\textwidth}
\begin{tabular}{@{}p{1.45cm}p{2.55cm}p{4.05cm}p{4.95cm}@{}}
\toprule
时间 & 方式 & 规模 & 用途与说明 \\
\midrule
2012-09 & IPO（深交所） & 公开发行 2，\allowbreak{}200 万股，\allowbreak{}发行价 23.90 元/\allowbreak{}股 & 水产品加工产能建设；发行时公司名称为百洋水产集团股份有限公司 \\
2017-09 & 非公开发行（配套融资） & 向孙忠义等 7 名特定投资者非公开发行 2，\allowbreak{}952.17 万股，\allowbreak{}发行价 18.19 元/\allowbreak{}股 & 支付收购火星时代的现金对价等；新增股份上市日为 2017 年 9 月 27 日 \\
2020-03-24 & 控制权协议转让 & 孙忠义、\allowbreak{}蔡晶将其持有的 1.04 亿股（占公司股份总数 29.90\%）以 9.45 元/\allowbreak{}股协议转让给青岛海洋创新产业投资基金有限公司、\allowbreak{}青岛市海洋新动能产业投资基金（有限合伙） & 公司控制权变更；2020 年 5 月 29 日完成股份交割过户 \\
2024-2025 & 限制性股票激励计划 & 2025 年 7 月 31 日为预留部分授予日，\allowbreak{}以 3.14 元/\allowbreak{}股向 21 名激励对象授予 98.80 万股限制性股票 & 核心团队长期激励；2024 年限制性股票激励计划的预留授予实施 \\
2025 & 股份回购 & 2024 年度支付回购库存股款项 5,\allowbreak{}258.90 万元 & 回购公司股份；2025 年度利润分配以 2025 年 12 月 31 日总股本扣除回购专户库存股后 3.46 亿股为基数 \\
\bottomrule
\end{tabular}
\end{adjustbox}
\end{center}

\pfl{历史融投资情况} 从现金流量表看，2025 年公司取得借款收到的现金 11.82 亿元、偿还债务支付的现金 7.31 亿元、吸收权益性投资 0.22 亿元，筹资活动现金流净额 4.73 亿元（2023 年为取得借款 8.46 亿元、偿还债务 4.82 亿元）。 2026 年上半年筹资活动现金流净额为 5.26 亿元（取得借款 9.31 亿元、偿还债务 4.73 亿元，未年化），仍处于净融入状态。 2025 年分配股利、利润或偿付利息支付的现金合计 0.41 亿元。 公开披露的融资记录共 5 笔，工具类型以 IPO、定向增发为主；金额与用途见下表。 其中规模最大的两笔为：2020 年 3 月，控制权协议转让，孙忠义、蔡晶将其持有的 1.04 亿股（占公司股份总数 29.90\%）以 9.45 元/股协议转让给青岛海洋创新产业投资基金有限公司、青岛市海洋新动能产业投资基金（有限合伙）。 股本方面，近三年披露股本变动 5 次（股份回购 1 次、其他 1 次）；总股本由 2021 年末的 3.49 亿股变为 3.46 亿股（-0.87\%）。 分红方面，近三年现金分红 1 次、合计每股派现 0.05 元，最近一次实施在 2025 年 12 月。 综合看，公司融资结构上以债务滚动为主、股权融资为补充；2025 年筹资活动现金流净额 4.73 亿元，年度内仍为净融入，与前述经营与投资现金流的方向基本一致。

\pfl{未来潜在融资需求分析} 2026 年 6 月末货币资金 5.32 亿元，而短期债务 15.11 亿元（现金短债比 0.35 倍），2026 年上半年经营现金流净额 -3.30 亿元。按“短期债务减去账面货币资金”的滚动偿付口径，静态缺口约 9.79 亿元；若再计入上半年亏损的年化额（1.65 亿元），维持一年正常经营所需的外部资金约 9.79---11.44 亿元。 按第 \ref{sec:financing-need} 节的三档划分，公司属于第一档“补血型”：2026 年上半年归母净利润 -0.83 亿元、6 月末资产负债率 66.4\%。 可行的融资路径按现实性排序为：债务重组或展期（与银行和债券持有人协商）、出售非核心资产与参股股权、定向增发（需股价配合，且控股股东需放弃优先认购或同步增资）。 公开披露的后续资金安排线索包括：资本开支安排：2025 年度，购建长期资产支付 1.87 亿元（上年同期 1.10 亿元）。 资金需求还要叠加在建工程：期末在建工程 1.18 亿元、2026 年上半年购建固定资产等支出 0.93 亿元（未年化），这部分支出在周期底部通常会被主动放缓。 综合看，公司的外部资金需求以债务接续为主，滚动偿付口径下的静态缺口约 9.79 亿元，能否落地取决于授信续作与发行窗口的配合。




\subsubsection{播恩集团（001366）}
\label{co:001366}

\pfl{公司基本情况} 2006 年设立至今，播恩集团是总部位于广州、注册在江西赣州的猪饲料企业，聚焦教槽料、乳猪料等幼小动物营养产品，2023 年 3 月登陆深交所主板；上市后连续两年亏损，2026 年起通过终止募投项目、以募集资金改道收购希杰三家饲料公司切入反刍料赛道。 公司于 2023-03-07 在深交所主板首发上市，注册资本 1.61 亿元，总股本 1.61 亿股。 截至 2026 年 9 月 24 日收盘价 13.05 元，流通市值 21.0 亿元，近一年+4.3\%，流通市值在三级行业“畜禽饲料”的 11 家公司中按由大到小排第\zhnumber{10}位，近一年涨跌幅高于全样本中位数（-10.1\%）14.4 个百分点。 按 2026 年半年报披露的行业构成，收入几乎全部来自饲料行业 100.0\%。

\pfl{历史股价} 本报告样本区间（2023 年 3 月 至 2026 年 9 月）内，股价由 18.21 元变为 13.05 元，累计 -28.3\%，区间最高 18.21 元（2023 年 3 月）、最低 9.92 元（2024 年 8 月）；当前价相当于区间最高价的 71.7\%，为区间最低价的 1.32 倍。 月度收益的年化波动率 36.7\%，低于全样本中位数 37.4\%，在三级行业“畜禽饲料”的 11 家公司中波动率由高到低排第\zhnumber{6}位。

\begin{center}
\begin{minipage}{\textwidth}\centering
\begin{tikzpicture}
\begin{axis}[
  width=12.6cm, height=3.9cm, scale only axis,
  xmin=2023.088, xmax=2026.888, ymin=9.23, ymax=19.48,
  xtick={2023,2024,2025,2026},
  yearticks,
  yticklabel style={font=\scriptsize},
  ylabel={元/股}, ylabel style={font=\scriptsize},
  grid=major, grid style={LightGray, line width=0.3pt},
  axis line style={InkGray, line width=0.5pt},
  every axis plot/.append style={line width=0.9pt},
]
\addplot[AgriGreen] table[x=x,y=y]{data/clean/profiles/px_001366.dat};
\addplot[SoftRed, dashed, line width=0.7pt] table[x=x,y=y]{data/clean/profiles/pxzz_001366.dat};
\addplot[only marks, mark=*, mark size=1.1pt, SoftRed] table[x=x,y=y]{data/clean/profiles/pxzz_001366.dat};
\end{axis}
\end{tikzpicture}

\captionof{figure}[播恩集团（001366）股价与周期骨架]{播恩集团（001366）月度收盘价（元/股）与 ZigZag 周期骨架（阈值 25\%，圆点为周期转折点）}
\end{minipage}
\end{center}

\pfl{过去周期分析} 以 25\% 的 ZigZag 阈值划分，样本区间内股价形成 2 个主升段与 3 个主跌段。 最大主跌段为 2023 年 3 月 至 2024 年 2 月，11 个月-38.8\%；最大主升段出现在 2024 年 8 月 至 2026 年 4 月，20 个月+56.2\%。 最近一次转折出现在 2026 年 6 月（下跌段自 2026 年 4 月起，2 个月-30.1\%），当前处于该段的延续之中。 公司股价较申万农林牧渔指数提前约 21 个月见底，是板块中较早定价周期反转的公司之一。 饲料夹在原料成本与养殖需求之间（第 \ref{sec:feed} 节），量的部分取决于存栏、利的部分取决于原料与成品的价差，公司 2026 年上半年实现盈利、半年 ROE 1.3\%（未年化），好于全样本中位数（0.75\%），下行周期对其报表的冲击小于同业。 综合区间形态看，该股单段涨跌幅度偏大、以趋势段为主（最大主升 +56.2\%、最大主跌 -38.8\%），年化波动率 36.7\% 与全样本中位数 37.4\% 相比波动与样本中位数接近。

\pfl{盈利情况} 2026 年上半年营业收入 6.69 亿元（同比 +9.0\%），归母净利润 0.10 亿元，实现扭亏。 以年报为趋势对照，2023---2025 年营业收入由 14.36 亿元变为 13.60 亿元（两年年均复合增速 -2.7\%），归母净利润由 0.49 亿元变为 -0.20 亿元。 按半年报口径，2026 年上半年的 ROE 为 1.3\%（未年化）、销售净利率 1.5\%、毛利率 13.7\%，ROE 在“畜禽饲料”11 家公司中按数值由高到低排第\zhnumber{1}位。 同期全样本半年 ROE 中位数 0.75\%，公司与之相差 0.52 个百分点（略好于中位数公司）。 从公司披露的原因看，业绩变动主要来自：2026 年 9 月，公司确定以 1.12 亿元收购上述三家标的公司。 综合看，公司在报告期维持盈利（高于全样本半年 ROE 中位数 0.75\%），利润的可持续性取决于畜禽饲料景气的延续与成本管控的执行。

\begin{center}\small\setlength{\tabcolsep}{3.4pt}
\captionof{table}[播恩集团（001366）关键财务指标]{播恩集团（001366）关键财务指标（合并报表口径；两个半年报期均未年化；现金短债比 $=$ 货币资金 $\div$ 短期债务）}
\begin{adjustbox}{max width=\textwidth}
\begin{tabular}{@{}lrrrrr@{}}
\toprule
指标 & 2023 年 & 2024 年 & 2025 年 & 2025 年半年报 & 2026 年半年报 \\
\midrule
营业收入（亿元） & 14.36 & 10.17 & 13.60 & 6.14 & 6.69 \\
归母净利润（亿元） & 0.49 & -0.29 & -0.20 & -0.15 & 0.10 \\
毛利率（\%） & 14.2 & 13.8 & 12.5 & 12.1 & 13.7 \\
ROE（\%） & 6.1 & -3.4 & -2.5 & -1.8 & 1.3 \\
资产负债率（\%） & 17.8 & 25.0 & 32.9 & 29.3 & 38.1 \\
经营活动现金流净额（亿元） & 0.46 & -0.31 & -0.28 & -0.95 & -1.39 \\
货币资金（亿元） & 3.52 & 2.91 & 2.93 & 2.21 & 2.60 \\
有息负债合计（亿元） & 0.75 & 1.51 & 2.56 & 2.21 & 3.68 \\
现金短债比（倍） & 35.39 & 3.13 & 1.69 & 1.37 & 1.10 \\
\bottomrule
\end{tabular}
\end{adjustbox}
\end{center}

\begin{center}
\begin{minipage}{\textwidth}\centering
\begin{tikzpicture}
\begin{groupplot}[
  group style={group size=2 by 1, horizontal sep=0.60cm},
  width=6.85cm, height=4.60cm,
  tick label style={font=\tiny},
  title style={font=\scriptsize\heiti},
  yticklabel style={font=\tiny},
  ylabel style={font=\tiny},
  grid=major, grid style={LightGray, line width=0.3pt},
  axis line style={InkGray, line width=0.5pt},
  enlarge x limits={abs=0.8},
  legend style={font=\scriptsize, at={(0.5,-0.20)}, anchor=north,
                legend columns=2, draw=none, fill=none, column sep=6pt},
]
\nextgroupplot[
  ybar, bar width=4pt,
  ymin=-2.09, ymax=19.87,
  xtick={0,1,2,3,4,5.7},
  xticklabels={2021,2022,2023,2024,2025,2026H1},
  ylabel={亿元},
]
\addplot[fill=AgriGreen!75, draw=AgriGreen, bar shift=-2.6pt] table[x=i, y=rev]{data/clean/profiles/fin_001366.dat};
\addplot[fill=SoftRed!60, draw=SoftRed, bar shift=2.6pt] table[x=i, y=ni]{data/clean/profiles/fin_001366.dat};
\addplot[fill=AgriGreen!30, draw=AgriGreen, pattern=north east lines, bar shift=-2.6pt] table[x=x, y=rev]{data/clean/profiles/finh_001366.dat};
\addplot[fill=SoftRed!20, draw=SoftRed, pattern=north east lines, bar shift=2.6pt] table[x=x, y=ni]{data/clean/profiles/finh_001366.dat};
\legend{营业收入（年/半年）, 归母净利润（年/半年）}
\nextgroupplot[
  ymin=-6.8, ymax=42.3,
  xtick={0,1,2,3,4,5.7},
  xticklabels={2021,2022,2023,2024,2025,2026H1},
  ylabel={\%},
]
\addplot[AgriGold, mark=*, mark size=1.3pt] table[x=i, y=roe]{data/clean/profiles/fin_001366.dat};
\addplot[OilBlue, mark=square*, mark size=1.3pt] table[x=i, y=debt]{data/clean/profiles/fin_001366.dat};
\addplot[only marks, AgriGold, mark=*, mark size=2.0pt] table[x=x, y=roe]{data/clean/profiles/finh_001366.dat};
\addplot[only marks, OilBlue, mark=square*, mark size=2.0pt] table[x=x, y=debt]{data/clean/profiles/finh_001366.dat};
\legend{ROE, 资产负债率}
\end{groupplot}
\end{tikzpicture}
\begin{tikzpicture}
\begin{axis}[
  name=finbot,
  width=13.75cm, height=3.10cm, scale only axis,
  ymin=-0.37, ymax=4.12,
  xtick={0,1,2,3,4,5.7},
  xticklabels={2021,2022,2023,2024,2025,2026H1},
  tick label style={font=\tiny},
  yticklabel style={font=\tiny},
  ylabel={亿元}, ylabel style={font=\tiny},
  ybar, bar width=4pt,
  grid=major, grid style={LightGray, line width=0.3pt},
  axis line style={InkGray, line width=0.5pt},
  enlarge x limits={abs=0.8},
  legend style={font=\scriptsize, at={(0.5,-0.26)}, anchor=north,
                legend columns=2, draw=none, fill=none, column sep=10pt},
]
\addplot[fill=AgriGreen!55, draw=AgriGreen, bar shift=-3.0pt] table[x=i, y=cash]{data/clean/profiles/fin_001366.dat};
\addplot[fill=SoftRed!45, draw=SoftRed, bar shift=3.0pt] table[x=i, y=ibd]{data/clean/profiles/fin_001366.dat};
\addplot[fill=AgriGreen!25, draw=AgriGreen, pattern=north east lines, bar shift=-3.0pt] table[x=x, y=cash]{data/clean/profiles/finh_001366.dat};
\addplot[fill=SoftRed!20, draw=SoftRed, pattern=north east lines, bar shift=3.0pt] table[x=x, y=ibd]{data/clean/profiles/finh_001366.dat};
\legend{货币资金（左轴）, 有息负债（左轴）}
\end{axis}
\begin{axis}[
  at={(finbot.south east)}, anchor=south east,
  width=13.75cm, height=3.10cm, scale only axis,
  axis y line*=right, axis x line*=none, xtick=\empty,
  ymin=-1.88, ymax=1.06,
  tick label style={font=\tiny},
  yticklabel style={font=\tiny}, scaled y ticks=false,
  legend style={font=\scriptsize, at={(0.5,-0.44)}, anchor=north,
                legend columns=1, draw=none, fill=none},
]
\addplot[OilBlue, mark=*, mark size=2.2pt, line width=1.3pt] table[x=i, y=ocf]{data/clean/profiles/fin_001366.dat};
\addplot[only marks, OilBlue, mark=*, mark size=3.2pt] table[x=x, y=ocf]{data/clean/profiles/finh_001366.dat};
\legend{经营活动现金流净额（右轴，亿元，蓝点）}
\end{axis}
\end{tikzpicture}

\captionof{figure}[播恩集团（001366）近年主要财务指标]{播恩集团（001366）近年主要财务指标（左上：营业收入与归母净利润；右上：ROE 与资产负债率；下：货币资金、有息负债为左轴柱状，经营活动现金流净额为其右轴的蓝点折线。
2021---2025 年为年报口径，2026H1 为 2026 年半年报/半年末，与其前一年报之间留出间隔、以斜纹或空心点区分；半年报数据未年化）}
\end{minipage}
\end{center}

\pfl{财务分析} 2026 年 6 月末资产负债率 38.1\%，较 2025 年末上升 5.2 个百分点（样本中位数 52.2\%，所处三级行业内按由低到高排第\zhnumber{1}位）。 2026 年 6 月末货币资金 2.60 亿元、短期债务（短期借款与一年内到期非流动负债）2.37 亿元、长期债务 1.31 亿元，有息负债合计 3.68 亿元，现金短债比 1.10 倍——覆盖充分。 净有息负债 1.08 亿元，相当于其 2026 年上半年营业收入的 16\%。 流动比率 2.08、速动比率 1.62，短期指标尚可。 2026 年上半年经营活动现金流净额 -1.39 亿元，净利润为正但经营现金流为负，利润的现金含量偏低。 2026 年 6 月末在建工程 0.63 亿元，较上年末减少 0.06 亿元，在建投入在收缩。 综合看，现金短债比 1.10 倍，短期偿付压力可控；资产负债率低于样本中位数 14.1 个百分点；有息负债中短期债务占比更高，期限结构仍以滚动续作为主。

\pfl{重大战略转型} 近三年公司的战略动作集中在：并购与重组（3 项）：2026 年 9 月，本次收购可新增牛、羊等反刍动物饲料产品线；2026 年 5 月，公司以现金收购希杰（沈阳）、希杰（天津）、希杰（聊城）饲料有限公司各 100\% 股权。 公告口径上，近三年（2023 年 9 月---2026 年 9 月）公司披露重大重组与并购类公告 9 条、对外投资与转型类公告 3 条、回购与股权激励类公告 1 条，合计公告 367 条。 第一大行业仍是“饲料行业”，收入占比 100.0\%，与 2021 年年报相比没有方向性变化。 综合判断，报告期内公司有 12 条与战略相关的公告落地（重大重组与并购 9 条、对外投资与转型 3 条），资本运作与主业调整之间存在直接关联。

\begin{center}\small\setlength{\tabcolsep}{4pt}
\captionof{table}[播恩集团（001366）历史融资与资本运作]{播恩集团（001366）历史融资与资本运作（据公司公告与公开报道整理；金额为披露值，含首次公开发行、再融资、债券、银行借款与重整投资等）}
\begin{adjustbox}{max width=\textwidth}
\begin{tabular}{@{}p{1.45cm}p{2.55cm}p{4.05cm}p{4.95cm}@{}}
\toprule
时间 & 方式 & 规模 & 用途与说明 \\
\midrule
2023-03 & IPO（深交所主板） & 发行 4，\allowbreak{}035 万股新股，\allowbreak{}发行价 9.32 元/\allowbreak{}股，\allowbreak{}募集资金总额 3.76 亿元 & 赣州年产 24 万吨饲料项目、\allowbreak{}浙江播恩年产 12 万吨饲料项目、\allowbreak{}播恩生物健康产业基地-维生素复合预混料项目、\allowbreak{}重庆八维生物年产 12 万吨饲料项目、\allowbreak{}研发中心建设项目；发行市盈率（发行后股本）22.97 倍；截至 2026 年中报累计投入 1.22 亿元 \\
2026-05-26 & 签署股权收购条款清单（Term Sheet） & 支付交易保证金 700.00 万元 & 拟收购希杰国际控股有限公司持有的希杰（沈阳）、\allowbreak{}希杰（天津）、\allowbreak{}希杰（聊城）饲料有限公司各 100\% 股权；暂定交易价格 1.25 亿元；约定若截至 2026 年 8 月 31 日未能签署正式股权转让协议等情形 \\
2026-09-14 & 现金收购股权（正式方案） & 股权转让款合计 1.12 亿元（较 5 月披露的 1.25 亿元暂定价格下调约 1 & 由全资子公司广东播恩投资有限公司收购希杰（沈阳）、\allowbreak{}希杰（天津）、\allowbreak{}希杰（聊城）饲料有限公司各 100\% 股权；三家公司将纳入合并报表范围；2025 年三家标的合计营收约 6.24 亿元 \\
\bottomrule
\end{tabular}
\end{adjustbox}
\end{center}

\pfl{历史融投资情况} 从现金流量表看，2025 年公司取得借款收到的现金 2.40 亿元、偿还债务支付的现金 1.35 亿元，筹资活动现金流净额 0.98 亿元（2023 年为取得借款 0.00 亿元、偿还债务 0.16 亿元）。 2026 年上半年筹资活动现金流净额为 1.35 亿元（取得借款 1.87 亿元、偿还债务 0.75 亿元，未年化），仍处于净融入状态。 2025 年分配股利、利润或偿付利息支付的现金合计 0.07 亿元。 公开披露的融资记录共 3 笔，工具类型以 IPO 为主；金额与用途见下表。 其中规模最大的两笔为：2023 年 3 月，IPO，发行 4，035 万股新股，发行价 9.32 元/股，募集资金总额 3.76 亿元；2026 年 9 月，现金收购股权，股权转让款合计 1.12 亿元（较 5 月披露的 1.25 亿元暂定价格下调约 1。 股本方面，近三年披露股本变动 5 次（A 股上市 1 次、限售股份上市 1 次）；总股本自 2021 年末以来未发生变化（1.61 亿股）。 分红方面，近三年现金分红 3 次、合计每股派现 0.35 元，最近一次实施在 2026 年 6 月。 综合看，公司融资结构上股权与债务融资并举；2025 年筹资活动现金流净额 0.98 亿元，年度内仍为净融入，与前述经营与投资现金流的方向基本一致。

\pfl{未来潜在融资需求分析} 2026 年 6 月末货币资金 2.60 亿元，而短期债务 2.37 亿元（现金短债比 1.10 倍），2026 年上半年经营现金流净额 -1.39 亿元。账面货币资金可以覆盖短期债务，缺口不体现在静态偿付层面。 公司 2026 年 6 月末资产负债率 38.1\%、2026 年上半年 ROE 1.3\%（未年化），资产端与负债端都没有明显压力，融资需求以相机抉择的扩张型用途为主。 公开披露的后续资金安排线索包括：资本开支安排：2026 年 8 月，募集资金结余 2.04 亿元；2023 年 3 月，IPO 募资改道并购受到媒体关注：公司 2023 年 3 月上市募得 3.19 亿元。 资金需求还要叠加在建工程：期末在建工程 0.63 亿元、2026 年上半年购建固定资产等支出 0.37 亿元（未年化），这部分支出在周期底部通常会被主动放缓。 综合看，公司在静态偿付层面不存在缺口，外部融资更多是配合扩张节奏与并购整合的主动性安排，而非偿债压力驱动。




\subsubsection{正虹科技（000702）}
\label{co:000702}

\pfl{公司基本情况} 正虹科技成立于 1994 年，湖南岳阳的老牌饲料企业，前身可追溯至 1986 年的国营屈原农场饲料厂，1997 年上市并有“中国饲料第一股”之称，2022 年控制权由屈原农垦转至岳阳国资观盛投资。 公司于 1997-03-18 在深交所主板首发上市，注册资本 3.47 亿元，总股本 3.47 亿股。 截至 2026 年 9 月 24 日收盘价 5.69 元，流通市值 19.7 亿元，近一年-20.9\%，流通市值在三级行业“畜禽饲料”的 11 家公司中按由大到小排第\zhnumber{11}位，近一年涨跌幅低于全样本中位数（-10.1\%）10.8 个百分点。 按 2026 年半年报披露的行业构成，收入占比前三位为饲料行业 89.8\%；饲养行业 8.8\%；其他(补充) 1.2\%。 股权结构方面，控股股东为岳阳观盛投资发展有限公司（持股 23.08\%）。

\pfl{历史股价} 本报告样本区间（2021 年 1 月 至 2026 年 9 月）内，股价由 5.23 元变为 5.69 元，累计 +8.8\%，区间最高 9.49 元（2024 年 10 月）、最低 4.08 元（2024 年 2 月）；当前价相当于区间最高价的 60.0\%，为区间最低价的 1.39 倍。 月度收益的年化波动率 42.4\%，高于全样本中位数 37.4\%，在三级行业“畜禽饲料”的 11 家公司中波动率由高到低排第\zhnumber{4}位。 把股价阶段与公司事件对齐看，2024 年 10 月 至 2026 年 6 月股价下跌-47.7\%，同期（2024 年 12 月）岳阳市屈原农垦有限责任公司持股 7.70\%（2，667.58 万股），其中冻结 1，330.00 万股。 整体上看，区间的几处大级别拐点都能在公司的资本运作、股权变动或风险事件上找到对应，股价波动有明显的公司层面驱动，而非单纯的行业 beta。

\begin{center}
\begin{minipage}{\textwidth}\centering
\begin{tikzpicture}
\begin{axis}[
  width=12.6cm, height=3.9cm, scale only axis,
  xmin=2020.922, xmax=2026.888, ymin=3.79, ymax=10.15,
  xtick={2021,2022,2023,2024,2025,2026},
  yearticks,
  yticklabel style={font=\scriptsize},
  ylabel={元/股}, ylabel style={font=\scriptsize},
  grid=major, grid style={LightGray, line width=0.3pt},
  axis line style={InkGray, line width=0.5pt},
  every axis plot/.append style={line width=0.9pt},
]
\addplot[AgriGreen] table[x=x,y=y]{data/clean/profiles/px_000702.dat};
\addplot[SoftRed, dashed, line width=0.7pt] table[x=x,y=y]{data/clean/profiles/pxzz_000702.dat};
\addplot[only marks, mark=*, mark size=1.1pt, SoftRed] table[x=x,y=y]{data/clean/profiles/pxzz_000702.dat};
\end{axis}
\end{tikzpicture}

\captionof{figure}[正虹科技（000702）股价与周期骨架]{正虹科技（000702）月度收盘价（元/股）与 ZigZag 周期骨架（阈值 25\%，圆点为周期转折点）}
\end{minipage}
\end{center}

\pfl{过去周期分析} 以 25\% 的 ZigZag 阈值划分，样本区间内股价形成 2 个主升段与 2 个主跌段。 最大主跌段为 2024 年 10 月 至 2026 年 6 月，20 个月-47.7\%；最大主升段出现在 2024 年 2 月 至 2024 年 10 月，8 个月+132.6\%。 最近一次转折出现在 2026 年 6 月（下跌段自 2024 年 10 月起，20 个月-47.7\%），当前处于该段的延续之中。 公司股价较申万农林牧渔指数提前约 27 个月见底，是板块中较早定价周期反转的公司之一。 饲料夹在原料成本与养殖需求之间（第 \ref{sec:feed} 节），量的部分取决于存栏、利的部分取决于原料与成品的价差，公司 2026 年 6 月末资产负债率 54.2\%、半年 ROE -6.2\%（未年化），在本轮下行中属于随行业同步调整的一类。 综合区间形态看，该股单段涨跌幅度偏大、以趋势段为主（最大主升 +132.6\%、最大主跌 -47.7\%），年化波动率 42.4\% 与全样本中位数 37.4\% 相比波动与样本中位数接近。

\pfl{盈利情况} 2026 年上半年营业收入 5.69 亿元（同比 +14.2\%），归母净利润 -0.21 亿元，亏损同比扩大 0.13 亿元。 以年报为趋势对照，2023---2025 年营业收入由 12.41 亿元变为 10.19 亿元（两年年均复合增速 -9.4\%），归母净利润由 -1.40 亿元变为 -0.14 亿元。 按半年报口径，2026 年上半年的 ROE 为 -6.2\%（未年化）、销售净利率 -4.9\%、毛利率 5.3\%，ROE 在“畜禽饲料”11 家公司中按数值由高到低排第\zhnumber{5}位。 按同一口径，三级行业“畜禽饲料”11 家公司的半年 ROE 中位数为 -6.25\%，该公司高于其中位数 0.08 个百分点。 面对这一局面，公司披露的举措包括：2026 年，饲料原料贸易业务适度调整以降低采购成本。 综合看，公司仍处亏损区间、亏损同比扩大，半年 ROE 在三级行业内按由高到低排第\zhnumber{5}位，单位盈利能力的修复依赖行业景气与自身成本端的改善，报表弹性主要体现为价格与销量的方向性变化。

\begin{center}\small\setlength{\tabcolsep}{3.4pt}
\captionof{table}[正虹科技（000702）关键财务指标]{正虹科技（000702）关键财务指标（合并报表口径；两个半年报期均未年化；现金短债比 $=$ 货币资金 $\div$ 短期债务）}
\begin{adjustbox}{max width=\textwidth}
\begin{tabular}{@{}lrrrrr@{}}
\toprule
指标 & 2023 年 & 2024 年 & 2025 年 & 2025 年半年报 & 2026 年半年报 \\
\midrule
营业收入（亿元） & 12.41 & 11.21 & 10.19 & 4.99 & 5.69 \\
归母净利润（亿元） & -1.40 & -0.65 & -0.14 & -0.08 & -0.21 \\
毛利率（\%） & 4.0 & 5.1 & 7.1 & 7.4 & 5.3 \\
ROE（\%） & -46.6 & -16.5 & -4.0 & -2.2 & -6.2 \\
资产负债率（\%） & 53.9 & 57.5 & 47.0 & 47.4 & 54.2 \\
经营活动现金流净额（亿元） & 0.61 & 0.11 & -0.31 & -0.18 & -0.63 \\
货币资金（亿元） & 2.49 & 2.67 & 0.66 & 0.92 & 0.60 \\
有息负债合计（亿元） & 2.84 & 2.82 & 1.41 & 1.33 & 1.97 \\
现金短债比（倍） & 0.91 & 0.96 & 0.51 & 0.71 & 0.33 \\
\bottomrule
\end{tabular}
\end{adjustbox}
\end{center}

\begin{center}
\begin{minipage}{\textwidth}\centering
\begin{tikzpicture}
\begin{groupplot}[
  group style={group size=2 by 1, horizontal sep=0.60cm},
  width=6.85cm, height=4.60cm,
  tick label style={font=\tiny},
  title style={font=\scriptsize\heiti},
  yticklabel style={font=\tiny},
  ylabel style={font=\tiny},
  grid=major, grid style={LightGray, line width=0.3pt},
  axis line style={InkGray, line width=0.5pt},
  enlarge x limits={abs=0.8},
  legend style={font=\scriptsize, at={(0.5,-0.20)}, anchor=north,
                legend columns=2, draw=none, fill=none, column sep=6pt},
]
\nextgroupplot[
  ybar, bar width=4pt,
  ymin=-3.90, ymax=16.30,
  xtick={0,1,2,3,4,5.7},
  xticklabels={2021,2022,2023,2024,2025,2026H1},
  ylabel={亿元},
]
\addplot[fill=AgriGreen!75, draw=AgriGreen, bar shift=-2.6pt] table[x=i, y=rev]{data/clean/profiles/fin_000702.dat};
\addplot[fill=SoftRed!60, draw=SoftRed, bar shift=2.6pt] table[x=i, y=ni]{data/clean/profiles/fin_000702.dat};
\addplot[fill=AgriGreen!30, draw=AgriGreen, pattern=north east lines, bar shift=-2.6pt] table[x=x, y=rev]{data/clean/profiles/finh_000702.dat};
\addplot[fill=SoftRed!20, draw=SoftRed, pattern=north east lines, bar shift=2.6pt] table[x=x, y=ni]{data/clean/profiles/finh_000702.dat};
\legend{营业收入（年/半年）, 归母净利润（年/半年）}
\nextgroupplot[
  ymin=-58.9, ymax=98.6,
  xtick={0,1,2,3,4,5.7},
  xticklabels={2021,2022,2023,2024,2025,2026H1},
  ylabel={\%},
]
\addplot[AgriGold, mark=*, mark size=1.3pt] table[x=i, y=roe]{data/clean/profiles/fin_000702.dat};
\addplot[OilBlue, mark=square*, mark size=1.3pt] table[x=i, y=debt]{data/clean/profiles/fin_000702.dat};
\addplot[only marks, AgriGold, mark=*, mark size=2.0pt] table[x=x, y=roe]{data/clean/profiles/finh_000702.dat};
\addplot[only marks, OilBlue, mark=square*, mark size=2.0pt] table[x=x, y=debt]{data/clean/profiles/finh_000702.dat};
\legend{ROE, 资产负债率}
\end{groupplot}
\end{tikzpicture}
\begin{tikzpicture}
\begin{axis}[
  name=finbot,
  width=13.75cm, height=3.10cm, scale only axis,
  ymin=-0.48, ymax=5.43,
  xtick={0,1,2,3,4,5.7},
  xticklabels={2021,2022,2023,2024,2025,2026H1},
  tick label style={font=\tiny},
  yticklabel style={font=\tiny},
  ylabel={亿元}, ylabel style={font=\tiny},
  ybar, bar width=4pt,
  grid=major, grid style={LightGray, line width=0.3pt},
  axis line style={InkGray, line width=0.5pt},
  enlarge x limits={abs=0.8},
  legend style={font=\scriptsize, at={(0.5,-0.26)}, anchor=north,
                legend columns=2, draw=none, fill=none, column sep=10pt},
]
\addplot[fill=AgriGreen!55, draw=AgriGreen, bar shift=-3.0pt] table[x=i, y=cash]{data/clean/profiles/fin_000702.dat};
\addplot[fill=SoftRed!45, draw=SoftRed, bar shift=3.0pt] table[x=i, y=ibd]{data/clean/profiles/fin_000702.dat};
\addplot[fill=AgriGreen!25, draw=AgriGreen, pattern=north east lines, bar shift=-3.0pt] table[x=x, y=cash]{data/clean/profiles/finh_000702.dat};
\addplot[fill=SoftRed!20, draw=SoftRed, pattern=north east lines, bar shift=3.0pt] table[x=x, y=ibd]{data/clean/profiles/finh_000702.dat};
\legend{货币资金（左轴）, 有息负债（左轴）}
\end{axis}
\begin{axis}[
  at={(finbot.south east)}, anchor=south east,
  width=13.75cm, height=3.10cm, scale only axis,
  axis y line*=right, axis x line*=none, xtick=\empty,
  ymin=-1.78, ymax=1.18,
  tick label style={font=\tiny},
  yticklabel style={font=\tiny}, scaled y ticks=false,
  legend style={font=\scriptsize, at={(0.5,-0.44)}, anchor=north,
                legend columns=1, draw=none, fill=none},
]
\addplot[OilBlue, mark=*, mark size=2.2pt, line width=1.3pt] table[x=i, y=ocf]{data/clean/profiles/fin_000702.dat};
\addplot[only marks, OilBlue, mark=*, mark size=3.2pt] table[x=x, y=ocf]{data/clean/profiles/finh_000702.dat};
\legend{经营活动现金流净额（右轴，亿元，蓝点）}
\end{axis}
\end{tikzpicture}

\captionof{figure}[正虹科技（000702）近年主要财务指标]{正虹科技（000702）近年主要财务指标（左上：营业收入与归母净利润；右上：ROE 与资产负债率；下：货币资金、有息负债为左轴柱状，经营活动现金流净额为其右轴的蓝点折线。
2021---2025 年为年报口径，2026H1 为 2026 年半年报/半年末，与其前一年报之间留出间隔、以斜纹或空心点区分；半年报数据未年化）}
\end{minipage}
\end{center}

\pfl{财务分析} 2026 年 6 月末资产负债率 54.2\%，较 2025 年末上升 7.2 个百分点（样本中位数 52.2\%，所处三级行业内按由低到高排第\zhnumber{2}位）。 2026 年 6 月末货币资金 0.60 亿元、短期债务（短期借款与一年内到期非流动负债）1.81 亿元、长期债务 0.16 亿元，有息负债合计 1.97 亿元，现金短债比 0.33 倍——覆盖偏紧。 净有息负债 1.37 亿元，相当于其 2026 年上半年营业收入的 24\%。 流动比率 1.00、速动比率 0.66，短期偿债能力处于正常区间。 2026 年上半年经营活动现金流净额 -0.63 亿元，经营现金流与利润同时为负，现金消耗较快。 综合看，账面货币资金对有息负债与短期债务的覆盖不足，滚动偿付对新增融资的依赖度较高；资产负债率高于样本中位数 2.0 个百分点；有息负债中短期债务占比更高，期限结构仍以滚动续作为主。 公司披露的股东与质押风险为：2024 年 12 月，岳阳市屈原农垦有限责任公司持股 7.70\%（2，667.58 万股），其中冻结 1，330.00 万股。

\pfl{重大战略转型} 公司的战略动作可归为以下几类：融资与资本运作（1 项）：2022 年 6 月，公司引入岳阳国资观盛投资作为控股股东，并筹划向观盛投资非公开发行股票募集 3 亿多元用于补充流动资金；其他动作（1 项）：2023 年生猪生产量同比减少 49.39\%。 公告口径上，近三年（2023 年 9 月---2026 年 9 月）公司披露重大重组与并购类公告 5 条、对外投资与转型类公告 3 条、回购与股权激励类公告 0 条，合计公告 308 条。 第一大行业“饲料行业”的收入占比由 2021 年年报的 61.9\% 变为 89.8\%（28.0 个百分点），收入结构出现再平衡。 综合判断，报告期内公司有 8 条与战略相关的公告落地（重大重组与并购 5 条、对外投资与转型 3 条），资本运作与主业调整之间存在直接关联。

\begin{center}\small\setlength{\tabcolsep}{4pt}
\captionof{table}[正虹科技（000702）历史融资与资本运作]{正虹科技（000702）历史融资与资本运作（据公司公告与公开报道整理；金额为披露值，含首次公开发行、再融资、债券、银行借款与重整投资等）}
\begin{adjustbox}{max width=\textwidth}
\begin{tabular}{@{}p{1.45cm}p{2.55cm}p{4.05cm}p{4.95cm}@{}}
\toprule
时间 & 方式 & 规模 & 用途与说明 \\
\midrule
1997-03 & IPO（深交所） & 发行社会公众股 5,\allowbreak{}000 万股（含职工内部股 500 万股），\allowbreak{}发行价 7.06 元/\allowbreak{}股 & 饲料产能建设；中国证监会证监发(1997)44 号、\allowbreak{}45 号文批准；上市时湖南省岳阳市屈原农垦集团公司代表国家持有 61.59\% 股权 \\
2022-06-17 & 控制权安排（股份协议转让 + 表决权委托 + 定增） & 屈原农垦向观盛农业转让 4，\allowbreak{}034.18 万股（占发行前总股本 15.13\%）；屈原农垦将其剩余 2 & 巩固控制权并补充流动资金；收购报告书披露变动数量合计 1.47 亿股；认购完成后观盛投资合计将支配公司 42.41\% 的表决权 \\
2022-06-17 & 发行对象与关联方情况 & 君泰控股为公司第一大股东观盛农业的少数股东；君泰科技及相关主体系君泰控股控制的企业 & 发行人实际控制人 2022 年变更后，\allowbreak{}董事、\allowbreak{}监事、\allowbreak{}高级管理人员变动较大 \\
2023-07 & 向特定对象发行股票（定增） & 向控股股东观盛投资发行不超过 7，\allowbreak{}999.04 万股 & 全部用于补充流动资金 \\
\bottomrule
\end{tabular}
\end{adjustbox}
\end{center}

\pfl{历史融投资情况} 从现金流量表看，2025 年公司取得借款收到的现金 2.66 亿元、偿还债务支付的现金 4.10 亿元、吸收权益性投资 0.07 亿元，筹资活动现金流净额 -1.56 亿元（2023 年为取得借款 3.15 亿元、偿还债务 4.80 亿元）。 2026 年上半年筹资活动现金流净额为 0.60 亿元（取得借款 0.54 亿元、偿还债务 0.05 亿元，未年化），仍处于净融入状态。 2025 年分配股利、利润或偿付利息支付的现金合计 0.07 亿元。 公开披露的融资记录共 4 笔，工具类型以 IPO、定向增发为主；金额与用途见下表。 其中规模最大的两笔为：2023 年 7 月，定向增发，向控股股东观盛投资发行不超过 7，999.04 万股。 股本方面，近三年披露股本变动 4 次（增发新股上市 1 次）；总股本由 2021 年末的 2.67 亿股变为 3.47 亿股（+30.00\%）。 分红方面，近三年未实施现金分红，在全部样本中属于分红能力偏弱的一端。 近三年“再融资”类公告 3 条，涉及定向增发。 综合看，公司融资结构上以债务滚动为主、股权融资为补充；2025 年筹资活动现金流净额 -1.56 亿元，融资节奏仍以维持存量债务周转为主，与前述经营与投资现金流的方向基本一致。

\pfl{未来潜在融资需求分析} 2026 年 6 月末货币资金 0.60 亿元，而短期债务 1.81 亿元（现金短债比 0.33 倍），2026 年上半年经营现金流净额 -0.63 亿元。按“短期债务减去账面货币资金”的滚动偿付口径，静态缺口约 1.21 亿元；若再计入上半年亏损的年化额（0.41 亿元），维持一年正常经营所需的外部资金约 1.21---1.63 亿元。 公司 2026 年上半年归母净利润为负（-0.21 亿元），但 6 月末资产负债率 54.2\% 尚未进入第一档（亏损且负债率 $>$65\%）的区间，当前融资需求以补充流动资金、支撑低谷期经营性支出为主。 公开披露的后续资金安排线索包括：股权融资线索：2022 年 6 月，后者正筹划公司控制权变更事项，并同步筹划向观盛投资非公开发行股票；2023 年 7 月，公司向观盛投资的定增（不超过 3.38 亿元）认购资金到位。 综合看，公司的外部资金需求以补充流动性与项目投入为主，滚动偿付口径下的静态缺口约 1.21 亿元，能否落地取决于授信续作与发行窗口的配合。




